\documentclass[10pt,a4paper]{article}

% ----------------- Paquetes -------------------%
\usepackage[T1]{fontenc}
\usepackage[utf8]{inputenc}
\usepackage[a4paper,margin=2.5cm]{geometry}
\usepackage{mathptmx}             
\usepackage{changepage} 
\usepackage{setspace}
\usepackage[spanish]{babel}
\usepackage[labelfont=bf,labelsep=space]{caption}
\usepackage{amssymb}
\usepackage{amsmath}
\usepackage{ebgaramond}
\usepackage{placeins}
\usepackage{etoolbox}
\usepackage{setspace}
\usepackage{parskip} 
\usepackage{tcolorbox}
\usepackage{needspace}
\usepackage{xparse}
\usepackage{ocgx2}
\usepackage{hyperref}
\usepackage{hypcap}
\usepackage{multirow}
\usepackage{soul}\sethlcolor{yellow}% resaltado amarillo: \hl{texto} (Panel TCEE, ⌃H)


%%%%%%%%%%%%%%%%%%%%%%%%%%%%%%%%%%%%%%%%%%%%%%%%%%%%%%%%%%%%%%%%
% --- Fija primero tus valores globales "buenos" ---
\setstretch{1.2}                 % Interlineado global
\setlength{\parskip}{0.7\baselineskip}
\setlength{\parindent}{0pt}
\newlength{\GlobalParskip}
\setlength{\GlobalParskip}{\parskip}

\newcounter{eqnum}[subsection]
\renewcommand{\theeqnum}{\arabic{eqnum}}
\newcounter{alphaitem}
\newcounter{numitem}

%% ------------------- Recuadros----------------------%%
\newcommand{\puntosclave}[1]{%
  \begin{tcolorbox}[
    colframe=black,
    title={Puntos clave},
    before upper={\setlength{\parindent}{0.5cm}\setlength{\parskip}{0.5\baselineskip}}
  ]
  #1
  \end{tcolorbox}
}

\newcommand{\modificaciones}[1]{
  \begin{tcolorbox}[
    colback=white,
    colframe=red,
    title={Modificaciones a realizar},
    before upper={\setlength{\parindent}{0.5cm}\setlength{\parskip}{0.5\baselineskip}}
  ]
  #1
  \end{tcolorbox}
}

%% ------------------- Listas----------------------%%
    % --- Lista alfabética ---
    \newenvironment{la}
      {\begin{list}{\alph{alphaitem})}{%
          \usecounter{alphaitem}%
          \setlength{\leftmargin}{1cm}%
          \setlength{\parskip}{3pt}% 
          \setlength{\itemsep}{0pt}% ← espacio entre ítems
          \setlength{\parsep}{0pt}%  ← párrafos dentro de un ítem
      }}
      {\end{list}}
      
    % --- Lista numérica ---
    \newenvironment{lnum}
      {\begin{list}{\arabic{numitem}.}{%
          \usecounter{numitem}%
          \setlength{\leftmargin}{1cm}%
          \setlength{\parskip}{3pt}% 
          \setlength{\itemsep}{0pt}% ← espacio entre ítems
          \setlength{\parsep}{0pt}%  ← párrafos dentro de un ítem
      }}
      {\end{list}}
      
% ------------------- Imágenes----------------------%
    \graphicspath{{Img/}}% Rutas por defecto 
    % --- Imagen adaptativa ---
    % Registros de dimensión definidos una sola vez
    \newdimen\imgcap
    \newdimen\imgmin
    \newdimen\imgmax
    \newdimen\imgremain
    \newdimen\imgh
    \newdimen\imgneed
    
    \newcommand{\imagenfit}[3]{%
      \addvspace{10pt}% separación antes
      \par\begingroup
        % Parámetros de composición (asignaciones, no \newdimen)
        \imgcap=1\baselineskip            % alto estimado de título+fuente
        \imgmin=0.15\textheight
        \imgmax=0.15\textheight
        \imgremain=\pagegoal
        \ifdim\pagegoal=\maxdimen \imgremain=0pt\fi
        \advance\imgremain by -\pagetotal
        \advance\imgremain by -\imgcap
        % Decisión de altura destino
        \ifdim\imgremain<\imgmin
          \newpage
          \imgh=0.20\textheight
        \else
          \imgh=\imgremain
          \ifdim\imgh>\imgmax \imgh=\imgmax \fi
        \fi
        % Reservar espacio para evitar cortes del bloque completo
        \imgneed=\imgh \advance\imgneed by \imgcap
        \Needspace{\imgneed}
        % Cuerpo
        \noindent\begin{minipage}{\textwidth}
          \centering
          \captionof{figure}{\textemdash\ #2}\par\vspace{0.3em}% título arriba
          \includegraphics[width=\linewidth,height=\imgh,keepaspectratio]{\detokenize{#1}}\par
          \vspace{0.3em}\small\textit{Fuente: }#3% fuente abajo
        \end{minipage}%
      \endgroup
      \par\addvspace{10pt}%
    }

%------------ Bloque de RECUADRO ESPECIAL -------------%
    \newenvironment{specialblock}[1]{%
      \begin{quote} % crea sangría a izquierda y derecha
      \small % texto más pequeño
      \noindent\textbf{#1}\\[-0.3em] % título en negrita
      \rule{\linewidth}{0.4pt}\\[0.6em]}{
      \end{quote}}
      

%-------------------------- Author Font -----------------------%
    \newcommand{\authorfont}[1]{{\fontfamily{EBGaramond-TLF}\selectfont #1}}

%-------------------------- Anexos -----------------------%
    \makeatletter
    \newcounter{anexo}
    \renewcommand{\theanexo}{\Roman{anexo}}
    \newcommand{\anexo}[1]{%
      \clearpage
      \phantomsection
      \refstepcounter{anexo}%
      \noindent\textbf{Anexo \theanexo.- }#1\par\medskip
      \addcontentsline{stc}{section}{Anexo \theanexo.- #1}%
    }
    \makeatother

    %-------------------------- Lista de figuras (personal) -----------------------%
\makeatletter
\newcommand{\MyListOfFigures}{%
  \clearpage
  \phantomsection
  {\centering\bfseries\Large Índice de figuras\par}%
  \medskip
  % Entrada en nuestro índice de contenido (stc)
  \addcontentsline{stc}{section}{Índice de figuras}%
  % Recuperación del listado (sin cabecera propia)
  \begingroup
    \parindent=0pt\parskip=2pt
    \@starttoc{lof}%
  \endgroup
}
\makeatother

%-------------------------- Bibliografía -----------------------%
\makeatletter
% Contador para las referencias
\newcounter{myref}
% Cabecera de la bibliografía, entra en el índice 'stc'
\newcommand{\MyBibliography}{%
  \clearpage
  \phantomsection
  {\centering\bfseries\Large Bibliografía y referencias\par}%
  \medskip
  \addcontentsline{stc}{section}{Bibliografía y referencias}%
  \setcounter{myref}{0}%
}
\newcommand{\MyBibItem}[4]{%
  \refstepcounter{myref}%
  \noindent[\themyref]\ #1.\ \emph{#2}. #3. #4\par
}
\makeatother

%--------- Establecer cuatro niveles de índice --------%
    \setcounter{tocdepth}{4}   
    \setcounter{secnumdepth}{4}% numera hasta \paragraph

%%%%%%%%%%%%%%%%%%%%%%%%%%%%%%%%

\makeatletter

    %--------- Interlineado Global --------%
    \edef\GlobalBaseLineStretch{\baselinestretch}
    
    %--------- Espaciado de seccinones --------%
    \renewcommand\section{\@startsection{section}
      {1}{\z@}
      {2.5ex \@plus 1ex \@minus .2ex} % espacio antes
      {1.5ex \@plus .2ex}             % espacio después
      {\normalfont\Large\bfseries}    % formato del título
    }
    \renewcommand\subsection{\@startsection{subsection}{2}{\z@}
      {3ex \@plus 0.8ex \@minus .2ex}
      {1.5ex \@plus .2ex}
      {\normalfont\large\bfseries}
    }
    \renewcommand\subsubsection{\@startsection{subsubsection}{3}{\z@}
      {3ex \@plus 0.5ex \@minus .2ex}
      {1ex \@plus .2ex}
      {\normalfont\normalsize\bfseries}
    }
    \renewcommand\paragraph{\@startsection{paragraph}{4}{\z@}%
    {-2ex \@plus -0.5ex \@minus -.2ex}% espacio antes
    {0.8ex \@plus .2ex}% espacio después
    {\normalfont\normalsize\itshape}}% ← cursiva en lugar de negrita

    %--------- Ecuaciones --------%   
    % Variables globales para almacenar títulos y notas
    \gdef\leq@current@section{}%
    \gdef\leq@current@subsection{}%
    \gdef\leq@last@printed@section{}%
    \gdef\leq@last@printed@subsection{}%
    
    % Comando para añadir nota (invisible en el cuerpo)
    \newcommand{\eqnote}[1]{%
      \addcontentsline{leq}{leqnote}{#1}%
    }
    
    % Comando para ecuaciones
    \newcommand{\eqblock}[2]{%
      % Verificar si necesitamos imprimir el título de sección
      \ifx\leq@current@section\leq@last@printed@section\else
        \addcontentsline{leq}{leqsection}{\leq@current@section}%
        \global\let\leq@last@printed@section\leq@current@section
        \gdef\leq@last@printed@subsection{}% Reset subsección al cambiar sección
      \fi
      % Verificar si necesitamos imprimir el título de subsección
      \ifx\leq@current@subsection\leq@last@printed@subsection\else
        \ifx\leq@current@subsection\@empty\else
          \addcontentsline{leq}{leqsubsection}{\leq@current@subsection}%
          \global\let\leq@last@printed@subsection\leq@current@subsection
        \fi
      \fi
      % Ecuación
      \refstepcounter{eqnum}%
      \begin{equation*}
        #1\tag{E.\theeqnum}
      \end{equation*}%
      \addcontentsline{leq}{leqt}{[E.\theeqnum] -- #2}%
      \addcontentsline{leq}{leqm}{#1}%
    }
    
    %%% ----- Índice de ecuaciones ----------%%%
    % Tamaño para []
    \newcommand{\leqIndexFont}{\normalsize}
    \newcommand{\leqTitleFont}{\tiny}
    \newcommand{\leqNoteFont}{\tiny}
    % Cabecera de sección 
    \newcommand*\l@leqsection[2]{%
      \addvspace{25pt}% Mayor separación antes de nueva sección
      \noindent\leqIndexFont
      \makebox[\linewidth]{\rule{.38\linewidth}{0.4pt}\ \ \textbf{  \MakeUppercase{#1}}\ \ \rule{.38\linewidth}{0.4pt}}%
      \par\vspace{-10pt}
    }
    % Cabecera de subsección 
    \newcommand*\l@leqsubsection[2]{%
      \par\addvspace{15pt}%
      \noindent\leqIndexFont #1
      \par\vspace{-7pt}
    }
    % Título de ecuación 
    \newcommand*\l@leqt[2]{%
    \par\addvspace{10pt}%
      \par\noindent\leqTitleFont #1\par
    }
    % Miniatura de ecuación
    \newcommand*\l@leqm[2]{%
      \vspace{0.3\baselineskip}%
      \par\scriptsize\noindent\hspace*{1.5em}\(#1\)\par%
    }
    % Nota de ecuación 
    \newcommand*\l@leqnote[2]{%
      \vspace{0.2\baselineskip}%
      \par\noindent\leqNoteFont\hspace*{1.5em}\textbf{Nota.--} #1\par\vspace{0.5\baselineskip}%
    }
    % Generación del índice
    \newcommand{\mylistofequations}{%
      \clearpage
      \phantomsection
      % Título visual de la lista (ajústalo a tu gusto):
      {\centering\bfseries\Large Índice de ecuaciones\par}
      % Entrada en nuestro índice 'stc':
      \addcontentsline{stc}{section}{Índice de ecuaciones}%
      % Llamada a tu lista real (usa el nombre exacto de tu macro):
      \@starttoc{leq}%
    }
    
    %--------- Captura de títulos de sección --------%
    \let\leq@original@section\section
    \let\leq@original@subsection\subsection
    % Flag para saber si estamos en una sección asterisco
    \newif\ifLeqStarSection
    \LeqStarSectionfalse
    
    % Redefinir \section CON CONTROL DE ASTERISCO
    \renewcommand{\section}{%
      \@ifstar{\leq@section@star}{\@dblarg\leq@section@nostar}%
    }
    \def\leq@section@star#1{%
      \global\LeqStarSectiontrue
      \gdef\leq@current@section{#1}%
      \gdef\leq@current@subsection{}%
      \leq@original@section*{#1}%
      \global\LeqStarSectionfalse
    }
    \def\leq@section@nostar[#1]#2{%
      \global\LeqStarSectionfalse
      \gdef\leq@current@section{#2}%
      \gdef\leq@current@subsection{}%
      \leq@original@section[#1]{#2}%
    }
    % Redefinir \subsection
    \renewcommand{\subsection}{%
      \@ifstar{\leq@subsection@star}{\@dblarg\leq@subsection@nostar}%
    }
    \def\leq@subsection@star#1{%
      \global\LeqStarSectiontrue
      \gdef\leq@current@subsection{#1}%
      \leq@original@subsection*{#1}%
      \global\LeqStarSectionfalse
    }
    \def\leq@subsection@nostar[#1]#2{%
      \global\LeqStarSectionfalse
      \gdef\leq@current@subsection{#2}%
      \leq@original@subsection[#1]{#2}%
    }

    % ----- Restauración de valores Globales de estilo ----------%
    % Flags
    \newif\ifInFootnote
    \InFootnotefalse
    \pretocmd{\@footnotetext}{\InFootnotetrue}{}{}
    \apptocmd{\@footnotetext}{\InFootnotefalse}{}{}
    % Profundidad de listas (soporta anidadas)
    \newcount\MyListDepth \MyListDepth=0
    \AtBeginEnvironment{la}{\advance\MyListDepth by 1}
    \AfterEndEnvironment{la}{\advance\MyListDepth by -1}
    \AtBeginEnvironment{lnum}{\advance\MyListDepth by 1}
    \AfterEndEnvironment{lnum}{\advance\MyListDepth by -1}
    \AtBeginEnvironment{itemize}{\advance\MyListDepth by 1}
    \AfterEndEnvironment{itemize}{\advance\MyListDepth by -1}
    \AtBeginEnvironment{enumerate}{\advance\MyListDepth by 1}
    \AfterEndEnvironment{enumerate}{\advance\MyListDepth by -1}
    % Restauración diferida: hazlo al INICIO del próximo párrafo real
    \newif\if@pendingRestore
    \@pendingRestorefalse
    \newcommand{\RequestRestore}{\global\@pendingRestoretrue}
    \everypar{%
      \if@pendingRestore
        \global\@pendingRestorefalse
        \ifInFootnote\else
          \ifnum\MyListDepth=0
            \setlength{\parskip}{\GlobalParskip}%
            \setstretch{\GlobalBaseLineStretch}%
          \fi
        \fi
      \fi
    }
    % Óscar final de bloques:
    \AfterEndEnvironment{equation}{\RequestRestore}
    \AfterEndEnvironment{equation*}{\RequestRestore}
    \AfterEndEnvironment{la}{\RequestRestore}
    \AfterEndEnvironment{lnum}{\RequestRestore}
    % Óscar final de macros:
    \apptocmd{\eqblock}{\RequestRestore}{}{}
    \apptocmd{\imagenfit}{\RequestRestore}{}{}
    
\makeatother

% ===================== ToC propio con 4 niveles (único bloque) =====================
\makeatletter

% --- Escritores: añadimos una línea al .stc en cada cabecera numerada ---
% Guardamos las versiones actuales para llamar al comportamiento normal
\let\MyTOC@orig@section\section
\let\MyTOC@orig@subsection\subsection
\let\MyTOC@orig@subsubsection\subsubsection
\let\MyTOC@orig@paragraph\paragraph

% SECTION
\renewcommand{\section}{%
  \@ifstar{\MyTOC@section@star}{\@dblarg\MyTOC@section@nostar}%
}
\def\MyTOC@section@star#1{%
  \MyTOC@orig@section*{#1}% sin entrada al .stc
}
\def\MyTOC@section@nostar[#1]#2{%
  \MyTOC@orig@section[{#1}]{#2}% comportamiento normal (escribe al .toc)
  % Escribimos también nuestra línea al .stc (texto con \numberline)
  \addcontentsline{stc}{section}{\protect\numberline{\thesection}#2}%
}

% SUBSECTION
\renewcommand{\subsection}{%
  \@ifstar{\MyTOC@subsection@star}{\@dblarg\MyTOC@subsection@nostar}%
}
\def\MyTOC@subsection@star#1{%
  \MyTOC@orig@subsection*{#1}%
}
\def\MyTOC@subsection@nostar[#1]#2{%
  \MyTOC@orig@subsection[{#1}]{#2}%
  \addcontentsline{stc}{subsection}{\protect\numberline{\thesubsection}#2}%
}

% SUBSUBSECTION
\renewcommand{\subsubsection}{%
  \@ifstar{\MyTOC@subsubsection@star}{\@dblarg\MyTOC@subsubsection@nostar}%
}
\def\MyTOC@subsubsection@star#1{%
  \MyTOC@orig@subsubsection*{#1}%
}
\def\MyTOC@subsubsection@nostar[#1]#2{%
  \MyTOC@orig@subsubsection[{#1}]{#2}%
  \addcontentsline{stc}{subsubsection}{\protect\numberline{\thesubsubsection}#2}%
}

% PARAGRAPH
\renewcommand{\paragraph}{%
  \@ifstar{\MyTOC@paragraph@star}{\@dblarg\MyTOC@paragraph@nostar}%
}
\def\MyTOC@paragraph@star#1{%
  \MyTOC@orig@paragraph*{#1}%
}
\def\MyTOC@paragraph@nostar[#1]#2{%
  \MyTOC@orig@paragraph[{#1}]{#2}%
  % Si no numeras los \paragraph, cambia \theparagraph por texto plano sin \numberline
  \addcontentsline{stc}{paragraph}{\protect\numberline{\theparagraph}#2}%
}

% --- Impresor del índice propio (lee el .stc) ---
\newcommand{\MyTOC}{%
  \par\bigskip
  {\centering\bfseries\Large Índice de contenido\par}%
  \medskip
  \begingroup
    \parindent=0pt\parskip=2pt
    \@starttoc{stc}%
  \endgroup
}

\makeatother
% ===================== fin ToC propio =====================


%%%%%%%%%%%%%%


% --- Datos del tema ---
\title{Teoría impositiva (IV): Imposición y oferta\\
\large Efectos incentivo de los impuestos}
\author{Marcelo Ortega}
\date{Agosto 2026}

\begin{document}
\maketitle
\newpage

% --- Página de resumen y modificaciones ---
\puntosclave{ %% Escribir puntos clave Apuntes CECO
     Dentro del bloque de Economía Pública del temario, éste es el tercero de los cuatro de los títulos de tributos (488 a 4811 ). Los dos primeros tienen un enfoque eminentemente (i) microeconómico y (ii) positivo, mientras que el del último es (i) microeconómico y (ii) nonnativo -al menos en gran medida-; así, éste difiere de todos ellos en tanto que se sitúa '·a caballo .. entre la microeconomía y la macroeconomía.

     El mensaje central de este tema es que los tributos tienen capacidad ele generar incentivos para afectar la conducta de los agentes individuales y que por tanto son una herramienta útil de política económica. Se ha optado por, en línea con la literatura. priorizar el estudio de la primera de las cuestiones; sin embargo, en el momento del ·'cante", la segunda parte debe de estar presente. Se recomienda al opositor que introduzca valoraciones orales sobre las implicaciones de los modelos y la evidencia empírica sobre el uso de los impuestos como herramientas de política económica, por ejemplo a modo de transición entre subapartados.
     
     \textcolor{red}{\textbf{OJO -> No está incluido pero parece interesante una vez se estructure el Tema:}} Adicionalmente. se incluyen dos anexos relacionados con esta cuestión de los incentivos que el opositor debe de conocer y que. potencialmente, puede incluir en mayor o menor medida en su exposición: 1) imposición pigouviana. 2) diseño de mecanismos.
    }
\vspace{1cm}

\modificaciones{
    
    }

\newpage
\MyTOC
\newpage

\mylistofequations
\clearpage

\MyListOfFigures
\clearpage


\pagenumbering{arabic}
\setcounter{page}{1}


\section{Introducción}



\subsection{Contextualización}


\subsubsection{Relevancia}




\subsubsection{Problemática}



\subsection{Estructura de la exposición}






\clearpage
\section{Dimensión conductual: Marco conceptual}
\puntosclave{
    
}

En los últimos años, numerosas economías han llevado a cabo o planean implementar reformas en su sistema tributario con objeto de adaptarlo a los retos actuales. Entre otros, el contexto actual se caracteriza por un trade-off entre la necesidad de garantizar la sostenibilidad de las finanzas públicas, un crecimiento débil y las restricciones que enfrenta la Política Monetaria.

Como podemos intuir, impuestos diferentes tienen efectos diferentes sobre las decisiones de los agentes, la eficiencia económica y las metas distributivas del Estado. De esta forma, resulta necesario asegurar que las figuras impositivas del sistema tributario se encuadran adecuadamente en el triángulo formado por estos objetivos de Política Económica y resultan suficientes para cubrir las necesidades de gasto público.

\textcolor{red}{El sector público necesita recaudar impuestos para llevar a cabo sus funciones. Tradicionalmente, ase asignan tres funciones principales al sector público: 1) asignativa, con o��jeto de contribuir a  una asignación eficiente de recursos como consecuencia de la existencia de fallos de mercado, 2) distributiva. con objeto de contribuir a una mayor equidad distributiva de la renta y la riqueza, y 3) estabilizadora, con objeto de contribuir al equilibrio interno (y externo). i.e. más empleo. estabilidad de precios y crecimiento económico.}

No obstante, en la práctica los impuestos causan una serie de distorsiones que afectan al comportamiento de los agentes modificando sus  decisiones de trabajo, ahorro y/o producción. Así, los impuestos pueden también reducir el bienestar social mediante dos vías: (i) un efecto renta, detrayendo renta de los individuos; (ii) un efecto sustitución, alterando los precios relativos de bienes y factores productivos dando lugar a una pérdida de eficiencia o efecto gravamen.

Adicionalmente, ante el establecimiento de impuestos los agentes tienen incentivos para evadirlos, aunque éstos pueden reducirse con mecanismos de control y sanción. A lo largo de la historia, diversas teorías han tratado de explicar los determinantes que influyen en la decisión de cumplir o no con las obligaciones tributarias, así como los diversos mecanismos de control y sanción para desincentivar los comportamientos estratégicos de los contribuyentes. La importancia de su estudio radica en que la evasión fiscal provoca efectos distributivos y atenta contra la equidad en tanto da lugar a transferencias desde quienes contribuyen hacia quienes defraudan.

Así, el objeto de este tema es analizar cómo los impuestos modifican las decisiones de oferta de los agentes económicos. Para ello, en primer lugar estudiaremos cómo distintos impuestos tendrán diferentes efectos económicos sobre las decisiones óptimas de consumo, de ahorro. de inversión. producción o de trabajo. En segundo lugar. se expondrán cómo estas distorsiones en el comportamiento ele los contribuyentes pueden afectar negativamente a la recaudación tributaria. Por último. nos centraremos en una serie de enfoques alternativos que buscan estudiar cómo los subsidios o transferencias monetarias afectan a las decisiones de oferta en el mercado de trabajoy hasta qué punto los programas de bienestar pueden reducir la participación en el mercado laboral al fomentar la dependencia del Estado.





\clearpage
\section{Dimensión conductual: Marco analítico}
\puntosclave{
    El sistema impositivo modifica las decisiones óptimas de los agentes. A pesar de ello, su influencia queda indeterminada ya que depende de la importancia relativa de los efectos sustitución y renta. 
    
    La introducción del tiempo permite el estudio del efecto de los impuestos sobre el ahorro, ampliar los resultados de los modelos estáticos y, sin embargo, no supone desviaciones significativas en las conclusiones. Por otro lado, los impuestos alteran las decisiones financieras de las empresas relativas al reparto de dividendos y las decisiones de financiación.

    Partiendo de las distorsiones sobre las ofertas factoriales. LAFFER introduce una idea controvertida en el debate político: reducciones del tipo impositivo no conducen necesariamente a pérdidas de recaudación. 
    
    El sistema impositivo se ve afectado por problemas de información asimétrica de forma que los agentes pueden tener incentivos a llevar a cabo comportamientos estratégicos para alcanzar sus propios intereses, lo cual puede tener un impacto negativo en la recaudación fiscal. Por tanto, a través de la regulación el sector público puede establecer distintos mecanismos que interactúen con los agentes para así incentivar la declaración de su renta verdadera. 
    
    Superadas las teorías de los Economistas de la Oferta, desarrollos más recientes estudian el efecto distributivo de la fiscalidad sobre los factores en términos agregados.
}

La función tributaria es el conjunto de actividades que realiza el Estado para aplicar el sistema de impuestos, recaudar fondos y financiar los gastos públicos según la capacidad económica de cada ciudadano


Entre 1987 y 1988, Islandia reformó por completo su sistema tributario. Antes de la reforma, los trabajadores pagaban impuestos sobre los ingresos del año anterior (de forma muy similar a como ocurría en Estados Unidos), y la carga tributaria media era del $14,5\%$ de la renta, con tipos marginales que llegaban hasta el $56,3\%$. El nuevo sistema pasó a ser de pago a medida que se perciben los ingresos (\textit{pay-as-you-go}), de modo que los trabajadores pagaban un impuesto sobre la renta de tipo único del $32,5\%$ conforme recibían cada salario. 

Lo interesnte, no obstante, fue que durante la transición hacia este nuevo sistema, los trabajadores pagaron en 1987 los impuestos correspondientes a su renta de 1986, y en 1988 pagaron los impuestos correspondientes a su renta de 1988. ¡La renta obtenida en 1987 nunca llegó a tributar! Durante un año, tanto el tipo medio como el tipo marginal sobre las rentas del trabajo fueron cero. ¿Afectó realmente este cambio radical de la carga tributaria a las decisiones de los trabajadores sobre cuántas horas trabajar? BIANCHI et al. (2001) investigaron esta cuestión examinando la respuesta de los trabajadores islandeses a este año de vacaciones fiscales (tax holiday), y encontraron importantes efectos sobre la oferta de trabajo y el crecimiento económico. En promedio, cada aumento del $1\%$ en los salarios después de impuestos llevó a los trabajadores a trabajar un $0,4\%$ más de semanas que anteriormente, de manera que, según su medida, el empleo total aumentó bruscamente en 1987, pasando del $78\%$ al $85\%$: 

\imagenfit{Fig1.png}{Islandia: Experimento de imposición sobre la oferta de trabajo}{}

El PIB real también pasó de una tasa anual de crecimiento del $4,3\%$ al $8,5\%$. No obstante, estos efectos fueron transitorios. En 1988, estas cifras volvieron a niveles comparables a los existentes antes de la reforma tributaria. Esta llamativa respuesta a un cambio en la política tributaria demuestra que los impuestos pueden influir en decisiones importantes, como cuánto trabajar. La existencia de tales respuestas pone de manifiesto la tensión entre la equidad y la eficiencia en el diseño de la política tributaria del gobierno. Como vimos en [\authorfont{Ver Tema 3.B.9}], el deseo de una sociedad de lograr una distribución más equitativa de la renta puede llevar a un mayor nivel deseado de impuestos, pero ese mayor nivel impositivo puede tener un efecto perjudicial sobre la economía: unos impuestos más altos pueden desincentivar la obtención de ingresos por parte de las personas y reducir el tamaño del pastel económico sobre el que pueden recaer esos impuestos. Que esta reducción de la base imponible llegue realmente a producirse depende de cuán sensible sea el tamaño de ese pastel a los impuestos que gravan las distintas actividades económicas.

Con el objetivo de comprender cómo afectan los tipos impositivos a la economía, abordaremos la cuestión de hasta qué punto las decisiones económicas individuales responden a la tributación. 

A continuación, estudiaremos la evidencia empírica sobre el impacto de la tributación en estos comportamientos. Finalmente, en cada caso, analizaremos las principales políticas tributarias vigentes en Estados Unidos que afectan a estas conductas.

\subsection{Aproximación microeconómica: Incidencia sobre las decisiones óptimas}
Analizaremos tres tipos de impuestos: los impuestos sobre el trabajo, los impuestos sobre el ahorro y los impuestos sobre la riqueza y la asunción de riesgos. En cada caso, comenzaremos examinando la teoría sobre cómo la tributación puede afectar a las decisiones individuales, como cuánto trabajar, cuánto ahorrar y cuánto riesgo asumir en las inversiones. 

A lo largo de este apartado se analizará cómo los impuestos modifican las decisiones de oferta de los agentes económicos ya que distintos impuestos tendrán diferentes efectos económicos sobre las decisiones óptimas de consumo, ahorro, inversión, producción o trabajo.

\textcolor{red}{Podemos distinguir dos tipos de impuestos: por un lado, los impuestos distorsionantes que alteran el comportamiento de los individuos al generar una variación de los precios relativos y, por otro lado, los impuestos neutrales o no distorsionantes, como los impuestos de suma fija, que no alteran  las decisiones individuales. Esto último sólo será posible cuando el impuesto sea independiente  de las decisiones de consumo o de las ofertas factoriales de los individuos lo cual incumple las  nociones de equidad aceptadas en las sociedades avanzadas, por lo que en la práctica la mayoría de los impuestos son del primer tipo.

Así. partimos de una situación "second best'" en tanto la estructura impositiva dará lugar a distorsiones en la economía y, en consecuencia, podría suponer una pérdida de eficiencia. La teoría del ·'second best'" (Lipsey y Lancaster, 1956) ya señalaba que en aquellos casos en que se incumple una de las condiciones de eficiencia exigidas por el criterio de Pareto puede no ser necesario y al mismo tiempo no ser deseable que se satisfagan las demás condiciones para alcanzar el estado económico que maximice el bienestar social. En consecuencia. dado que el Estado necesita ingresos públicos la pregunta esencial es qué distorsiones introducir en la economía para alcanzar tal fin.

A continuación, asumimos dos tipos de agentes: hogares (los cuales son demandantes de bienes y oferentes de trabajo) y las empresas (las cuales son oferentes de bienes y demandantes de trabajo).}

Este capítulo, el segundo de nuestra serie de tres capítulos sobre tributación y comportamiento individual, se centra en las decisiones de ahorro. 


\subsubsection{Oferta de trabajo: ¿Cómo afectan los impuestos sobre la renta salarial?}
Comencemos con un análisis de la tributación de los ingresos obtenidos por el trabajo, lo que los economistas denominan generalmente oferta de trabajo (\textit{labor supply}). 

Si los individuos reducen considerablemente su esfuerzo de trabajo a medida que aumentan los impuestos sobre la renta, los impuestos sobre la renta podrían imponer grandes pérdidas de peso muerto a la sociedad. Esto ilustra, de manera simple, la importancia de comprender cómo responde la oferta de trabajo a los cambios en la tributación.

El marco teórico para evaluar cómo los impuestos sobre la renta afectan a la oferta de trabajo es muy similar al utilizado en cualquier modelo de Economía de la Imposición siguiendo un enfoque parcial.

\paragraph{Trabajadores: Impuesto de la Renta de las Personas Físicas}
La Figura 21-2 ilustra las posibles elecciones de Ava entre horas de ocio y€ de consumo. Recuérdese que resolvemos para el ocio óptimo, y después calculamos el trabajo como el total de horas posibles por año menos las horas de ocio. Recuérdese también que antes de impuestos la pendiente de la restricción presupuestaria, la tasa salarial, es el precio del ocio porque es el coste de oportunidad de tomar ocio en lugar de trabajar. Antes de que se impongan los impuestos, Ava disfruta de un nivel inicial de ocio de $900$ y un nivel de consumo de $C1 = 13.750€$ (punto A). Supongamos ahora que se impone un impuesto $t$ del $30\%$ sobre cada dólar de salarios ganado.

La pendiente de la restricción presupuestaria es ahora el salario después de impuestos, $12,50\cdot(1 − 0,3) = 8,75€$, porque esto es lo que ella (y todos los trabajadores con el mismo salario bruto) ven como el rendimiento monetario de su trabajo. Un impuesto de tipo único sobre los ingresos del trabajo a una tasa $t = 0,3$ hace que la restricción presupuestaria pivote hacia dentro desde $BC1$ hasta $BC2$. Con la misma cantidad de trabajo, $900$ horas, y por tanto la misma cantidad de ocio, puede ahora consumir menos bienes, $C2 = 9.625$, porque parte de sus ingresos se destina a pagos de impuestos.

Efectos sustitución y renta sobre la oferta de trabajo

No podemos decir con certeza qué ocurrirá con la oferta de trabajo como resultado de este impuesto porque tiene dos efectos compensatorios. El salario después de impuestos es el precio efectivo del ocio. Debido a que el salario después de impuestos es menor que el salario antes de impuestos, el precio del ocio ha disminuido. La disminución del precio del ocio inducirá un efecto sustitución hacia más ocio y menos trabajo. Sin embargo, la disminución de los rendimientos del trabajo también significa que es más pobre para cualquier nivel dado de oferta de trabajo. Esta reducción de la renta tendrá un efecto renta que hace que compre menos de todos los bienes normales, incluido el ocio; y menos horas de ocio significan más horas de trabajo.\footnote{%
    Para comprender la intuición del efecto renta sobre la oferta de trabajo, a veces resulta útil pensar en el objetivo de renta de un individuo, su objetivo de ganar una cantidad fija de renta. Imaginemos que la única razón por la que trabaja es para comprar un CD cada semana. Si gana $5€$ por hora, y un CD cuesta $20€$, trabajará cuatro horas cada semana. Si el gobierno impone un impuesto del $20\%$ sobre los ingresos del trabajo, su salario después de impuestos caerá a $4€$. Para comprar ese mismo CD, tendrá ahora que trabajar cinco horas por semana. Así, trabaja más, aunque su salario después de impuestos haya caído, porque tiene una renta objetivo que quiere ganar. En este caso, el efecto renta domina al efecto sustitución, y la cantidad de trabajo ofrecida aumenta.
} Debido a que los efectos sustitución y renta sobre la oferta de trabajo tiran en direcciones opuestas, no podemos predecir claramente si la oferta de trabajo aumenta o disminuye en respuesta al impuesto $t$.

Los dos paneles de la Figura 21-3 ilustran dos posibles efectos de gravar la renta del trabajo. En el panel (a), el efecto sustitución de la tributación (que reduce el precio del ocio, llevando a Ava a desear más ocio) es mayor que el efecto renta de una menor renta después de impuestos (que conduce a menos ocio). En este caso, el ocio de Ava aumenta de 900 horas a 1.200 horas, lo que implica una menor oferta de trabajo. En el panel (b), el efecto renta de la menor renta después de impuestos es mayor que el efecto sustitución de la tributación, y trabaja más para ganar más renta de modo que pueda permitirse más consumo. En este caso, el ocio de Ava cae de $900$ a $600$, lo que implica una mayor oferta de trabajo.

\imagenfit{Fig2.png}{¿Qué efecto domina? Dos supuestos: Efecto renta vs. Efecto sustitución}{\textit{Public Finance}. GRUBER (20XX)}

Estas dos posibilidades implican formas diferentes para la curva de oferta en el mercado de trabajo. Si dominan los efectos sustitución, como en el primer caso, entonces la curva de oferta de trabajo tiene la típica forma con pendiente ascendente que hemos discutido hasta ahora. Si dominan los efectos renta, como en el segundo caso, entonces la curva de oferta de trabajo tendrá pendiente descendente, con salarios más altos llevando a los individuos a ofrecer una menor cantidad de trabajo. Parece muy improbable que las curvas de oferta de trabajo puedan tener pendiente descendente en todas partes porque los efectos renta sobre la oferta de trabajo son proporcionales a las horas trabajadas antes del cambio salarial. Si los individuos no han estado trabajando en absoluto, y los salarios son gravados, entonces hay un efecto sustitución sobre su decisión de oferta de trabajo, pero ningún efecto renta: no pueden ser más pobres porque no estaban ganando nada antes de la tributación. Así, a niveles bajos de oferta de trabajo, parece improbable que los efectos renta puedan ser mayores que los efectos sustitución. A niveles más altos de oferta de trabajo, donde hay una mayor pérdida de renta por la tributación de los salarios, los efectos renta podrían llegar a ser mayores que los efectos sustitución, y la tributación podría conducir a más, no menos, oferta de trabajo.

\paragraph{Mercado de trabajo: Cotización a la Seguridad Social}
Considérese el impuesto sobre las nóminas utilizado para financiar el sistema de la Seguridad Social. Como se señaló en el Capítulo 11, un impuesto igual al 7,65 por ciento de los ingresos de los trabajadores debe ser pagado por sus empleadores y un impuesto al mismo tipo pagado por los propios trabajadores—un total del 15,3 por ciento. Esta división tiene una larga historia y es una consecuencia de la creencia de nuestros legisladores de que el impuesto sobre las nóminas debería ser compartido por igual por empleadores y empleados. Pero la distinción legal entre trabajadores y jefes es irrelevante. Como se sugirió anteriormente, la incidencia de este impuesto sobre el trabajo está determinada únicamente por la cuña que el impuesto introduce entre lo que reciben los empleados y lo que pagan los empleadores. Este punto se ilustra en el gráfico siguiente: 

\imagenfit{Fig3.png}{Incidencia de un impuesto sobre la nómina: Oferta inelástica}{ROSEN Y GRABER (2011)}

Señalamos $D_L$ como la demanda de trabajo y $S_L$ como la oferta de trabajo, siendo esta perfectamente inelástica (por motivos ilustrativos). Antes de la tributación, el salario es $w_0$. El impuesto ad valorem sobre el trabajo desplaza la curva de demanda efectiva a $D′_L$. Como de costumbre, la distancia entre $D′_L$ y $D_L$ es la cuña entre lo que se paga por un bien y lo que reciben quienes lo ofrecen. Después de que se imponga el impuesto, el salario recibido por los trabajadores cae a $w_n$. Por otro lado, $w_g$, el precio pagado por los empleadores, permanece en $w_0$. En este ejemplo, a pesar de la división legal del impuesto, la tasa salarial recibida por los trabajadores cae exactamente en la cuantía del impuesto: los trabajadores soportan toda la carga.

Por supuesto, podríamos haber obtenido justamente el resultado opuesto dibujando la curva de oferta como perfectamente elástica. El punto clave que hay que recordar es que no se puede saber nada acerca de la incidencia de un impuesto sin información sobre las elasticidades de comportamiento relevantes. De hecho, aunque las estimaciones de la elasticidad de la oferta de trabajo varían, muchos economistas creen que está cerca de cero (FUCHS et al., 1998). Óscar menos a corto plazo, el trabajo probablemente soporta la mayor parte del impuesto sobre las nóminas, a pesar de los intentos de dividir la carga.

\paragraph{Limitaciones comunes: Restricciones y normas laborales}
La teoría básica del efecto de la tributación sobre la oferta de trabajo supone una visión idealizada del mercado de trabajo, donde los individuos pueden ajustar sus horas libremente e incrementalmente a medida que cambia la política tributaria del gobierno. En la mayoría de los mercados de trabajo, sin embargo, los individuos no pueden ajustar libremente sus horas de trabajo para encontrar la tangencia exacta de su curva de indiferencia y su restricción presupuestaria. Las empresas pueden, por ejemplo, exigir a los trabajadores que estén en el trabajo durante un determinado número de horas.

Esta restricción podría deberse a complementariedades de producción, características del proceso de producción que hacen importante tener a muchos trabajadores en el trabajo al mismo tiempo. Un trabajador que forma parte de una línea de producción en una planta manufacturera no puede trabajar solamente $32$ horas por semana, incluso si ese es su óptimo, cuando todos los demás trabajadores de la línea están trabajando $40$ horas por semana.

Otra restricción al aumento de las horas de los trabajadores es la existencia de reglas de pago de horas extraordinarias, que obligan a que a los trabajadores pagados por hora se les pague una vez y media su remuneración horaria regular si trabajan más de $40$ horas por semana. Estas leyes hacen que sea muy caro para las empresas hacer que los trabajadores trabajen más de $40$ horas por semana, lo que significa que serán reacias a permitir que los trabajadores lo hagan incluso si ese es el óptimo para el trabajador. En conjunto, tales restricciones de horas forzarán a los trabajadores hacia horarios de trabajo uniformes, reduciendo así la capacidad de respuesta de las horas trabajadas a los salarios después de impuestos.

\subsubsection{Oferta de ahorro}
La decisión sobre cuánto de la propia renta ahorrar y cuánto gastar, al igual que la decisión sobre cuánto trabajo ofrecer, es otro comportamiento individual importante que se ve afectado por la tributación. En esta sección, analizamos cómo los impuestos podrían afectar a las decisiones de ahorro.

El 31 de enero de 2003, la administración Bush propuso la creación de Cuentas de Ahorro de por Vida (Lifetime Savings Accounts) y Cuentas de Ahorro para la Jubilación (Retirement Savings Accounts), que permitirían a las familias colocar hasta 30.000€ cada año en cuentas especiales de ahorro sobre las cuales los intereses obtenidos no estarían sujetos al impuesto sobre la renta. Estas nuevas cuentas presentaban menos restricciones que los instrumentos de ahorro existentes subsidiados fiscalmente: los individuos podían retirar dinero de su Cuenta de Ahorro de por Vida por cualquier motivo, a cualquier edad, y de su Cuenta de Ahorro para la Jubilación por cualquier motivo después de los 58 años. Cuando se retiraran los fondos, cualquier interés obtenido en estas cuentas, o los aumentos en el valor de los activos de las cuentas, estarían exentos de tributación. Los opositores a los nuevos planes afirmaban que las nuevas cuentas drenarían ingresos fiscales (potencialmente enormes cantidades) del gobierno en el futuro, cuando el dinero fuera retirado y no se pagaran impuestos. Además, muchos opositores, como el representante demócrata Charles Rangel, argumentaban que la nueva desgravación fiscal sería utilizada por personas con medios para proteger de los impuestos lo que ya ahorrarán, de modo que el ahorro en realidad no aumentaría. Políticamente, la mayoría de los demócratas y muchos republicanos coincidían con la opinión de Rangel de que los costes de ingresos de tal plan superarían los beneficios obtenidos de cualquier posible aumento del ahorro. Tras una semana del anuncio, la administración Bush se dio cuenta de que las nuevas cuentas tenían muy poco apoyo político y comenzó a centrarse en otras medidas de reducción de impuestos.

El presidente Barack Obama adoptó un enfoque muy diferente el 5 de septiembre de 2009, cuando anunció su plan para ayudar a los estadounidenses a ahorrar para la jubilación. Consistía en iniciativas basadas en investigaciones de economía conductual que han mostrado que los empujones sistemáticos son muy eficaces para aumentar la inversión de un individuo en ahorro para la jubilación.\footnote{%
    La primera iniciativa animaba a las pequeñas empresas a inscribir automáticamente a sus empleados en cuentas de ahorro para la jubilación, como los 401(k) o las Cuentas Individuales de Jubilación simples (\textit{Individual Retirement Accounts}, IRA), reduciendo las costosas comisiones que anteriormente habían desalentado el establecimiento de planes de ahorro de pequeñas empresas. Además, también se animaría a los empleadores a aumentar automáticamente la tasa de ahorro de un empleado cada año o cada vez que el empleado recibiera un aumento salarial. En el momento de escribir esto, el Congreso no ha aprobado ninguna legislación sobre cuentas de jubilación obligatorias a nivel federal, pero algunos estados han asumido por sí mismos la obligación de exigir a los empleadores que inscriban automáticamente a sus empleados en tales planes. La segunda iniciativa, sobre la que posteriormente actuó el IRS, permite a los trabajadores recibir cualquier cantidad de su devolución de impuestos en forma de bonos de ahorro de Estados Unidos, en lugar de en efectivo. La tercera iniciativa facilitaría a los trabajadores depositar en sus cuentas de ahorro los salarios obtenidos por horas extraordinarias o el tiempo de vacaciones no utilizado. El IRS dictaminó que esto está permitido, pero que quedaba a discreción del empleador ofrecer este servicio a sus empleados. Por último, el plan pedía la creación de una guía de ahorro fácil de entender para empleadores y empleados que detallara las opciones disponibles para el ahorro para la jubilación con tratamiento fiscal favorable. En respuesta, el IRS creó una versión más simple y más fácil de usar de su sitio web de planificación de la jubilación.
} Aunque la administración Obama tuvo cierto éxito, los críticos han culpado al plan de no abordar los principales problemas del ahorro para la jubilación, incluida la insuficiente financiación de los planes de jubilación estatales y locales.

Estos enfoques contrapuestos para fomentar el ahorro para la jubilación plantean una serie de preguntas. ¿Reduce la estructura existente del impuesto sobre la renta en Estados Unidos la cantidad que los individuos ahorran en este país? ¿Debería el gobierno utilizar el código fiscal para animar a los estadounidenses a ahorrar más? ¿O debería el gobierno, en cambio, basarse en cambios en la estructura de los planes de ahorro para la jubilación que proporcionen incentivos conductuales para que las personas ahorren más?

Como sabemos, la cantidad que ahorra una sociedad, que hace que haya más capital disponible para las empresas, es un determinante clave del crecimiento económico. Así, una fuente importante de debates de política en todo el mundo es el papel apropiado de la tributación de la renta del capital, los impuestos aplicados a los rendimientos del ahorro.

Comenzamos el apartado anterior analizando la teoría tradicional de cómo la tributación sobre la renta afecta a la oferta de trabajo. De manera similar, comencemos con una discusión de la teoría tradicional de cómo la tributación sobre la renta afecta a las decisiones de ahorro. Sin embargo, a diferencia de nuestros resultados sobre la oferta de trabajo, existe relativamente poca evidencia sobre la cuestión clave planteada por este modelo: ¿cuán sensible es el ahorro al tipo de interés después de impuestos? A continuación, pasamos a dos modelos alternativos de ahorro, el modelo precautorio (en el que el ahorro sirve principalmente como autoaseguro frente al riesgo) y el modelo de autocontrol (en el que el ahorro se determina mediante una competencia entre la impaciencia a corto plazo y la paciencia a largo plazo).

\paragraph{Modelo elección intetemporal: Suavización del consumo}
En la teoría tradicional del ahorro, el papel del ahorro es suavizar el consumo entre períodos de tiempo. Los individuos tienen más renta en algunos períodos (como durante sus vidas laborales) y menos en otros (como durante la jubilación). Sin ahorro, los individuos se verían obligados a consumir mucho menos en los períodos en los que la renta es menor. Sin embargo, debido a la utilidad marginal decreciente, es poco probable que tal patrón de consumo maximice la utilidad: los individuos preferirían suavizar su consumo a lo largo del tiempo, consumiendo una cantidad constante en lugar de darse un festín en algunos períodos y pasar hambre en otros. Óscar reducir el gasto en los períodos en los que es alto y aumentar el gasto en los períodos en los que la renta es baja, el ahorro desempeña ese papel de suavización del consumo.

Bajo este modelo, llamado modelo de elección intertemporal, la elección sobre cuánto ahorrar es realmente la elección sobre cómo asignar el propio consumo a lo largo del tiempo. Óscar igual que al modelizar la oferta de trabajo, no modelizamos directamente el ahorro, sino que modelizamos el consumo a lo largo del tiempo. Debido a que el trabajo es un mal, analizamos el bien complementario, el ocio, y después calculamos la oferta de trabajo como la diferencia entre el tiempo disponible y el tiempo de ocio tomado. De manera similar, el ahorro es un mal: es algo que hacemos hoy para financiar el consumo futuro, pero no obtenemos ninguna utilidad directa del ahorro. (Óscar menos no en los modelos estándar; algunos economistas han propuesto modelos alternativos de ahorro que sugieren que los individuos obtienen realmente utilidad de tener ahorros). El bien correspondiente es el consumo de hoy. El ahorro es la diferencia entre la renta de una persona y su consumo actual.

Un modelo simplificado

Supongamos que Óscar vive durante dos períodos. El período uno es su vida laboral, durante la cual obtiene una renta Y; el período dos es su jubilación, durante la cual no obtiene nada. Su consumo durante su vida laboral, $C_W$, es su renta en el período uno menos cualquier ahorro ($S$) que realice. Esos ahorros se depositan en el banco y obtienen un tipo de interés $r$. En el período dos, la jubilación, el consumo de Óscar es su ahorro del período uno más los intereses obtenidos sobre esos ahorros, $C_R=S \times (1+r)$. Dada esta información, Óscar elige su nivel óptimo de consumo en ambos períodos. Gráficamente:

\imagenfit{Fig5.png}{Imposición y la decisión intertemporal de consumo}{\textit{Public Finance and Policy Making}. GUBER (2016)}

Analizando la elección de Óscar entre el consumo en los períodos uno ($C_W$) y dos ($C_R$), podemos derivar el ahorro utilizando el modelo de elección intertemporal. Como podemos intuir, las preferencias de Óscar respecto al consumo actual frente al futuro, junto con su deseo de suavización del consumo a lo largo del tiempo, determinan la forma y la posición de sus curvas de indiferencia. Antes de cualquier tributación del ahorro, Óscar maximiza su utilidad sujeto a una restricción presupuestaria intertemporal ($BC_1$), que relaciona su consumo del primer y segundo período ($IC_1$) con sus decisiones de renta y ahorro. Si Óscar no ahorra nada, puede consumir $Y$ en el primer período. Si ahorra toda su renta, puede consumir $Y(1+r)$ en el segundo período. El precio relativo del consumo del primer período es, por tanto, $1+r$, la pendiente de la restricción presupuestaria. Esto se debe a que el coste de oportunidad del consumo del primer período son los intereses no obtenidos sobre el ahorro destinado al consumo del segundo período, más el consumo del segundo período al que se renuncia. Del mismo modo que el precio del ocio en el modelo de oferta de trabajo es el salario (el coste de oportunidad del ocio), el precio del consumo del primer período en este modelo de dos períodos es el coste de oportunidad del consumo del primer período, $1+r$.

Podemos considerar qué ocurre cuando el gobierno grava los pagos de intereses sobre el ahorro.\footnote{%
    \textbf{NOTA.-} Es importante recordar que \textbf{no estamos modelizando directamente el ahorro}: del mismo modo que modelizamos el ocio y resolvimos el trabajo como un residuo, aquí modelizamos el consumo mientras se trabaja, $C_W$, y resolvemos el ahorro como un residuo.
} Supongamos ahora que el gobierno grava la renta por intereses. Tal impuesto haría que el rendimiento después de impuestos del ahorro cayera de $r\to r\times(1-t)$. Esto haría más plana la restricción presupuestaria, pasando de $BC_1$ a $BC_2$, como podemos ver en la figura, reduciendo la pendiente de la restricción presupuestaria de $(1+r)\to [1+(r\times[1-t])]$.

Que la restricción presupuestaria sea más plana refleja el hecho de que el precio del consumo del primer período, $C_W$, ha caído: el coste de oportunidad de consumir en el primer período ha caído porque cada dólar de ahorro proporciona menos consumo en el segundo período. Para un nivel dado de consumo del primer período, Óscar puede permitirse ahora menos consumo del segundo período, ya que sus ahorros obtienen una tasa de rendimiento menor.

Como siempre, el cambio de precio que resulta del impuesto sobre los intereses del ahorro tendrá dos efectos. El menor tipo de interés después de impuestos (menor precio de $C_W$) hará que el consumo en el período uno aumente a través del efecto sustitución. Esto, a su vez, hará que el ahorro disminuya. Existe también, sin embargo, un efecto renta derivado de una menor renta después de impuestos. Óscar es ahora más pobre para todos los niveles de ahorro porque la cantidad de intereses que conserva por cada dólar de ahorro ha caído, y un nivel dado de ahorro le permite comprar menos $C_R$. Con el mismo nivel de consumo del primer período (y, por tanto, el mismo nivel de ahorro) de $C_{W1}$, Óscar solo puede permitirse un consumo del segundo período de $C_{R2}$. Esta caída de la renta hará que $C_{W1}$ disminuya y que $S$ aumente.\footnote{%
    Óscar pensar en el efecto renta de los cambios en el tipo de interés después de impuestos sobre el ahorro, es útil reflexionar sobre el caso extremo de un nivel objetivo de consumo para la jubilación en el período dos, como hicimos antes con los ingresos objetivo en el caso de la oferta de trabajo. Si Óscar tiene una determinada cantidad de consumo que quiere en el período dos (como un determinado nivel de vida durante su jubilación), entonces, cuando cae el tipo de interés después de impuestos, debe ahorrar más y reducir $C_W$ en el período uno para alcanzar ese objetivo.
} Debido a que los efectos sustitución y renta operan en direcciones opuestas, el efecto neto de la tributación de los intereses sobre el ahorro es incierto:

\imagenfit{Fig6.png}{¿Qué efecto domina? Dos supuestos: Efecto renta vs. Efecto sustitución}{\textit{Public Finance}. GRUBER (20XX)}

Si el efecto sustitución del menor tipo de interés después de impuestos (el menor precio del consumo del período uno, $C_W$, hace que Óscar prefiera más consumo del período uno) es mayor que el efecto renta de la menor renta después de impuestos (una renta menor hace que Óscar consuma menos de todo, incluido el consumo del período uno), tendríamos la situación que se ve en el panel (a). Óscar consume inicialmente en el punto $A=(C_{W1}, C_{R1})$. Después de que se imponga un impuesto sobre los intereses (que se observa en el desplazamiento de $BC_1 \to BC_2$), Óscar se desplaza al punto $B$, con mayor consumo del primer período, $C_{W2}$, menor ahorro (caída de $S_1 \to S_2$) y un consumo del segundo período mucho menor, $C_{R2}$.

Si el efecto renta de la menor renta después de impuestos es mayor que el efecto sustitución del menor tipo de interés después de impuestos, y Óscar se desplaza al punto $C$, situación que se observa en el panel (b). El consumo del primer período cae de $C_{W1}$ a $C_{W3}$, lo que implica un aumento del ahorro de $S_1 \to S_3$. El consumo del segundo período sigue cayendo en este caso (a $C_{R3}$), pero no tanto como en el panel (a). Esto se debe a que el efecto renta no es suficientemente grande como para llevar a Óscar a deshacer completamente la reducción de la renta ahorrando más; si $C_R$ no cambiara en absoluto, este sería el caso de renta objetivo descrito anteriormente, con el ahorro aumentando exactamente lo suficiente como para compensar la reducción del rendimiento del ahorro.\footnote{%
    \textbf{¡IMPORTANTE!} En este ejemplo simplificado, Óscar tiene ingresos solo en el primer período. En un modelo más general en el que Óscar tiene ingresos tanto en el primer como en el segundo período, y en el que puede pedir prestado contra los ingresos del segundo período para consumir en el primer período, existe otro efecto que refuerza aún más el efecto negativo de los impuestos sobre el ahorro: el efecto de riqueza humana destacado por SUMMERS (1981). El consumo de Óscar en ambos períodos en este modelo está determinado por el valor presente descontado de todo su flujo de renta a lo largo de la vida. Óscar descuenta sus ingresos laborales del segundo período utilizando el tipo de interés después de impuestos: desde la perspectiva de hoy, tener ingresos en el futuro vale menos con un tipo de interés más alto. El aumento de los impuestos reduce la tasa de descuento y eleva el valor presente neto de los ingresos laborales del segundo período. Con un mayor valor presente de los ingresos laborales a lo largo de la vida, Óscar consumirá más y ahorrará menos en respuesta a la mayor tributación.
}

\begin{specialblock}{\textbf{Caso ampliado.-} La inflación y la tributación del ahorro}
    Durante varios años de las décadas de 1970 y 1980, Estados Unidos tuvo que hacer frente a tasas de inflación muy elevadas, que alcanzaron un máximo del $13,5\%$ en 1980. Antes de 1981, los tramos impositivos sobre los que se basa la tributación estaban denominados en€ constantes que no cambiaban con la inflación. Esta práctica dio lugar a un fenómeno conocido como desplazamiento de tramo o progresividad en frío (\textit{bracket creep}), mediante el cual los individuos experimentaban un aumento de su tipo impositivo a pesar de no haber aumentado su renta real (en€ constantes).

    De 1979 a 1980, por ejemplo, los precios aumentaron un $11,3\%$. Si las rentas hubieran aumentado en esa misma proporción, los consumidores habrían mantenido ingresos constantes en términos reales, de modo que su poder adquisitivo no habría cambiado. Sin embargo, los importes de los tramos impositivos permanecieron iguales en 1979 y en 1980, lo que significaba que los individuos cuyas rentas aumentaban únicamente para compensar la inflación pagarían más impuestos, porque una mayor parte de su renta aparecería en tramos impositivos superiores.
    
    Supongamos, por ejemplo, que Óscar tenía una renta de $16.500€$ en 1979. Ese año, habría estado sujeto a un tipo marginal del $21\%$ y habría pagado 2.370€ en impuestos, quedándole una renta después de impuestos de 14.130€. En 1980, la renta de Óscar aumentó hasta 18.365€, compensando exactamente la inflación del $11,3\%$ de ese año, de modo que la renta real de Óscar (la cantidad de bienes que podía comprar con su renta) no cambió. Los tramos impositivos no estaban indexados en 1980, por lo que Óscar habría estado sujeto a un tipo marginal superior, del 24\%, porque el tramo del 21 % terminaba en una renta de 16.600€. Habría pagado 2.815€ en impuestos y habría tenido una renta después de impuestos de 15.550€. Obsérvese que la renta de Óscar después de impuestos aumentó solo un 10 %, mientras que los precios aumentaron un 11,3 %. Debido a los tramos impositivos nominales fijos, el aumento de su salario fue demasiado pequeño para mantener el ritmo de la inflación. La solución a este problema fue bastante sencilla: en 1981, el Congreso indexó los tramos del impuesto ordinario sobre la renta (aunque no los del Impuesto Mínimo Óscarternativo), de modo que el tipo marginal se basara en la renta real y no en la renta nominal.
    
    La inflación y la tributación del capital
    
    La indexación de los tramos del impuesto sobre la renta no eliminó completamente el impacto de la inflación sobre la tributación de la renta, porque las reglas relativas a la tributación de las rentas del capital permanecieron iguales. El tipo de interés obtenido sobre una cuenta bancaria está determinado por el tipo de interés nominal, mientras que la mejora real del poder adquisitivo derivada del ahorro está determinada por el tipo de interés real. En lugar de preocuparse por cuánto dinero recibirá el año próximo, usted debería preocuparse por cuántos bienes podrá consumir con ese dinero el año próximo.
    
    La Tabla 22-1 ilustra el impacto de la tributación del capital en un entorno inflacionario. Supongamos que Robin va a ahorrar 100€ a un tipo de interés nominal del $10\%$, y que gasta todo su dinero en caramelos Skittles, a un dólar por bolsa. Inicialmente, no existe ningún impuesto sobre los intereses obtenidos. Robin gana diez€ en intereses, de modo que sus recursos después de impuestos son $110€$. Con un precio de un dólar por bolsa de Skittles, sin inflación (primera fila de la tabla), puede comprar $110$ bolsas como resultado de este ahorro. En la segunda fila, se introduce un impuesto sobre la renta del capital del $50\%$, lo que significa que Robin conserva solo la mitad de los intereses obtenidos, es decir, cinco€. Por tanto, después de la tributación solo puede comprar $105$ bolsas de Skittles.
    
    \begin{center}
    \begin{tabular}{p{0.07\linewidth}p{0.07\linewidth}p{0.07\linewidth}p{0.07\linewidth}p{0.07\linewidth}p{0.07\linewidth}p{0.07\linewidth}p{0.07\linewidth}p{0.07\linewidth}}
    \hline
      {\scriptsize\textbf{Caso}} & {\scriptsize\textbf{Inflación}} & {\scriptsize\textbf{Impuesto ($r$)}} & {\scriptsize\textbf{Ahorro}} & {\scriptsize\textbf{Interés ($i$)}} & {\scriptsize\textbf{Ingresos $S\,(1+r)$}} & {\scriptsize\textbf{Recursos (after-tax)}} & {\scriptsize\textbf{Precio}} & {\scriptsize\textbf{Cantidad}} \\
    \hline
      {\scriptsize\multirow{2}{*}{\textbf{Sin inflación}}} & $0\%$ & $0\%$ & $100$ & $10\%$ & $10€$ & $110€$ & $1€$ & $110$ \\
       & $0\%$ & $50\%$ & $100$ & $10\%$ & $10€$ & $105€$ & $1€$ & $105$ \\
       {\scriptsize\multirow{2}{*}{\textbf{Inflación}}} & $10\%$ & $0\%$ & $100$ & $10\%$ & $10€$ & $110€$ & $1.1€$ & $100$ \\
       & $10\%$ & $50\%$ & $100$ & $10\%$ & $10€$ & $105€$ & $1.1€$ & $95.5$\\
    {\scriptsize\textbf{Constante}} & $10\%$ & $0\%$ & $100$ & $21\%$ & $21€$ & $121€$ & $1.1€$ & $110$ \\
       {\scriptsize\textbf{(real)}}& $10\%$ & $50\%$ & $100$ & $21\%$ & $21€$ & $110.5€$ & $1.1€$ & $100.5$ \\
    \hline
    \end{tabular}
    \end{center}
    
    Supongamos ahora que existe una inflación del $10\%$, de modo que una bolsa de Skittles cuesta $1,10€$. En este caso, sin tributación (tercera fila), Robin solo puede comprar $100$ bolsas con sus $110€$ de recursos, exactamente la misma cantidad que hoy. Su poder adquisitivo no ha aumentado porque sus ganancias por intereses del $10\%$ han sido compensadas por una inflación del $10\%$. Con una tributación de las rentas del capital del $50\%$, ahora solo dispone de $105$€ después de impuestos y solo puede comprar $95,5$ bolsas de Skittles.
    
    Cuando existe inflación, sin embargo, esta afectará al tipo de interés que Robin puede obtener por sus ahorros. Podemos definir la relación entre los tipos de interés real y nominal como: $\text{Tipo de interés real }(r)=\frac{1+\text{Tipo de interés nominal }(i)}{1+\text{Tasa de inflación }(p)}-1$
    
    Cuando los individuos toman sus decisiones de ahorro, les importa el tipo de interés real, o cuánto puede mejorar su poder adquisitivo gracias a sus ahorros (en este ejemplo, cuántas bolsas de Skittles podrán comprar el año siguiente). Como consecuencia, cuando aumenta la inflación, los bancos y las empresas tendrán que pagar un tipo de interés nominal más elevado para atraer el ahorro de los individuos.
    
    En principio, los bancos y las empresas compensarán la inflación aumentando los tipos de interés nominales para mantener constante el nivel de los tipos reales. Si la inflación fuera del $10\%$, un banco pagaría un tipo de interés nominal del $21\%$ para proporcionar el mismo rendimiento real del $10\%$ que se obtenía antes de la inflación (porque $121€/1,10€ = 110$ bolsas de Skittles). Como muestra la tercera fila de la tabla, con un tipo de interés del $21\%$ y sin tributación, Robin puede volver a comprar $110$ bolsas de Skittles a un precio de $1,10€$ (como en la primera fila). Así, si los bancos aumentan los tipos de interés nominales lo suficiente para compensar completamente la inflación, en ausencia de tributación la inflación no erosionará el poder adquisitivo del ahorro.
    
    El problema es que los impuestos se aplican sobre los intereses nominales, no sobre los intereses reales. La última fila reintroduce los impuestos sobre los intereses obtenidos. Ahora, cuando Robin gana 21€ en intereses, paga 10,50€ en impuestos sobre esos intereses, lo que le deja con 110,50€ de recursos después de impuestos. Con esos 110,50€ solo puede comprar 100,5 bolsas de Skittles, una cantidad inferior a la que podía comprar después de impuestos antes de la inflación (105 bolsas en la segunda fila).
    
    Debido a que los impuestos se aplican sobre los tipos de interés nominales, el aumento de los tipos de interés nominales por parte de los bancos durante los períodos de inflación no compensará completamente la disminución de los tipos de interés reales. Así, aunque el desplazamiento de tramo (bracket creep) terminó en 1981, los efectos de la inflación sobre el sistema tributario siguen siendo importantes. En particular, una inflación más elevada reduce el rendimiento real del ahorro después de impuestos.
\end{specialblock}


En el modelo estándar de elección intertemporal, el ahorro se lleva a cabo únicamente para suavizar el consumo a lo largo del tiempo. La investigación económica reciente se ha centrado en dos ampliaciones de este modelo que sugieren otros factores que podrían ser importantes para determinar el ahorro, y que el tipo de interés después de impuestos puede ser menos importante de lo que propone la teoría tradicional.

\paragraph{Modelos precautorios: Incertidumbre}
Otro factor que determinará el ahorro es la cantidad de incertidumbre a la que se enfrentan los individuos en sus perspectivas financieras y su deseo de utilizar el ahorro para asegurarse contra shocks financieros negativos. Del mismo modo que el ahorro puede suavizar el consumo a lo largo del tiempo, también puede suavizar el consumo entre estados del mundo, como destacó nuestra discusión del autoaseguro en el Capítulo 12. De hecho, cuando se preguntó a las personas por sus razones para ahorrar, la respuesta principal, junto con ahorrar para la jubilación, fue ahorrar para «emergencias», como el desempleo o los problemas de salud.⁷ Este resultado ha llevado al desarrollo del modelo de ahorro precautorio, en el que el ahorro está motivado por el deseo no solo de suavizar el consumo a lo largo del tiempo, sino también de autoasegurarse contra el riesgo.

En este modelo, las personas se enfrentan al riesgo de que ocurran acontecimientos adversos en el futuro: shocks adversos de salud (como sufrir un ataque al corazón), desempleo, divorcio, etcétera. Este modelo supone que los individuos son incapaces de pedir prestado si experimentan un shock adverso porque se enfrentan a restricciones de liquidez, barreras al endeudamiento. Tales restricciones de liquidez pueden surgir, por ejemplo, porque los bancos no están dispuestos a prestar a un individuo con una enfermedad grave o que acaba de perder su empleo. Como resultado, los individuos deben acumular una reserva de ahorro que esté disponible en caso de que ocurra uno de estos shocks adversos, de modo que puedan suavizar el consumo entre períodos de buena suerte y períodos de shocks adversos. El deseo de tener una reserva es otra razón, además de la suavización intertemporal, por la que los individuos tienen ahorros.



\paragraph{Modelos de autocontrol: Impaciencia}

Una formulación alternativa de la decisión de ahorro se basa en el tipo de modelos de autocontrol analizados en el Capítulo 6 en el contexto del tabaquismo. En este modelo, los individuos se enfrentaban al conflicto entre preferencias impacientes de corto plazo («necesito un cigarrillo hoy») y preferencias pacientes de largo plazo («me gustaría dejarlo mañana»). Es probable que un problema de autocontrol de este tipo también sea importante en el contexto del ahorro. Los individuos tienen una preferencia de largo plazo por asegurar suficiente ahorro para suavizar el consumo durante toda su vida, pero sus preferencias impacientes de corto plazo pueden hacer que consuman toda su renta y no ahorren lo suficiente para períodos futuros. En este modelo, un determinante clave del comportamiento de ahorro es la capacidad de los individuos para encontrar maneras de comprometerse a ahorrar, de modo que puedan mantener su renta fuera de las manos de su «yo de corto plazo» impaciente.


Incentivos fiscales para el ahorro para la jubilación

Óscargunos economistas y responsables de políticas sostienen que ahorramos demasiado poco en Estados Unidos y que nuestro crecimiento económico sufre como resultado. Además, a pesar de la existencia del programa de la Seguridad Social, algunos siguen preocupados por que la falta de visión de futuro de los trabajadores pueda hacer que estos ahorren demasiado poco para la jubilación. Como resultado de estas preocupaciones, el gobierno de Estados Unidos ha introducido una serie de subsidios fiscales para fomentar el ahorro para la jubilación. En esta sección, revisamos la estructura y los efectos de estos subsidios.

Subsidios fiscales disponibles para el ahorro para la jubilación

En Estados Unidos, existen diversos mecanismos subsidiados fiscalmente disponibles para el ahorro para la jubilación. En esta sección, revisamos brevemente los cuatro principales incentivos fiscales para el ahorro para la jubilación.

Subsidio fiscal a las pensiones proporcionadas por el empleador

Uno de los mayores beneficios complementarios proporcionados por los empleadores a sus empleados es el plan de pensiones, mediante el cual los empleadores ahorran para proporcionar ingresos de jubilación a sus trabajadores.

Tradicionalmente, los planes de pensiones proporcionados por el empleador eran planes de pensiones de prestación definida, en los que los trabajadores acumulaban derechos de pensión durante su permanencia en la empresa y, cuando se jubilaban, la empresa les pagaba una prestación que era función de la permanencia de ese trabajador en la empresa y de sus ingresos. Con el tiempo, los planes de pensiones proporcionados por el empleador han pasado a ser planes de pensiones de aportación definida, en los que los empleadores reservan una determinada proporción de los ingresos de un trabajador (como el 5 %) en una cuenta de inversión, y el trabajador recibe estos ahorros y cualquier rendimiento de inversión acumulado cuando se jubila.

De manera similar al seguro de salud proporcionado por el empleador, las aportaciones que los empleadores realizan a los planes de pensiones no se gravan como renta de los empleados. Del mismo modo, cualquier interés que se acumule sobre este ahorro para la pensión no se grava a medida que se acumula. En cambio, los empleados tributan por sus ahorros de pensión como renta ordinaria cuando se retiran durante la jubilación.

Cuentas 401(k)

La forma de ahorro para la jubilación que crece más rápidamente es la cuenta 401(k), un programa de ahorro controlado individualmente ofrecido a través del lugar de trabajo. Estas cuentas 401(k) permiten a los individuos ahorrar para su jubilación con un tratamiento fiscal favorable mediante una deducción de la nómina, y los empleadores a menudo igualan las aportaciones de los empleados.

Una opción 401(k) típica en una empresa permite al trabajador aportar hasta el 10 % de su renta a una cuenta de jubilación, y cualquier dólar aportado no se cuenta como renta imponible. Además, el empleador iguala, por ejemplo, el primer 5 % de la renta del empleado que este aporta, contribuyendo a la cuenta con una cantidad igual a la aportación del trabajador. Existe un límite a las aportaciones a los 401(k) de 18.000€ (en 2015). Cuando se retiran durante la jubilación, los saldos de las cuentas 401(k) se gravan como renta ordinaria.

Cuentas Individuales de Jubilación

El problema de las pensiones proporcionadas por el empleador o de los planes 401(k) como vehículos de ahorro para la jubilación es que a muchos individuos sus empleadores no les ofrecen tales planes; solo el 48 % de los trabajadores del sector privado participa en un plan de pensiones proporcionado por el empleador.⁹

En 1974, el Congreso introdujo la Cuenta Individual de Jubilación (Individual Retirement Account, IRA), un vehículo de ahorro para la jubilación con tratamiento fiscal favorable para individuos no cubiertos por pensiones proporcionadas por el empleador. 

El subsidio fiscal del ahorro a través de una IRA se amplió en 2001 mediante la introducción del «crédito fiscal del ahorrador» (saver’s tax credit). Este programa iguala las aportaciones a cuentas IRA realizadas por familias de bajos ingresos. Por ejemplo, para matrimonios que presentan conjuntamente la declaración de impuestos con ingresos inferiores a 36.500€, el crédito del ahorrador equivale al 50 % de su aportación a una IRA hasta 2.000€, de modo que el coste después de impuestos de aportar 2.000€ a una IRA es de solo 1.000€. La tasa de igualación disminuye después, y para declarantes conjuntos con ingresos entre 39.501 y 61.000€, la tasa del crédito es de solo el 10 %. El crédito del ahorrador también puede aplicarse a las aportaciones a un 401(k), bajo una fórmula algo más complicada.

IRA de Pensión Simplificada para Empleados

La última gran forma de ahorro para la jubilación subsidiado fiscalmente es la IRA de Pensión Simplificada para Empleados (Simplified Employee Pension IRA, SEP IRA), que es una opción de ahorro para la jubilación específicamente para los trabajadores por cuenta propia.

Los individuos con cuentas SEP-IRA pueden ahorrar hasta 53.000€ al año procedentes de sus ingresos por trabajo por cuenta propia libres de impuestos, para su retirada (y tributación) en el momento de la jubilación. Así, las cuentas SEP-IRA funcionan de la misma manera que las cuentas 401(k) (sin igualación), excepto que no se gestionan a través de empleadores.

¿Por qué los subsidios fiscales elevan el rendimiento del ahorro?

Todos los subsidios fiscales que se acaban de describir tienen una estructura similar: los individuos están protegidos de la tributación sobre sus ahorros, así como sobre cualquier ingreso por intereses obtenido sobre esos ahorros. Además, los individuos tributan por sus ahorros para la jubilación como renta ordinaria cuando retiran los fondos de sus cuentas de ahorro para la jubilación.

Si los ahorros para la jubilación se gravan de todos modos cuando los individuos están jubilados, ¿cómo constituye esto un subsidio fiscal? Esto es un subsidio porque, en lugar de pagar los impuestos sobre los ahorros por adelantado, usted difiere el pago de impuestos (tanto sobre la aportación inicial como sobre cualquier interés obtenido) hasta que retira esos ahorros para la jubilación. Recuérdese de nuestra discusión del valor presente descontado (VPD) en el Capítulo 4 que el dinero recibido hoy vale más que el dinero que se recibirá en el futuro porque usted puede obtener intereses sobre el dinero si lo tiene hoy. Del mismo modo, los impuestos pagados en el futuro son menos costosos que los impuestos pagados hoy porque usted puede obtener intereses sobre los pagos de impuestos que evita hoy.

En otras palabras, con ahorros que no tienen un tratamiento fiscal preferente, usted paga impuestos a medida que gana el dinero. A su vez, el gobierno puede tomar estos pagos de impuestos, depositarlos y obtener intereses sobre ellos. Con ahorros para la jubilación con tratamiento fiscal preferente, usted puede conservar cualesquiera impuestos que habría pagado tanto sobre su aportación inicial como sobre cualquier ingreso por intereses, y puede obtener los intereses sobre el dinero que, de otro modo, habría sido pagado en impuestos. La diferencia entre pagar los impuestos por adelantado (y permitir que el gobierno obtenga los intereses sobre ellos) y pagar los impuestos en el momento de la retirada (de modo que usted obtiene los intereses sobre ellos) puede ser bastante grande.

Supongamos que Ted es una persona de 70 años que ha ganado 100€ en su trabajo y que quiere ahorrar durante un año en una cuenta bancaria y después retirarlos. Su tipo impositivo es del 25 %, y el tipo de interés que puede obtener en la cuenta bancaria es del 10 %. Está intentando decidir si denominar su cuenta bancaria como una IRA; las implicaciones de su decisión se muestran en la Tabla 22-2. Si no denomina la cuenta como una IRA, tiene que pagar impuestos sobre sus ingresos. En ese caso, conserva 75€ de ingresos después de impuestos, los deposita en el banco, obtiene 7,50€ de intereses, pero después paga 1,88€ en impuestos sobre esos intereses. Cuando retira su dinero después de un año, tiene 80,62€. Sin embargo, si denomina su cuenta como una IRA, puede invertir los 100€ completos de ingresos (porque las aportaciones a una IRA son deducibles fiscalmente). Obtiene 10€ de intereses sobre esta inversión. Después, cuando retira el dinero tras un año, el gobierno recauda el 25 % de sus 110€, o 27,50€. Termina con 82,50€, 1,88€ más que la cantidad retirada de la cuenta que no es IRA.

Además, la ventaja de la IRA de diferir los pagos de impuestos aumenta cuanto más tiempo se mantiene el activo. Por ejemplo, si usted mantiene un activo durante 30 años a un tipo de interés del 10 % y un tipo impositivo del 25 %, una IRA le dejaría con el doble de dinero que una cuenta que no fuera IRA. (De hecho, incluso las IRA no deducibles disponibles para los contribuyentes de renta más alta después de 1986 son oportunidades valiosas para los inversores a largo plazo porque permiten a los inversores diferir la tributación sobre los ingresos por intereses hasta que se retira la inversión. De hecho, a largo plazo, la mayor parte de los beneficios de una IRA no procede de la deducibilidad inicial de las aportaciones a la IRA, sino de la acumulación libre de impuestos de los intereses obtenidos por la cuenta). Otra ventaja es que muchos contribuyentes se encontrarán en un tramo impositivo más bajo cuando se jubilen porque su renta es menor que cuando estaban trabajando. Como resultado, diferir los impuestos hasta la jubilación reduce la cantidad de impuestos pagados.

Esta misma lógica se aplica a todas las formas de ahorro para la jubilación subsidiado fiscalmente. Así, estos tipos de incentivos fiscales pueden aumentar de manera drástica la tasa de rendimiento después de impuestos del ahorro para la jubilación.



\textbf{APLICACIÓN PRESENTE EN EL MANUAL}

Efectos teóricos del ahorro para la jubilación subsidiado fiscalmente

Desde el punto de vista teórico, los subsidios fiscales al ahorro para la jubilación actúan de manera opuesta a como actuaba el impuesto sobre los ingresos por intereses en la Figura 22-1. Los efectos se ilustran en la Figura 22-3, que vuelve a mostrar la disyuntiva del consumo intertemporal.

Partimos de la restricción presupuestaria después de impuestos, BC_2, de esa figura, con una pendiente de -(1+r\times[1-t]); recuérdese que la pendiente es el tipo de interés después de impuestos, el precio del consumo del primer período (vida laboral), C_W.

Sin embargo, cuando el ahorro para la jubilación adopta una forma subsidiada fiscalmente, el ahorro se grava mucho más ligeramente porque los impuestos no se pagan hasta la jubilación (reduciendo el valor presente descontado, VPD, de los pagos de impuestos). Este retraso en el pago de impuestos reduce la carga fiscal sobre el ahorro de t a t\times r, donde r es la proporción de la carga tributaria que permanece después de tener en cuenta el retraso en el pago de impuestos. Supóngase, por ejemplo, que t=0,3 y que r es igual a 0,33; esto implicaría que los individuos que ahorran mediante planes de ahorro para la jubilación subsidiados fiscalmente tienen un tipo impositivo efectivo sobre el ahorro de 0,3\times0,33=0,1 una vez considerados los subsidios.

En este caso, la disponibilidad de ahorro subsidiado fiscalmente eleva la pendiente de la restricción presupuestaria a -(1+r\times[1-t\times r]), de modo que la restricción presupuestaria se desplaza de BC_2 a BC_3, elevando el precio del consumo del primer período. (Con un tipo de interés más alto, existe un mayor coste de oportunidad de consumir en lugar de ahorrar). Esto conduce a un efecto sustitución hacia un mayor ahorro y a un efecto renta hacia un menor ahorro. Los individuos ahorrarán más porque el precio del consumo presente (el tipo de interés después de impuestos) ha aumentado (el efecto sustitución); pero ahorrarán menos porque ahora pueden alcanzar sus objetivos de ahorro para la jubilación con mayor facilidad (el efecto renta).

Como se analizó anteriormente, el efecto de este cambio en el tipo de interés después de impuestos es ambiguo. La Figura 22-3 muestra dos posibles resultados de los subsidios al ahorro. Si el efecto sustitución es mayor que el efecto renta, entonces dichos subsidios aumentarán el ahorro (de S_2 a S_3), con una disminución del consumo del primer período de C_{W2} a C_{W3} (desplazamiento del punto A al punto B). Si el efecto renta es mayor que el efecto sustitución, entonces dichos subsidios reducirán el ahorro (de S_2 a S_4), con un aumento del consumo del primer período de C_{W2} a C_{W4} (desplazamiento del punto A al punto C).

Limitaciones del ahorro para la jubilación subsidiado fiscalmente

La mayoría de los instrumentos de ahorro para la jubilación subsidiado fiscalmente que revisamos anteriormente presentan un límite a las aportaciones anuales, como el límite de 5.500€ para las IRA. Estos límites complican el análisis teórico.

Supóngase que Andrea está tomando decisiones de ahorro en un mundo con y sin una IRA disponible. La Figura 22-4 muestra el impacto de la IRA sobre su restricción presupuestaria. Su restricción presupuestaria original después de impuestos, BC_2, entre el consumo del período uno y el consumo del período dos es la línea AB, con una pendiente de -(1+r\times[1-t]), donde r es el tipo de interés obtenido sobre sus ahorros y t es el tipo impositivo pagado sobre los ingresos por intereses. Sin la IRA, el precio del consumo del primer período, en términos del consumo del segundo período al que se renuncia, es 1+r\times(1-t), porque esa es la cantidad de consumo del segundo período que podría tener si consumiera un dólar menos hoy.

Si Andrea decide depositar dinero en una IRA, la tasa de rendimiento después de impuestos del ahorro cambia. La pendiente de la nueva restricción presupuestaria, BC_3, para los primeros 5.500€ de ahorro (que corresponden al consumo del primer período de C_{W2}) aumenta hasta -(1+r\times[1-t\times r]), donde t\times r vuelve a ser el tipo impositivo efectivo sobre el ahorro subsidiado fiscalmente. Este cambio se refleja en una nueva restricción presupuestaria, más inclinada, entre los puntos E y B. Por encima de los 5.500€ de ahorro (con un consumo del primer período inferior a C_{W2}), no existe ningún cambio en la tasa de rendimiento del ahorro debido al límite de las aportaciones a la IRA. Por encima de los 5.500€ de ahorro, por tanto, la pendiente de la restricción presupuestaria vuelve a su valor original.

¿Qué efecto tiene la IRA sobre el ahorro?

Considérense dos tipos de individuos, mostrados en la Figura 22-5.

En el panel (a), el Sr. Saltamontes ahorra poco antes de que se introduzca la IRA (punto A), consumiendo C_{W1} y ahorrando únicamente S_1=1.000€. Para el Sr. Saltamontes, el efecto de la IRA sobre el ahorro es ambiguo debido a los efectos sustitución (que tienden a aumentar el ahorro cuando aumenta el tipo de interés después de impuestos) y a los efectos renta (que tienden a reducir el ahorro cuando aumenta la renta después de impuestos), que actúan en direcciones opuestas. Si predominan los efectos sustitución, el Sr. Saltamontes se desplazará del punto A al punto B, reducirá el consumo del primer período a C_{W2} y aumentará el ahorro hasta S_2=1.500€. Si predominan los efectos renta, el Sr. Saltamontes se desplazará del punto A al punto C, aumentará el consumo del primer período hasta C_{W3} y reducirá el ahorro hasta S_3=500€.

El panel (b) muestra el análisis para la Sra. Hormiga, que era una gran ahorradora antes de que se introdujera la IRA, consumiendo C_{W1} y ahorrando 6.500€ (punto A). Para la Sra. Hormiga, la introducción de la IRA no modifica el precio del consumo del primer período, que sigue siendo 1+r\times(1-t), de modo que las restricciones presupuestarias con (BC2) y sin (BC1) IRA tienen la misma pendiente para ella.

Sin embargo, la introducción de las IRA sí tiene un efecto renta, porque ahora ella es más rica al disponer de un subsidio fiscal para sus primeros 5.500€ de ahorro: con el mayor rendimiento, ya no necesita ahorrar tanto (6.500€) para alcanzar el mismo nivel de consumo en el segundo período. Por tanto, utilizará su mayor renta para comprar más consumo en el primer período: el consumo del primer período aumenta de C_{W1} a C_{W2} cuando se desplaza del punto A al punto B. Su ahorro disminuye de S_1 a S_2=5.500€.

El efecto renta para los grandes ahorradores como la Sra. Hormiga surge porque estos «reasignan» sus activos existentes a las IRA; toman 5.500€ de los ahorros que ya estaban realizando y los vuelven a etiquetar como ahorro IRA con tratamiento fiscal preferente. Para los grandes ahorradores, por tanto, la IRA constituye un subsidio al ahorro ya existente: el gobierno les proporciona una mayor tasa de rendimiento sobre unos ahorros que ya tenían previsto realizar.

Si es probable que se produzca esta reasignación, entonces es posible que las IRA puedan reducir realmente el ahorro privado total a través de este efecto renta.

Este análisis utilizó las IRA como ejemplo, pero se aplica también a otras formas de ahorro para la jubilación. La diferencia con estas otras formas es que los límites son mucho más elevados, por lo que es probable que existan más individuos como el Sr. Saltamontes que como la Sra. Hormiga; por ejemplo, relativamente pocos trabajadores por cuenta propia habrían ahorrado más del límite de 53.000€ de una cuenta SEP-IRA para su jubilación sin disponer de una cuenta SEP-IRA.

Implicaciones de los modelos alternativos

Como se ha señalado, existe una sólida evidencia que sugiere que las decisiones de ahorro están determinadas tanto por consideraciones precautorias como de autocontrol. Ambas alternativas sugieren que los incentivos fiscales a la jubilación pueden tener efectos positivos sobre el ahorro más fuertes de lo que implica la teoría tradicional.

Ahorro precautorio

Esta reasignación de activos que acabamos de analizar será probable si se cumplen dos condiciones: (1) si una gran proporción de los contribuyentes a IRA habría ahorrado de todos modos 5.000€ en ausencia de la IRA y (2) si los ahorros en IRA y los ahorros fuera de IRA se consideran altamente sustitutivos. Considérese a alguien que iba a ahorrar 5.000€ para la jubilación incluso sin una IRA. Tal individuo no se sentiría limitado por el hecho de que no pudiera retirar el dinero antes de la jubilación, de modo que los ahorros en IRA y fuera de IRA serían sustitutos entre sí, y la reasignación sería probable (de modo que no habría ahorro nuevo neto).

Por otro lado, considérese a personas que tenían más de 5.000€ en ahorros pero estaban utilizando sus ahorros como precaución frente a la pérdida del empleo. Debido a que no se puede acceder a ellos si pierden su empleo antes de los 59 años y medio, los ahorros en una IRA no serían considerados muy sustitutivos para ellos, de modo que no reasignarían simplemente 5.000€ de sus ahorros (precautorios) a la IRA. El análisis de las Figuras 22-3, 22-4 y 22-5 se refiere únicamente al ahorro para la jubilación: quienes tienen ahorros precautorios superiores a 5.000€ no reasignan, pero quienes tienen ahorros para la jubilación superiores a 5.000€ sí lo hacen. El hecho de que muchos contribuyentes a cuentas de ahorro para la jubilación subsidiadas fiscalmente tengan elevados ahorros no implica necesariamente que vayan a reasignar esos ahorros a la cuenta subsidiada fiscalmente, si los ahorros son por razones precautorias en lugar de para la jubilación. Así, puede haber más ahorro nuevo debido a los incentivos a la jubilación de lo que sugiere el modelo tradicional analizado anteriormente.

Modelos de autocontrol

El rasgo distintivo de los modelos de autocontrol del ahorro es la búsqueda de mecanismos de compromiso que proporcionen autocontrol. Las cuentas de jubilación, como las cuentas de pensiones y 401(k), proporcionan excelentes mecanismos de compromiso porque las aportaciones se retiran directamente de la nómina y los individuos no pueden acceder a su dinero hasta la jubilación. Así, no se puede acceder al dinero de estas cuentas para satisfacer la impaciencia a corto plazo, permitiendo al individuo comprometerse con la paciencia a largo plazo. Las cuentas de ahorro para la jubilación pueden, por tanto, aumentar el ahorro más de lo que sugiere el análisis anterior. Además de la demanda de estas cuentas debido a los incentivos fiscales, también habrá una demanda que surge porque los individuos pueden comprometerse efectivamente a ahorrar.


\subsubsection{¿Oferta de crédito? Imposición y asunción de riesgo}
Hemos presentado las políticas gubernamentales que afectan a las decisiones individuales sobre cuánto ahorrar alterando directamente la tasa de rendimiento del ahorro. Centrémonos ahora en otros dos aspectos de la tributación que podrían afectar a las decisiones de ahorro de los contribuyentes:
\begin{la}
    \item \textbf{Asunción de riesgos}. El primero es la tributación de la asunción de riesgos. Los individuos no solo deciden cuánto ahorrar, sino también en qué forma mantener sus ahorros. ¿Debería ahorrar, por ejemplo, en forma de bonos del Estado, que son muy seguros pero ofrecen una rentabilidad media baja, o en forma de acciones de empresas, que son mucho más arriesgadas pero ofrecen una rentabilidad media más elevada? Del mismo modo que los impuestos pueden influir en cuánto ahorran los individuos, también pueden influir en la forma que adopta ese ahorro. Una política tributaria concreta que puede afectar a la asunción de riesgos es la tributación de las ganancias de capital, es decir, de los beneficios obtenidos por la venta de activos de capital; y,
    \item \textbf{Riqueza acumulada}. Los individuos también pueden ser gravados no solo sobre el rendimiento de sus ahorros en cada período, sino también sobre la cantidad de riqueza que han acumulado mediante ahorros pasados. En Estados Unidos, dicha tributación adopta dos formas. La primera forma de impuesto sobre la riqueza consiste en los impuestos sobre las transmisiones (transfer taxes), un conjunto de impuestos sobre las transferencias de bienes de una persona a otra, incluido el controvertido impuesto sobre sucesiones que acabamos de analizar. El segundo tipo de impuesto sobre la riqueza es el impuesto sobre los bienes inmuebles (property tax), que constituye la mayor fuente de ingresos de las administraciones locales en muchos países. Dado que la principal forma de ahorro para la mayoría de los estadounidenses es la vivienda, la tributación de la propiedad también puede afectar de manera significativa al nivel de ahorro.
\end{la}

Como vemos, el sistema tributario puede influir en la asunción de riesgos desde una doble óptica: (i) puede afectar a las decisiones de cartera de los hogares y, por tanto, a la disponibilidad de fondos; (ii) puede afectar a las decisiones reales de inversión de las empresas y los particulares.

La imposición puede afectar tanto a las tasas de rendimiento como a los niveles de riesgo de los diferentes activos. Siguiendo el trabajo de Domar y Musgrave (1944), es posible que los impuestos fomenten la asunción de riesgos, ya que el gobierno comparte el riesgo. Posteriormente. Arrow (1965) concluyó, a partir de series temporales sobre la demanda de dinero en Estados Unidos, que la elasticidad renta de la demanda del activo de riesgo es positiva pero inferior a la unidad. Esto implica en un modelo de dos activos tanto la imposición sobre la riqueza como la imposición sobre la renta aumentan la asunción social de riesgos. 

Así. el papel del gobierno en el reparto del riesgo resulta de especial interés. Si, por ejemplo, el mercado privado ofrece una distribución completa del riesgo, salvo para los riesgos sociales (como los riesgos asociados al ciclo económico) entonces el argumento de que el gobierno puede aumentar el riesgo compartido a través del sistema fiscal es menos convincente. En un caso extremo, suponemos que la recaudación impositiva se redistribuye a los individuos en forma de pagos unifo1mes o de suma fija y que inicialmente existe un activo de riesgo común que es comprado por todos los individuos, es decir, todos los riesgos están perfectamente correlacionados. Adicionalmente. si todos los individuos son idénticos y, por tanto, tienen una misma renta inicial, el impuesto se anula con la transferencia a tanto alzado, es decir, el impuesto no tiene ningún efecto sobre la asunción de riesgos. En todo caso. esta conclusión se justifica en el supuesto de que los resultados de la inversión de riesgo están perfectainente correlacionados para todos los individuos. Por lo tanto, no tendría sentido que el gobierno, a través de los impuestos, mejore la eficiencia con la que la economía gestiona el riesgo. Por el contrario. sí asl��mimos mercados de capitales imperfectos, el retorno de los diferentes activos estaría imperfectamente co1Telacionado por lo que el gobierno podría me:.:jorar el reparto del riesgo de forma más eficaz que el propio mercado.

Adicionalmente, los acuerdos de riesgo compartido disponibles en la economía pueden verse afectados por el sistema fi scal. Consideramos un ej emplo sencillo de la determinación del equilibrio general de los contratos de riesgo compm1ido formado por capitalistas (neutrales al riesgo) y gestores (aversos al riesgo). Los capitalistas proporcionan capital a los gestores, que lo emplean en una actividad productiva de riesgo y distribuyen parte del rendimiento a los capitalistas. Los capitalistas no pueden supervisar directamente las acciones de los gestores, por lo que recurren a incentivos para asegurarse de que éstos adopten las medidas que maximicen la rentabilidad de las inversiones. Si asumimos que el retorno esperado se determina en función del esfuerzo de los gestores multiplicado por el riesgo. el contrato óptimo para los capitalistas depende del grado de aversión al riesgo y de la elasticidad de la oferta de trabajo de los gestores. Por consiguiente, la introducción de un impuesto proporcional sobre la renta con compensación íntegra de pérdidas disminuiría el retorno obtenido por los capitalistas. Adicionalmente, bajo el supuesto de que los gestores presentan una aversión relativa al riesgo constante, éstos reducen su oferta laboral y, por tanto, el retorno recibido por los capitalistas disminuye en mayor proporción que el impuesto. En resumen, el impuesto tenderá a desincentivar la asunción de riesgos y



\paragraph{Fiscalidad y asunción de riesgos}
Nos enfrentamos al riesgo y a la incertidumbre en muchas dimensiones de la vida. Óscargunos son riesgos de resultados puramente adversos (accidentes de tráfico, enfermedades graves o lesiones), y tratamos de asegurarnos contra ellos. Otros riesgos están más equilibrados, con cierta probabilidad de un resultado positivo y cierta probabilidad de un resultado negativo. Un ejemplo de un riesgo más equilibrado es el riesgo asociado a las inversiones financieras: ¿Tendrá éxito o fracasará un negocio? ¿Subirá o bajará una cartera de acciones? La decisión de realizar o no inversiones arriesgadas puede verse afectada por la fiscalidad. 

El modelo básico de fiscalidad y asunción de riesgos fue desarrollado por DOMAR y MUSGRAVE en 1944 en el contexto del riesgo de las inversiones financieras. En su modelo, los individuos pueden elegir entre invertir en un activo seguro que no produce ningún rendimiento real e invertir en un activo arriesgado que produce una determinada tasa positiva de rendimiento. El gobierno grava cualquier rendimiento positivo del activo arriesgado, pero permite una deducción contra la renta imponible por el importe íntegro de cualquier rendimiento negativo. En esta situación, DOMAR y MUSGRAVE señalaron que gravar los rendimientos del activo arriesgado aumentaría la asunción de riesgos porque cualquier impuesto sobre esos rendimientos podría neutralizarse completamente asumiendo un mayor riesgo.

Este punto se ilustra mejor mediante un ejemplo. Supongamos que Óscar ha invertido 100€ en una pequeña empresa que tiene un 50\% de probabilidad de aumentar su valor hasta 120€ (de modo que obtiene una ganancia de 20€) y un 50\% de probabilidad de disminuir su valor hasta 80€ (de modo que pierde 20€). Esta inversión tiene un rendimiento esperado de 0€: el rendimiento de una inversión exitosa multiplicado por la probabilidad de éxito más el rendimiento de una inversión fallida multiplicado por la probabilidad de fracaso es igual a 0€ ([(20 $ × 0,5) + (−20 $ × 0,5) = 0 $]). Podemos resumir todas las combinaciones posibles con la siguiente tabla:

\begin{center}
    \begin{tabular}{p{0.09\linewidth}p{0.09\linewidth}p{0.09\linewidth}p{0.09\linewidth}p{0.09\linewidth}p{0.09\linewidth}p{0.09\linewidth}p{0.09\linewidth}p{0.09\linewidth}}
    \hline
      {\scriptsize\textbf{Política implementada}} & {\scriptsize\textbf{Inversión}} & {\scriptsize\textbf{Payoff (éxito)}} & {\scriptsize\textbf{Payoff (fracaso)}} & {\scriptsize\textbf{Impuesto ($t_{\text{éxito}}$)}} & {\scriptsize\textbf{Deducción $t_{\text{fracaso}}$}} & {\scriptsize\textbf{Recursos (after-tax)}} & {\scriptsize\textbf{Recursos (after-deduction)}} \\
    \hline
      {\scriptsize (1) Sin impuesto} & $100€$ & $20€$ & $-20€$ & $0\%$ & $0\%$ & $20€$ & $-20€$ \\
    \hline
      {\scriptsize (2) Con impuesto} & $100€$ & $20€$ & $-20€$ & $50\%$ & $50\%$ & $10€$ & $-10€$ \\
    \hline
      {\scriptsize (3) Compensación fiscal de pérdidas} & $200€$ & $40€$ & $-40€$ & $50\%$ & $50\%$ & $20€$ & $-20€$ \\
    \hline
      {\scriptsize (4) Sin compensación fiscal de pérdidas} & $200€$ & $40€$ & $-40€$ & $50\%$ & $0\%$ & $20€$ & $-40€$ \\
    \hline
      {\scriptsize (5) Progresividad fiscal} & $200€$ & $40€$ & $-40€$ & $75\%$ & $50\%$ & $15€$ & $-20€$ \\
    \end{tabular}
\end{center}

Inicialmente no existe ningún impuesto, por lo que estos rendimientos antes de impuestos son también sus rendimientos después de impuestos. Supongamos que el gobierno anuncia que va a introducir un impuesto del $50\%$ sobre la renta procedente de las inversiones, con una deducción contra la renta imponible por cualquier pérdida. Todo rendimiento positivo de esta inversión queda gravado, por lo que Óscar conserva la mitad de sus beneficios y paga la otra mitad en impuestos. Como cualquier pérdida derivada de la inversión se deduce de la renta imponible, Óscar solo soporta la mitad de la pérdida; las pérdidas se restan de la renta imponible, de modo que la reducción de impuestos compensa la mitad del importe de la pérdida.\footnote{%
    Como el gobierno recibe el $50\%$ de cada dólar de renta imponible, permitir una deducción contra la renta imponible significa que el gobierno también soporta el $50\%$ de cada dólar de pérdida.
} Bajo esta política (2), Óscar obtendrá un beneficio neto de solo $10€$ si la inversión sale bien, porque tiene que pagar $10€$ adicionales en impuestos; perderá solo $10€$ si la inversión sale mal, porque puede deducir la pérdida de $20€$ de la renta imponible, de modo que sus impuestos se reducen en $10€$. Este nuevo resultado tiene el mismo rendimiento esperado, $0€$€, pero deja a Óscar con menos riesgo del que desea; como revela su decisión inicial de inversión, él quería enfrentarse al riesgo de ganar $20€$ o perder $20€$, no al riesgo de ganar $10€$ o perder $10€$ (riesgo que podría haber tenido inicialmente invirtiendo solo $50€$, pero que decidió no asumir).

Sin embargo, Óscar puede neutralizar completamente los efectos de esta política fiscal sobre su cartera invirtiendo más dinero en la inversión arriesgada. Supongamos que duplica su inversión arriesgada hasta $200€$, (3). Después de impuestos, vuelve a existir un $50\%$ de probabilidad de ganar $20€$ si las cosas salen bien (porque Óscar gana $40€$ antes de impuestos y luego paga la mitad en impuestos) y un $50\%$ de probabilidad de perder $20€$ si las cosas salen mal (porque Óscar pierde $40€$ antes de impuestos y deduce la mitad de la pérdida de sus impuestos). Óscar invertir más en este activo arriesgado, Óscar ha neutralizado completamente el sistema fiscal del gobierno y ha regresado a las ganancias y pérdidas después de impuestos (1). Este resultado nos muestra que, si un agente económico puede deshacer las intervenciones del gobierno para volver a su equilibrio original, lo hará.

El gobierno es, en esencia, un socio silencioso en esta inversión. Como el gobierno asume parte del riesgo de éxito o fracaso, un individuo querrá aumentar su inversión arriesgada. En este caso, la tributación de las inversiones arriesgadas incrementa efectivamente la asunción de riesgos.

\textcolor{red}{Complicaciones del mundo real

Bajo el modelo simplificado de DOMAR-MUSGRAVE, la tributación de la renta procedente de inversiones aumentaría efectivamente la asunción de riesgos. En la realidad, esto puede no ser cierto debido a dos complicaciones importantes que el modelo no tiene en cuenta.

Compensación fiscal inferior al importe total de la pérdida}

En el modelo, los individuos pueden deducir íntegramente las pérdidas de su renta imponible al calcular los impuestos. Cuando Óscar pierde su apuesta, puede deducir la totalidad de las pérdidas de su renta imponible, del mismo modo que, cuando gana la apuesta, la totalidad de sus ganancias se añade a su renta imponible. En la realidad, la mayoría de los sistemas fiscales solo permiten a los individuos deducir una parte de sus pérdidas de la renta imponible al calcular los impuestos. La cantidad de una pérdida que puede deducirse se denomina compensación fiscal de pérdidas (tax loss offset).\footnote{%
    Por ejemplo, en Estados Unidos, los individuos solo pueden deducir hasta $3.000€$ de pérdidas por inversiones en un año fiscal de su otra renta imponible. Más concretamente, los individuos pueden compensar todas las pérdidas de inversión de un año con las ganancias de inversión obtenidas durante ese mismo año; cualquier pérdida neta restante puede deducirse de otras rentas imponibles, hasta un máximo de $3.000€$. Estas normas existen para impedir que las personas generen pérdidas fiscales realizando inversiones claramente perdedoras con el fin de eliminar parte de su renta imponible. Sin embargo, estas normas tienen además la implicación de que un contribuyente no puede simplemente neutralizar la tributación del gobierno realizando inversiones cada vez más arriesgadas, porque las pérdidas derivadas de esas inversiones no serán plenamente deducibles a efectos fiscales.
}

Continuando con el ejemplo, supongamos que el gobierno no permitiera a los individuos deducir ninguna pérdida por inversiones, como en el escenario (4). En ese caso, Óscar no podría simplemente neutralizar la política fiscal del gobierno invirtiendo más, como hizo en (3). Si Óscar eleva su inversión hasta $200€$ en la pequeña empresa, existe un $50\%$ de probabilidad de que gane $40€$, que se convierten en $20€$ después de impuestos, y un $50\%$ de probabilidad de que pierda $40€$, pérdida que no puede deducirse de la renta imponible y que, por tanto, sigue siendo una pérdida después de impuestos.\footnote{%
    Además, como se puede comprobar, sin posibilidad de deducir las pérdidas Óscar no puede simplemente compensar la Política Fiscal aumentando su inversión, porque eso lo situaría en una posición desfavorable: un $50\%$ de probabilidad de ganar $20€$ después de impuestos y un $50\%$ de probabilidad de perder $40€$ después de impuestos, con un rendimiento esperado después de impuestos de $(20\cdot  0,5) + (−40\cdot 0,5) = −10€$.
} Por tanto, como Óscar no puede simplemente deshacer la política del gobierno asumiendo más riesgo, no necesariamente aumentará su asunción de riesgos en respuesta a la tributación. 

Es imposible predecir con certeza qué efecto tendrán las limitaciones a la compensación de pérdidas sobre la asunción de riesgos, pero una compensación fiscal inferior al importe total de la pérdida limita claramente la aplicabilidad del sencillo modelo de DOMAR-MUSGRAVE.

Tributación redistributiva

El modelo idealizado de la asunción de riesgos también supone un tipo impositivo constante sobre la renta procedente de las inversiones. En la realidad, los sistemas fiscales suelen ser progresivos, con tipos marginales más elevados a medida que aumenta la renta. Bajo un sistema de este tipo, si los inversores obtienen grandes ganancias en una apuesta arriesgada, pueden pasar a un tramo impositivo superior (un tipo marginal más alto); si sufren grandes pérdidas, pueden pasar a un tramo impositivo inferior. Así, las ganancias de las apuestas exitosas pueden gravarse a un tipo superior al tipo al que se deducen las pérdidas de las apuestas fracasadas. Óscar igual que ocurre con la compensación limitada de pérdidas, esto puede llevar a los inversores a reducir la asunción de riesgos.

Supongamos que, en lugar de un tipo fijo del $50\%$, el gobierno establece un tipo impositivo del $50\%$ sobre la renta hasta $20€$ y del $75\%$ sobre la renta que exceda de $20€$. Supongamos también que existe una compensación total de pérdidas, de modo que cualquier pérdida puede deducirse de la renta imponible al tipo del $50\%$. En este caso, Óscar vuelve a no poder neutralizar los efectos del impuesto sobre sus ganancias, (5). Si aumenta su inversión hasta $200€$ y el resultado de la inversión es favorable, solo se lleva a casa $15€$ después de impuestos: el $50\%$ de los primeros $20€$ de ganancias y el $25\%$ de los siguientes $20€$. Si, por el contrario, el resultado de la inversión es desfavorable, pierde $20€$ después de impuestos, tal como se muestra. Una vez más, Óscar no compensará simplemente la Política Fiscal aumentando su inversión, porque eso lo situaría en una posición desfavorable, con un $50\%$ de probabilidad de ganar $15€$ después de impuestos y un $50\%$ de probabilidad de perder $20€$ después de impuestos (un rendimiento esperado de $[(15\cdot 0,5) + (−20 \cdot 0,5)] = −2,50€)$.



Evidencia sobre la tributación y la asunción de riesgos

Debido a estas complicaciones, así como a otras, no existe una predicción clara sobre cómo afectará la tributación a la asunción de riesgos en el mundo real. En última instancia, el efecto de los impuestos sobre la asunción de riesgos es una cuestión empírica. Desafortunadamente, sin embargo, existe muy poca evidencia sobre el efecto de la tributación de las rentas del capital sobre la asunción de riesgos, lo que convierte a este tema en otro de los importantes misterios de la economía de la tributación.

\begin{specialblock}{\textbf{Otras aplicaciones}: Asunción de riesgos relativos a la invesión en trabajo}
    Las inversiones financieras no son las únicas inversiones arriesgadas que realizan los individuos; también pueden invertir en capital humano mediante la educación u otra formación laboral. Invertir en capital humano es arriesgado porque los individuos realizan un sacrificio inicial a cambio de una expectativa de mayores ingresos en el futuro. Sin embargo, el grado en que esta inversión en capital humano conducirá a mayores ingresos es incierto.

    Considérese la decisión de asistir un año a la universidad. La asistencia tiene dos costes: el coste del año de estudios (matrícula, libros, etc.) y los ingresos dejados de percibir que el individuo renuncia al asistir a la universidad durante un año en lugar de trabajar. El beneficio es la perspectiva de obtener mayores ingresos como consecuencia de una mayor educación.\footnote{%
        Por ejemplo, en Estados Unidos existe un aumento aproximado del $7\%$ en los ingresos por cada año adicional de educación, en promedio. Sin embargo, esta es solo una estimación media y, para cualquier individuo concreto, el rendimiento de una mayor educación podría ser inferior o superior.
    } ¿Cómo afectan los impuestos sobre la renta a la decisión de acumular capital humano? 
    
    El análisis es similar al de las inversiones financieras. El rendimiento neto de la inversión en capital humano es el aumento de los salarios menos los costes directos de la educación y el coste indirecto de oportunidad derivado de los ingresos dejados de percibir. Supongamos inicialmente que existe un impuesto sobre la renta de tipo único y fijo y que los costes financieros de adquirir capital humano (como la matrícula) son plenamente deducibles de la renta imponible. Como los impuestos sobre la renta solo se aplican a los ingresos obtenidos, el coste de oportunidad de la educación también es plenamente deducible porque obtener educación significa dejar de percibir salarios, reduciendo así los impuestos. En este caso, al igual que ocurre con la inversión financiera, una mayor tributación simplemente incrementaría la inversión en capital humano para que los individuos pudieran preservar la cantidad deseada de rendimiento esperado neto (después de impuestos) de la inversión.
    
    Existe evidencia de que los individuos reasignan su cartera hacia activos con tratamiento fiscal preferente cuando aumentan los tipos impositivos; más recientemente, un estudio de ALAN et al. (2010) encuentra efectos bastante pequeños de los impuestos sobre la asignación de la cartera.
    
    Una vez más, sin embargo, las características de la tributación en el mundo real complican el panorama. Además, el sistema tributario progresivo vigente en la mayoría de los países implica que las ganancias derivadas de un aumento exitoso del capital humano serán gravadas a un tipo superior al que se aplicaría a los salarios obtenidos si se optara por trabajar en lugar de estudiar. Si no vas a la universidad y, en cambio, trabajas en el sector de la comida rápida, renuncias a la posibilidad de obtener un salario elevado, pero al menos pagas pocos impuestos; si renuncias a los ingresos del trabajo en la comida rápida para ir a la universidad, los rendimientos de esa educación tributan a un tipo elevado. Esto hace que las inversiones en capital humano sean menos atractivas. Así pues, al igual que ocurre con las inversiones financieras, el impacto neto de la tributación sobre la renta en las inversiones en capital humano es ambiguo.
\end{specialblock}








\textcolor{red}{Tributación de las ganancias de capital

Los análisis de los dos capítulos anteriores se han centrado en la tributación de la renta del trabajo, es decir, de la renta generada ya sea por los ingresos laborales o por los ingresos procedentes de intereses. Sin embargo, muchos activos proporcionan a los inversores un rendimiento que no adopta la forma de ingresos anuales por intereses, sino de una ganancia de capital, es decir, la diferencia entre el precio de compra de un activo y su precio de venta. Esta es la forma de rendimiento correspondiente a las inversiones en obras de arte y viviendas, así como la principal fuente de rendimiento de muchas inversiones en empresas o acciones. 

La cuestión de cómo gravar estos rendimientos ha sido uno de los temas más controvertidos de la política fiscal durante las últimas décadas.

Tratamiento fiscal actual de las ganancias de capital

Los activos (como las cuentas bancarias o los bonos del Estado) que generan intereses se gravan sobre la base del devengo; esto es, los impuestos se pagan en cada período sobre los intereses obtenidos en ese período. Las ganancias de capital, por el contrario, se gravan en el momento de su realización; esto es, los impuestos solo se pagan cuando el activo se vende, y el importe del impuesto se basa en la diferencia entre el precio de venta del activo y su precio de compra. La tributación en el momento de la realización, en lugar de en el momento del devengo, conduce por lo general a una reducción de las obligaciones tributarias del titular del activo. La intuición que subyace a esta afirmación es la misma que utilizamos en nuestra discusión sobre las cuentas de ahorro para la jubilación con ventajas fiscales. Óscar pagar los impuestos cuando se vende el activo, en lugar de hacerlo a medida que su valor se va devengando, se pueden obtener los intereses correspondientes a lo que, de otro modo, habría constituido la obligación tributaria.

Supongamos, por ejemplo, que compras un cuadro por $100€$ y que este aumenta de valor un $10\%$ anual, de modo que, al cabo de siete años, lo vendes por $195€$. En ese momento pagas impuestos sobre la diferencia de $95€$ entre el precio de venta y el precio de compra. Si el tipo impositivo sobre las ganancias de capital es del $20\%$, pagarás $19€$ en concepto de impuesto sobre las ganancias de capital y obtendrás una ganancia neta de $76€$. Si, en lugar de ello, hubieras invertido esos $100€$ en una cuenta bancaria que rindiera un $10\%$ anual, con un tipo impositivo del $20\%$ sobre los ingresos por intereses, al final de los siete años solo tendrías $71€$ en ingresos por intereses. Se obtiene un rendimiento menor cuando se pagan impuestos sobre el devengo porque el gobierno recauda sus ingresos antes y, por tanto, es él quien obtiene los intereses correspondientes a esos ingresos, en lugar de que los obtengas tú. 

Esto constituye un tratamiento fiscal preferente para el ahorro en forma de activos que generan ganancias de capital.

La preferencia fiscal por las ganancias de capital es difícil de eliminar para muchos bienes de capital por dos razones. En primer lugar, puede que no sea posible medir el devengo para muchos activos. En el caso de las acciones, es factible que el gobierno pudiera gravar el devengo anual utilizando la variación del valor de las acciones de un año a otro. Esto podría funcionar para las acciones, pero ¿cómo estimamos el devengo anual del valor de una vivienda o de un cuadro? En principio, podríamos obtener una valoración pericial de cada activo de capital al final de cada año, pero la ventaja del sistema basado en la realización es que se apoya en valoraciones de mercado, y no en la opinión de expertos, que es tanto imperfecta como susceptible de manipulación.

El segundo problema es que, incluso si el gobierno pudiera medir adecuadamente el devengo, es posible que los individuos no tuvieran capacidad para financiar el pago del impuesto exigido. Supongamos que una acción muy volátil acaba de duplicar su valor. El incremento resultante en la riqueza de una persona puede ser tan grande que no exista forma de pagar la factura fiscal anual sin vender realmente una parte importante de sus acciones. Podría ser ineficiente obligar a los individuos a desprenderse de activos productivos simplemente para poder pagar impuestos.

Mejora de la eficiencia de las transacciones de capital

Un segundo argumento importante a favor de tipos impositivos más bajos sobre las ganancias de capital es que los individuos retrasarán la venta de sus activos de capital para reducir el valor actual descontado de su carga tributaria. La ventaja de poseer un activo de capital que tributa en el momento de su realización, en comparación con otra inversión (como una cuenta bancaria) que tributa sobre el devengo, es que el pago de impuestos puede aplazarse hasta el momento de la venta. Cuanto más largo sea este aplazamiento, menor será el valor actual descontado de los pagos de impuestos (porque el dinero pagado en el futuro vale menos que el dinero pagado hoy). De hecho, una persona podría reducir a cero los impuestos sobre las ganancias de capital de su familia conservando un activo hasta su fallecimiento y transmitiéndolo a sus hijos. Este comportamiento puede dar lugar a un efecto de inmovilización (\textit{lock-in effect}), por el cual los individuos retrasan la venta de sus activos de capital para minimizar el valor actual descontado de los pagos del impuesto sobre las ganancias de capital.

Los efectos de inmovilización son costosos porque gran parte del éxito de los mercados de capitales descansa en su fluidez, que permite a los inversores destinar sus activos al uso más productivo posible. Supongamos que tienes una gran idea para un nuevo producto, pero debes vender tu colección de obras de arte para financiar el desarrollo de ese producto, lo que desencadenaría el pago del impuesto sobre las ganancias de capital derivadas de la venta de las obras. Si no vendes la colección de arte para financiar este nuevo producto, dicha colección pasará a tus hijos, evitando por completo la tributación de las ganancias de capital. Por tanto, podrías decidir no vender tus cuadros para minimizar el impuesto sobre las ganancias de capital y, como resultado, quizá nunca pongas en marcha esa nueva empresa. Esto puede representar un coste de eficiencia para la sociedad si la eficiencia social fuera mayor porque hubieras creado la empresa en lugar de conservar tus cuadros. Este constituye un argumento a favor de reducir los impuestos sobre las ganancias de capital para disminuir el efecto de inmovilización.}

Fomento de la actividad emprendedora

Cuando te gradúas en la universidad y creas tu propia empresa, el principal rendimiento que esperas obtener no son los ingresos que percibirás cada año, sino el aumento del valor de la propia empresa. Los emprendedores, individuos que crean sus propios negocios, obtienen la mayor parte de su riqueza no de ingresos devengados durante las primeras etapas de vida de la empresa, sino del aumento del valor del activo empresarial subyacente a lo largo del tiempo. El tipo impositivo relevante para los emprendedores no es, por tanto, el impuesto sobre la renta sino el impuesto sobre las ganancias de capital, ya que su principal rendimiento procederá de la venta de sus activos cuando estos adquieran valor. En consecuencia, un tipo impositivo más elevado sobre las ganancias de capital puede desalentar el emprendimiento. El impuesto sobre las ganancias de capital es, en esencia, un impuesto sobre la asunción de riesgos, pero la asunción de riesgos constituye un motor del crecimiento económico. Por ello, fomentar el emprendimiento es una de las razones para mantener un tipo impositivo más bajo sobre las ganancias de capital.

Existen, sin embargo, tres argumentos que contrarrestan esta afirmación. En primer lugar, como se analizó anteriormente, no está claro si gravar la asunción de riesgos fomenta o desincentiva dicha asunción de riesgos. Hasta la fecha no existe evidencia concluyente acerca de si los impuestos sobre las ganancias de capital aumentan o reducen la asunción de riesgos. En segundo lugar, solo una fracción muy pequeña de las ganancias de capital corresponde a emprendedores.\footnote{%
    Por ejemplo, en Estados Unidos, según una estimación únicamente el $10\%$ de las inversiones de capital riesgo se veían afectadas por los impuestos sobre las ganancias de capital, porque la mayor parte de la financiación de capital riesgo no procede de individuos, sino de otras entidades, como los fondos de pensiones.
} Por último, aunque las reducciones actuales de los tipos impositivos sobre las ganancias de capital pueden incrementar la actividad emprendedora tanto en el presente como en el futuro, también proporcionan enormes beneficios a quienes realizaron inversiones de capital en el pasado. Un impuesto más bajo sobre las ganancias de capital no solo constituye un incentivo para que los inversores asuman riesgos empresariales; también representa una recompensa por haber asumido riesgos en el pasado. Como resultado, una gran parte de la recaudación perdida debido a la reducción de los impuestos sobre las ganancias de capital no está fomentando la asunción de riesgos, sino simplemente recompensando riesgos asumidos anteriormente.

Este punto añade un matiz a los análisis sobre los efectos marginales e inframarginales de los incentivos fiscales. En los casos considerados anteriormente, las cuentas IRA y las donaciones benéficas, el efecto marginal era el nuevo comportamiento incentivado (nuevo ahorro o nuevas donaciones benéficas), mientras que el efecto inframarginal era el comportamiento que habría tenido lugar de todos modos pero que ahora recibía un tratamiento fiscal preferente (ahorros o donaciones previamente previstos). Ambos efectos están presentes en el caso de los impuestos sobre las ganancias de capital: el gobierno está subvencionando inversiones que se habrían realizado de todas formas, al mismo tiempo que incentiva nuevas inversiones en el presente.

Sin embargo, en el caso de las ganancias de capital existe un efecto inframarginal adicional: el gobierno no solo subvenciona las inversiones que de todos modos se realizarían en la actualidad, sino que también subvenciona los rendimientos de inversiones efectuadas en el pasado. Este constituye un coste adicional de recaudación que debe ponderarse frente a las ganancias marginales derivadas del fomento de la asunción de riesgos empresariales.

En otras palabras, reducir el tipo impositivo sobre las ganancias de capital es un instrumento muy poco preciso para fomentar el emprendimiento. Puede aumentar los rendimientos del emprendimiento en el futuro, pero al elevado coste recaudatorio de recompensar también las inversiones realizadas en el pasado. Para verlo con claridad, compárese una reducción general del impuesto sobre las ganancias de capital con una reducción prospectiva del impuesto sobre las ganancias de capital, en la que el tipo impositivo se reduce únicamente para las ventas de inversiones realizadas a partir de ese momento, y no para las ventas de inversiones efectuadas en el pasado. Tanto una reducción general del impuesto sobre las ganancias de capital como una reducción prospectiva tendrían efectos idénticos sobre la actividad emprendedora futura, ya que todas las inversiones realizadas a partir de ese momento se beneficiarían del tipo impositivo más bajo sobre las ganancias de capital. Sin embargo, la reducción prospectiva sería mucho menos costosa porque no concedería también una ventaja fiscal a las inversiones realizadas en el pasado.

Evidencia sobre la tributación y las ganancias de capital

Para resumir este análisis, existen dos argumentos principales a favor de fijar un tipo impositivo sobre las ganancias de capital inferior al aplicado a otras formas de renta del capital: liberar las ganancias acumuladas en el pasado y fomentar el emprendimiento. Reducir el tipo aplicable a todas las ganancias de capital, incluidas las correspondientes a inversiones realizadas en el pasado, constituye una forma muy costosa de fomentar el emprendimiento, en comparación con reducir el tipo únicamente para las inversiones realizadas a partir del momento presente. El principal argumento a favor de una reducción general del impuesto sobre las ganancias de capital que incluya también las ganancias ya obtenidas anteriormente es, por tanto, que un tipo impositivo más bajo dará lugar a un mayor número de realizaciones de ganancias de capital (ventas de activos). Esto tiene dos ventajas: libera los activos para que se destinen a su uso más productivo y aumenta la recaudación tributaria sobre esas ganancias.

Si existe un importante efecto de liberación (\textit{unlocking effect}) —es decir, si un impuesto más bajo sobre las ganancias de capital anima a las personas a vender sus activos ahora, en lugar de esperar o de no venderlos en absoluto—, es posible que una reducción del tipo impositivo sobre las ganancias de capital llegue incluso a aumentar la recaudación: estaríamos aplicando un tipo impositivo más bajo sobre una base imponible mucho mayor de ganancias de capital. Es decir, el gobierno podría encontrarse en el lado equivocado de la curva de LAFFER (p. 641) con respecto a los impuestos sobre las ganancias de capital. Sin embargo, si las reducciones de este impuesto no afectan a la venta de activos, o si únicamente modifican el momento de la venta de un año al siguiente, entonces las reducciones del tipo impositivo sobre las ganancias de capital simplemente disminuirán la recaudación tributaria y proporcionarán una transferencia a aquellos titulares de ganancias de capital que, de todas formas, iban a realizarlas.

La siguiente figura muestra la evolución temporal de las realizaciones de ganancias de capital realizadas para un periodo muestral en Estados Unidos:

\imagenfit{Fig8.png}{Ganancias de capital: Realización de las ganancias acumuladas (1960-2010)}{\textit{Public Financie and Policy Making}. GRUBER (2017)}

Como vemos, queda bastante patente que los impuestos desempeñan un papel importante en la determinación de dichas realizaciones. Se produjo un enorme incremento de las realizaciones en 1986: estas ascendieron a $172.000\text{M}\$$ en 1985, $328.000\text{M}\$$ en 1986 y $148.000\text{M}\$$ en 1987. Este aumento fue impulsado por el importante incremento relativo de la tributación de las ganancias de capital cuya entrada en vigor estaba prevista para 1987, pero que fue anunciada en 1986; como se señaló anteriormente, en 1987 las ganancias de capital pasaron de tributar a un tipo mucho menor que otras formas de renta a tributar del mismo modo que las demás rentas. Los individuos anticiparon este incremento relativo del impuesto y vendieron sus activos antes de que el aumento impositivo entrara en vigor.

Sin embargo, el hecho de que las realizaciones de ganancias de capital volvieran a su nivel original en 1987 sugiere que la respuesta al aumento del impuesto sobre las ganancias de capital fue únicamente temporal. De hecho, Burman y Randolph (1994) estimaron que prácticamente toda la respuesta de las realizaciones de ganancias de capital a la tributación es transitoria, o de corto plazo, y que existe muy poco incremento a largo plazo en la tasa de realizaciones de ganancias de capital cuando los tipos impositivos son más bajos. Cuando se reducen los impuestos sobre las ganancias de capital, los individuos simplemente adelantan la venta de activos que de todos modos iban a vender. La recaudación puede aumentar en el corto plazo a medida que esos activos se venden, pero disminuye en el largo plazo porque la base restante de ganancias es menor y sobre esa base se aplica un tipo impositivo más bajo.

Burman y Randolph estiman que la elasticidad de largo plazo de las ganancias de capital con respecto al tipo impositivo es de tan solo $–0,18$; con una elasticidad de este valor, una reducción del tipo impositivo sobre las ganancias de capital del $1\%$ reduciría la recaudación aproximadamente en un $0,82\%$ (porque el aumento compensatorio de las realizaciones sería muy pequeño). Así pues, nos encontramos claramente en el lado correcto de la curva de LAFFER en lo que respecta a los impuestos sobre las ganancias de capital. Además, no parece existir una asignación ineficiente significativa de los activos a largo plazo que pueda corregirse reduciendo el tipo impositivo sobre las ganancias de capital; más bien, los cambios en dicho tipo parecen afectar principalmente al reajuste temporal, a corto plazo, del momento en que se venden los activos.

Esta interpretación se ve respaldada por el comportamiento observado a finales de 2012. Los tipos impositivos aplicables a los contribuyentes de mayores ingresos en Estados Unidos aumentaron a comienzos de 2013; así, a medida que el año 2012 llegaba a su fin, se produjo una enorme actividad orientada a realizar ingresos en 2012 en lugar de hacerlo a los tipos más elevados vigentes en 2013. A finales de 2012, el New York Times informó de que, para las personas con mayores patrimonios, \textit{el objetivo primordial es registrar este año la mayor parte posible de sus ingresos futuros}. Como afirmó Kenneth Bezozo, socio de un despacho de abogados de Nueva York especializado en este tipo de operaciones, a finales de 2012: En mis 30 años de ejercicio profesional, nunca había visto tal avalancha de deseo y de actuaciones para transferir una empresa y hacer efectiva su venta.

Además, parte de la recaudación obtenida mediante tipos impositivos más bajos sobre las ganancias de capital se consigue a costa de otras fuentes de ingresos tributarios. Por ejemplo, muchos trabajadores con remuneraciones muy elevadas pueden elegir, hasta cierto punto, si desean ser remunerados en forma de salarios en efectivo o mediante acciones que generen rendimientos sujetos al impuesto sobre las ganancias de capital. Cuando los impuestos sobre las ganancias de capital se reducen en relación con los impuestos sobre los salarios, los individuos optarán por esta segunda alternativa, percibiendo menos ingresos en forma de salarios (gravados a tipos elevados) y más remuneración en forma de acciones (gravadas a tipos reducidos). Este desplazamiento provocará un aumento de la recaudación procedente del impuesto sobre las ganancias de capital cuando esas acciones sean vendidas, pero ese incremento de la recaudación se producirá a costa de una menor recaudación del impuesto sobre una base salarial más reducida.

¿Cuáles son los argumentos en contra de las preferencias fiscales para las ganancias de capital?

Existen dos argumentos en contra del tratamiento preferente que reciben las rentas derivadas de las ganancias de capital en la mayoría de los países. En primer lugar, los impuestos sobre las ganancias de capital son muy progresivos: los ingresos por ganancias de capital corresponden principalmente a los contribuyentes más ricos. Los individuos con mayores ingresos acumulan de forma desproporcionada las ganancias de capital: en Estados Unidos, en 2010, los $400$ contribuyentes con mayores ingresos concentraban el $16\%$ de todas las ganancias de capital. Por tanto, cualquier reducción de estos impuestos beneficia principalmente a los contribuyentes con mayores niveles de renta.

En segundo lugar, unos tipos impositivos más bajos sobre las ganancias de capital vulneran el principio de HAIG-SIMONS para los sistemas tributarios. El objetivo de la tributación debería ser proporcionar un terreno de juego equilibrado entre las distintas decisiones económicas, y no favorecer una elección frente a otra, salvo que exista algún argumento de equidad o de eficiencia que lo justifique (como una externalidad positiva que justifique una preferencia fiscal). Mantener un tipo impositivo sobre las ganancias de capital inferior al aplicado a otras formas de renta introduce una cuña fiscal que proporciona incentivos para que los individuos adopten decisiones ineficientes con el fin de acceder al tipo impositivo más bajo.

Como se señaló anteriormente, por ejemplo, cuando los impuestos sobre las ganancias de capital se reducen en relación con los impuestos sobre los salarios, los individuos optan por recibir una menor parte de su remuneración en forma de salarios y una mayor parte en forma de acciones; cambio que puede no ser eficiente, en la medida en que fomenta una asunción excesiva de riesgos e incluso, posiblemente, comportamientos fraudulentos por parte de los directivos. Sin embargo, este desplazamiento se ve incentivado por la cuña fiscal existente entre la tributación de las ganancias de capital y la tributación de otras formas de renta, como los salarios.

Del mismo modo, gravar las ganancias de capital a un tipo inferior al aplicado a los intereses obtenidos puede llevar a los individuos a orientar sus inversiones hacia activos más arriesgados que generan ganancias de capital (como las acciones), en lugar de hacia activos que generan ingresos por intereses (como los bonos). La eficiencia de la inversión se maximiza cuando las decisiones de inversión se toman en función del valor intrínseco que dichas inversiones tienen para el individuo, y no en función de una cuña fiscal. Cuando los individuos eligen una cartera demasiado arriesgada debido a una preferencia fiscal por las ganancias de capital, la eficiencia global de la inversión disminuye.












Conclusión

El impacto del sistema tributario sobre las decisiones relativas a cuánto ahorrar y en qué forma hacerlo seguirá ocupando siempre un lugar central en los debates sobre la reforma fiscal. En este capítulo hemos revisado algunos de los principales aspectos del sistema tributario que afectan a la asunción de riesgos y a la acumulación de riqueza en Estados Unidos. En primer lugar, destacamos que la tributación no necesariamente reduce la asunción de riesgos y que, bajo determinados supuestos, la incrementa de forma inequívoca.

Este resultado teórico tiene importantes implicaciones para el debate sobre la tributación de las ganancias de capital. Los argumentos más sólidos a favor del tratamiento fiscal preferente de las ganancias de capital son que: (1) unos tipos impositivos más bajos sobre las ganancias de capital «liberarán» activos productivos (unlock productive assets), y (2) unos impuestos más bajos sobre las ganancias de capital fomentarán el emprendimiento.

La evidencia existente sobre el primero de estos argumentos sugiere que dicho efecto de liberación no es importante a largo plazo. La discusión teórica sugiere que las predicciones relativas al emprendimiento no son concluyentes. Además, incluso si una menor tributación de las ganancias de capital promoviera la asunción de riesgos y el emprendimiento, lo haría a un coste muy elevado, consistente en proporcionar importantes subvenciones a inversiones realizadas con anterioridad.

También analizamos el impuesto sobre sucesiones, que constituye un impuesto muy progresivo, pero que plantea difíciles cuestiones relacionadas con la equidad horizontal y con su aplicación efectiva. Por último, analizamos la incierta incidencia del impuesto sobre la propiedad y las importantes cuestiones relativas a la tributación de las empresas frente a las viviendas residenciales, así como de los terrenos frente a las mejoras realizadas sobre dichos terrenos.


\subsubsection{Oferta de bienes y servicios: ¿?}
Por último, antes de pasar a presentar los efectos de la imposición sobre la oferta desde una perspectiva macroeconómica, debemos presentar el marco analítico a través del cual estudiar los efectos de la imposición societaria.

¿Por qué existe un impuesto sobre la renta de las sociedades? Después de todo, las empresas no son entidades en sí mismas, sino combinaciones de factores de producción. Por consiguiente, cuando gravamos a las empresas, en última instancia estamos gravando los factores de producción que las integran. ¿No sería más sencillo gravar directamente esos factores (trabajo y capital), en lugar de obtener ingresos públicos a través del complejo (e incierto) mecanismo del impuesto sobre sociedades? Existen al menos dos razones por las que puede ser conveniente disponer de un impuesto específico,

Tributación de los beneficios puros. En primer lugar, en la medida en que las sociedades poseen poder de mercado, obtendrán beneficios puros, es decir, rendimientos que superan las remuneraciones pagadas a sus factores de producción. Como demostraron DIAMOND y MIRRLEES (1971) en un análisis fundamental, un impuesto sobre los beneficios puros constituye una forma mucho mejor de recaudar ingresos públicos que un impuesto sobre los factores de producción. Ello se debe a que \textbf{un impuesto sobre los beneficios puros no distorsiona la toma de decisiones del productor}, mientras que los impuestos sobre el trabajo o el capital generan efectos distorsionadores.

Una empresa elige los precios y los niveles de producción con el objetivo de maximizar sus beneficios. Si el gobierno anunciara mañana que se apropiará de una parte de esos beneficios, la elección óptima de precios y cantidades por parte de la empresa no cambiaría; la decisión que maximiza los beneficios antes de impuestos también maximiza los beneficios después de impuestos: $(1 − \text{impuesto}) \cdot \text{beneficios}$. Por tanto, los impuestos sobre los beneficios puros permiten recaudar ingresos sin distorsionar el comportamiento. Además de no generar distorsiones, un impuesto sobre los beneficios puros es muy progresivo, porque quienes perciben los beneficios derivados de la producción suelen pertenecer a los grupos con mayor nivel de renta.\footnote{%
    Aunque el impuesto sobre los beneficios puros parece una buena idea, no es así como funciona el impuesto sobre sociedades, por dos razones. En primer lugar, los impuestos sobre sociedades no son impuestos sobre beneficios puros: las empresas pueden minimizar su carga tributaria modificando la utilización que hacen de sus factores de producción. Como los impuestos sobre sociedades inducen a las empresas a apartarse de su patrón óptimo de producción mediante la sustitución de unos factores por otros, generan una distorsión que produce ineficiencia. En segundo lugar, un impuesto sobre los beneficios puros debería gravar los beneficios económicos, es decir, la diferencia entre los ingresos de la empresa y sus costes económicos. Ello implicaría utilizar un tipo de cálculo de costes que valorase adecuadamente los recursos según su coste de oportunidad (su utilización en la mejor actividad alternativa disponible). Sin embargo, las empresas pagan impuestos sobre sus beneficios contables, definidos como la diferencia entre los ingresos y los costes registrados contablemente. Los costes registrados pueden diferir de los costes económicos tanto porque utilizan precios en lugar de costes de oportunidad como porque las empresas pueden manipular sus prácticas contables para variar el importe que declaran como costes.
}

Beneficios retenidos. Otra justificación del impuesto sobre sociedades es similar a los argumentos analizados anteriormente respecto a las ganancias de capital: si las sociedades no tributaran por sus beneficios, los individuos propietarios de acciones podrían simplemente evitar los impuestos haciendo que las sociedades nunca distribuyeran sus beneficios. Estos beneficios se acumularían libres de impuestos dentro de la sociedad, dando lugar a una importante subvención fiscal a los beneficios empresariales en comparación con otras formas de ahorro (o, en términos más generales, con otras actividades económicas). Si las sociedades distribuyeran esos beneficios muchos años después, el valor actual descontado de la carga tributaria sería muy reducido.

\textcolor{red}{El 4 de mayo de 2009, Barack Obama propuso reformar el sistema del impuesto sobre sociedades poniendo fin a los incentivos fiscales para las empresas estadounidenses que habían invertido de forma significativa en operaciones en el extranjero. Afirmó que estas empresas estaban «eludiendo» sus responsabilidades y constituían un factor importante de las «inequidades de un sistema tributario roto que recompensaba a las empresas por crear empleos en Bangalore en lugar de Buffalo, Nueva York». En particular, Obama propuso reformar las normas de diferimiento (deferral rules), de modo que las empresas no pudieran deducir gastos en sus declaraciones fiscales estadounidenses hasta haber pagado los impuestos correspondientes sobre los beneficios obtenidos en el extranjero. Esta disposición, cuya entrada en vigor estaba prevista para 2011, recaudaría 60.000 millones de€ entre 2011 y 2019. El segundo componente de las propuestas de Obama recaudaría un total de 87.000 millones de€ entre 2011 y 2019 mediante restricciones más estrictas sobre los «paraísos fiscales» (tax havens) en el extranjero, donde los tipos impositivos son especialmente bajos.

Los defensores de la propuesta coincidían en que el gobierno debía intervenir para arreglar el sistema tributario. TODER, experto en política fiscal y antiguo funcionario del Departamento del Tesoro durante la presidencia Clinton, afirmó que los planes de Obama cerrarrían lagunas que erosionan la base imponible del impuesto sobre sociedades y que no tendrían ningún efecto sobre el empleo en Estados Unidos. Anna Burger, secretaria-tesorera del \textit{Service Employees International Union}, sostuvo que, por fin, se igualarían las condiciones para aquellas empresas que habían recibido un tratamiento fiscal injustificadamente favorable.

Otros argumentaron, sin embargo, que, en lugar de crear más oportunidades en Estados Unidos, el plan de Obama tendría un efecto contraproducente y reduciría la creación de empleo en el país. John Castellani, presidente de la Business Roundtable, afirmó que el sistema vigente era beneficioso para las empresas estadounidenses porque sus competidores internacionales recibían un tratamiento fiscal más favorable cuando desarrollaban actividades en el extranjero. Drew Lyon, de PricewaterhouseCoopers, lamentó que los efectos de las propuestas de Obama pudieran ser drásticos: \textit{Está afectando realmente a la mayoría de las empresas del Fortune 100, que dependen en gran medida del crecimiento de los mercados extranjeros para incrementar sus beneficios totales}.

El sector tecnológico mostró, en particular, su desacuerdo con la propuesta de Obama. El director ejecutivo de Microsoft, Óscarn Ballmer, amenazó con trasladar empleados al extranjero si el Congreso seguía adelante con los planes de Obama de imponer impuestos más elevados sobre los beneficios obtenidos fuera del país. Phil Bond, presidente de TechAmerica, el mayor grupo de defensa de la industria de alta tecnología del país, afirmó: \textit{Esto} [las disposiciones fiscales actuales que afectan a las empresas que operan en el extranjero] \textit{es la piedra angular de la competitividad estadounidense.} [...]\textit{ Las disposiciones fiscales relativas a los ingresos obtenidos en el extranjero son fundamentales para permitir que nuestras empresas compitan y tengan éxito fuera del país}. Quizá como consecuencia de esta fuerte oposición, el plan del presidente Obama no había estado, en el momento de redactarse este texto, ni siquiera cerca de convertirse en realidad.

Las intensas reacciones suscitadas por el debate sobre la tributación de las sociedades se debieron, en parte, al reconocimiento de que el papel del impuesto sobre sociedades en la recaudación pública ha cambiado de manera drástica en los últimos años. En Estados Unidos, las sociedades tributan por sus beneficios netos, es decir, por la diferencia entre lo que ingresan y lo que gastan en factores de producción. Según la legislación, la mayoría de las sociedades están sujetas a un tipo impositivo marginal del $35\%$. En la práctica, numerosas características del sistema del impuesto sobre sociedades (a menudo denominadas lagunas fiscales o loopholes) hacen que los tipos impositivos efectivos sean mucho más bajos que ese porcentaje, y el uso cada vez mayor de estas lagunas ha contribuido, al menos en parte, al importante descenso de la recaudación procedente del impuesto sobre sociedades. En 1960, por ejemplo, casi una cuarta parte de los ingresos federales procedía del impuesto sobre sociedades, como muestra el panel izquierdo de la Figura 24-1. Esta proporción ha cambiado de forma drástica durante las últimas décadas, y actualmente el impuesto sobre sociedades aporta el 15,1 % de los ingresos federales, como muestra el panel derecho.

Para sus detractores, el impuesto sobre sociedades constituye un importante obstáculo para la productividad del sector empresarial, y la reducción de la carga tributaria soportada por las sociedades ha sido beneficiosa para la economía, al haber llevado a las empresas a incrementar su inversión en activos productivos. Para sus defensores, en cambio, el impuesto sobre sociedades constituye una salvaguardia fundamental del carácter progresivo de nuestro sistema tributario. Los partidarios de este impuesto sostienen que, al permitir que el sistema del impuesto sobre sociedades se haya ido erosionando con el tiempo, hemos enriquecido a los propietarios del capital a costa del resto de los contribuyentes.

A continuación evaluaremos estos argumentos. Comenzamos analizando la naturaleza de las sociedades y las razones para establecer un impuesto sobre sociedades. Presentamos la estructura del impuesto sobre sociedades en Estados Unidos y examinamos las dificultades que plantea la definición de los gastos deducibles de las empresas. Asimismo, analizamos cómo aplicar los principios expuestos en el Capítulo 19 para evaluar la incidencia última del impuesto sobre sociedades. 
}

Existen dos importantes efectos sobre el comportamiento derivados del impuesto sobre sociedades: (i) el impacto de la tributación de las sociedades sobre las decisiones de inversión de las empresas; y, (ii) el efecto conjunto de la tributación de las sociedades y de la tributación individual sobre las decisiones de las empresas acerca de cómo financiar sus actividades empresariales. A continuación, presentaremos el marco general a través del cual podemos estudiar el segundo efect, por estar este estrechamente vinculado al crecimiento económico y la oferta de bienes y servicios en la economía.\footnote{
\textbf{¡IMPORTANTE!} Disclaimer.- Óscar analizar la tributación de las sociedades, es importante señalar que no todos los bienes y servicios son producidos por el sector societario, más aún en España. Existe un amplio y sólido sector no societario, integrado por trabajadores por cuenta propia, sociedades personalistas y otras formas de organización distintas a las societarias. El sector no societario representa aproximadamente una quinta parte de la producción española.
}

Como sabemos, las empresas crecen realizando inversiones y obteniendo los rendimientos derivados de esas inversiones. Sin embargo, para invertir en bienes de capital, como nueva maquinaria, una empresa debe financiar, es decir, obtener los fondos necesarios para llevar a cabo dicha inversión. Existen tres posibles vías de financiación empresarial:

\imagenfit{Fig12.png}{Inversión corporativa: Fuentes de financiación}{}

En primer lugar, la empresa puede pedir préstamos a prestamistas, como los bancos: financiación mediante deuda. Este endeudamiento suele realizarse mediante la emisión de bonos corporativos, que son promesas de la sociedad de efectuar pagos periódicos de intereses (junto con la devolución final del principal) a los tenedores de los bonos (los prestamistas). En segundo lugar, la empresa puede emitir acciones a cambio de dinero que pueda utilizarse para realizar inversiones: financiación mediante capital propio. Los inversores que compran acciones de una empresa pueden obtener rendimientos de dos formas distintas. Una consiste en recibir un dividendo, un pago periódico realizado por la empresa por cada acción poseída. La otra consiste en obtener una ganancia de capital sobre las acciones mantenidas cuando el precio de estas aumenta por encima del precio al que fueron adquiridas. Por último, las empresas pueden financiar sus inversiones con cargo a sus propios beneficios retenidos: financiación por fondos propios. De hecho, en los últimos años existe cada vez más evidencia de que una gran parte, quizá la mayoría, de las nuevas inversiones se financia mediante fondos propios.

Analicemos las consecuencias del impuesto sobre sociedades en términos de eficiencia. En particular, existe una preocupación importante por el hecho de que el impuesto sobre sociedades pueda reducir el volumen de inversión realizado por las empresas. Para comprender el efecto de la tributación de las sociedades sobre las decisiones de inversión, comenzamos modelizando la decisión de inversión en un mundo sin impuesto sobre sociedades. 

La decisión de inversión viene determinada por el hecho de que las empresas igualan, en cada período, los beneficios marginales y los costes marginales de la inversión; la empresa estima el rendimiento que obtendrá de su inversión en cada período (el beneficio), lo compara con el coste de la inversión en ese mismo período y solo invierte si los beneficios son superiores a los costes.

Supongamos que cada dólar invertido en una máquina produce $\text{PMg}_K$ céntimos adicionales de producción en cada período, que constituye el beneficio de la inversión. En cada período, la máquina pierde valor en una cantidad lineal por dólar invertido (por ejemplo, 10 centavos por cada dólar invertido en maquinaria). Sin embargo, esta depreciación no representa el coste total de la máquina para la empresa durante un período; la empresa también debe financiar la compra de la máquina. Como analizamos anteriormente, una empresa puede financiar sus inversiones de diversas maneras. Supongamos que esta empresa financia su inversión mediante la venta de acciones de la empresa. A cambio, debe realizar pagos de dividendos de $b_v$ por cada dólar obtenido mediante esa financiación en cada período. En este caso, el coste total de la máquina en cada período es: $\delta_t+v_t=\delta+b_v$.

Si la tasa de depreciación es del $10\%$, y la empresa distribuye cada período dividendos equivalentes al $10\%$ del valor de la inversión (la rentabilidad por dividendo), entonces el coste por período de invertir $1€$ en una máquina es: $\delta + b_v = 0,10+ 0,10= 0,20$€. Gráficamente, podemos también mostrar cómo las curvas de coste marginal y beneficio marginal de la inversión determinan el volumen de inversión realizado ($K$).

\imagenfit{Fig9.png}{Impuesto de Sociedades: Decisiones sobre inversión}{\textit{Public Finance and Policy making}. GRUBER (2017)}

La curva de beneficio marginal mide el rendimiento efectivo obtenido por cada dólar invertido en cada período, $\text{PMg}_K$. El beneficio marginal disminuye a medida que aumenta la inversión, debido al supuesto de producto marginal decreciente; cada euro adicional invertido genera un rendimiento menor porque el producto marginal de cada euro adicional de capital es menor. La curva de coste marginal mide el rendimiento exigido para cada dólar adicional invertido en cada período, es decir, cuánto debe producir esa inversión en cada período para cubrir sus costes (de depreciación y financiación). El coste marginal es constante e igual a: $\delta + b_v = 0,20$, como muestra la curva $MC_1$.

Las empresas invierten hasta que el dólar marginal invertido genera unos costes y unos beneficios iguales, lo que ocurre en el punto $A$, con un nivel de inversión $K_!$. Las empresas invierten hasta que el siguiente euro invertido produce exactamente el rendimiento suficiente ($0,20$) para cubrir sus costes en cada período. Si una empresa invierte menos que esa cantidad (a la izquierda de $K_1$), el rendimiento marginal de un dólar adicional de inversión supera su coste marginal de $0,20$€, por lo que la empresa debería invertir más. Si la empresa invierte más que esa cantidad (a la derecha de $K_1$), el rendimiento marginal de ese dólar adicional es inferior a su coste marginal de $0,20$€, por lo que esa inversión no debería realizarse.


\paragraph{Las consecuencias del impuesto sobre sociedades para la inversión}
¿Qué ocurre si introducimos impuestos en este análisis? Supongamos, en primer lugar, que el impuesto sobre sociedades consiste simplemente en un impuesto a un tipo $t$ aplicado sobre los ingresos monetarios menos los costes laborales (sin ningún tipo de deducción fiscal por las inversiones realizadas).

Los ingresos monetarios generados por cada euro invertido en la máquina durante cada período son $MP_K$. Una vez establecido el impuesto, los ingresos obtenidos por cada dólar invertido pasan a ser: $MP_K \cdot (1 − t)$, porque debe pagarse el nuevo impuesto sobre cada dólar de ingresos.

Esta reducción del rendimiento efectivo provoca un desplazamiento hacia abajo de la curva de beneficio marginal hasta $MB_2$: la tributación de los beneficios empresariales ha reducido el beneficio marginal de invertir. Gráficamente:

\imagenfi{Fig10.png}{Impuesto de Sociedades: Decisiones de inversión en presencia de carga fiscal}{}

Los costes por dólar invertido permanecen en $\delta + b_v$, por lo que la curva de coste marginal permanece en su nivel inicial. La nueva decisión óptima de inversión se encuentra en el punto $B$, y la inversión disminuye hasta $K_2$. Así, una empresa invierte menos cuando el gobierno se apropia de parte de su rendimiento mediante el impuesto sobre sociedades. ¿Por qué? Porque la tasa efectiva de rendimiento después de impuestos obtenida por la empresa debe seguir siendo suficiente para alcanzar la tasa de rendimiento exigida $\delta + b_v$. Como consecuencia, la tasa de rendimiento antes de impuestos debe ser mayor que en ausencia de tributación, y ello solo ocurre si la empresa invierte menos. Por ejemplo, con un tipo impositivo del $50\%$, la empresa debe obtener un rendimiento de $0,40€$ por cada euro invertido para poder cubrir su coste de $0,20$€ correspondiente a la depreciación y a la financiación. Por tanto, la empresa debe invertir menos: debe dejar de invertir en el punto en el que el euro marginal de inversión produce un rendimiento de $0,40€$, en lugar de continuar invirtiendo hasta que dicho euro marginal produzca un rendimiento de $0,20€$. 

En este escenario, la tributación de las sociedades conduce a un menor nivel de inversión.

\begin{specialblock}{\textbf{Adicional.-} Los efectos de las amortizaciones fiscales y del crédito fiscal por inversión sobre la inversión empresarial}
    Esta descripción del efecto de los impuestos sobre la inversión empresarial no incorpora la influencia de las deducciones fiscales a la inversión, como las amortizaciones fiscales o el crédito fiscal por inversión (ITC). Estas deducciones fiscales actúan como descuentos sobre el precio de las inversiones, reduciendo el coste marginal al compensar parte de los costes de financiación y depreciación.

    Recordemos que las amortizaciones fiscales suelen distribuirse entre el año de adquisición del activo y los años posteriores y que, para valorar estos flujos de beneficios, debemos considerar su valor actual descontado (VAD).\footnote{%
        El valor de un determinado calendario de amortización, $z$, es el valor actual descontado del conjunto de deducciones por depreciación asociadas a la compra de una nueva máquina, expresado como fracción del precio de compra de dicha máquina.
    } Si la empresa pudiera imputar inmediatamente como gasto el coste íntegro de la máquina (deduciendo todo su valor en el año de adquisición), entonces: $z = 1,0$, porque la deducción equivaldría al $100\%$ del precio de compra. A medida que las amortizaciones fiscales se distribuyen a lo largo de los años futuros, $z$ disminuye, ya que el valor actual descontado de dichas deducciones disminuye cuanto más alejadas se encuentran en el tiempo.
    
    La depreciación se resta de los beneficios de la empresa al calcular la base imponible del impuesto sobre sociedades. Por tanto, cada euro de amortización fiscal ahorra a la empresa $t€$ en impuestos sobre sociedades, porque reduce la base imponible en un euro. Si las amortizaciones fiscales de una empresa tienen un valor actual descontado (VAD) igual a $z$, su valor para la empresa es de $(t \cdot z)€$. Si una empresa pudiera imputar inmediatamente como gasto su inversión (por ejemplo, si $z = 1$) y el tipo del impuesto sobre sociedades fuera del $35\%$, la amortización correspondiente a cada euro del precio de compra de la máquina tendría para la empresa un valor de $0,35€$ en forma de reducción del impuesto a pagar.
    
    Por tanto, la existencia de las amortizaciones fiscales y del crédito fiscal actúa como un reembolso sobre el coste de invertir. En esencia, el coste de cada dólar invertido en una máquina durante cada período, que anteriormente era $\delta + b_v$, pasa a ser: $(\delta+ b_v) × (1 − [t \cdot z] − a)$ porque la empresa puede compensar estos incentivos fiscales con parte de los costes derivados de la inversión. Gráficamente:
    
    \imagenfit{Fig11.png}{Impuesto de Sociedades: Decisiones de inversión con carga fiscal y amortizaciones}{}
    
    La curva de beneficio marginal relevante para la empresa es la curva de beneficio marginal después de impuestos, $MB_2$. La reducción del rendimiento requerido derivada de las amortizaciones fiscales y del crédito fiscal desplaza la curva de coste marginal desde $MC_1$ hasta $MC_2$, cuyo valor pasa a ser: $(\delta + b_v) × (1 − [t \cdot z] − a)$.
    
    Para ilustrar estos cambios numéricamente, volvamos al ejemplo utilizado anteriormente: con una tasa de depreciación por período del $10\%$ ($\delta$) y una rentabilidad por dividendo del $10\%$ ($b_v$), la curva de coste marginal antes de impuestos se situaba en $0,20€$. Supongamos que el calendario de amortización correspondiente a esta inversión es tal que el valor actual descontado de las amortizaciones fiscales equivale a la mitad del precio de compra de la máquina ($z = 0,5$), que el tipo del impuesto sobre sociedades es del $35\%$, y que existe un crédito fiscal del $10\%$. En ese caso, el coste de cada euro invertido en la máquina durante cada período (es decir, el rendimiento exigido por el inversor sobre ese dólar invertido) pasa a ser: $0,20 \cdot (1 − [0,35 \cdot 0,5] − 0,1) = 0,145€$.
    
    El nuevo equilibrio se encuentra en el punto $C$, con un nivel de inversión $K_3$, punto en el que los costes marginales después de impuestos y los beneficios marginales después de impuestos correspondientes a cada euro invertido en maquinaria durante cada período son iguales; es decir, donde los rendimientos efectivos (curva $MB_2$) son iguales a los rendimientos exigidos (curva $MC_2$) para el euro marginal de inversión. Por tanto, el nivel de inversión es superior al existente antes de introducir las amortizaciones fiscales y el crñedito fiscal ($K_2$), pero sigue siendo inferior al existente antes de la tributación ($K_1$) (aunque, en algunos casos, puede llegar a ser superior al existente antes de la tributación). La reducción de los beneficios marginales derivada de que los rendimientos sean percibidos después de impuestos es mayor que la reducción de los costes marginales derivada de las amortizaciones fiscales y del ITC. Por ello, la inversión sigue disminuyendo en términos globales como consecuencia de la tributación de las sociedades.
\end{specialblock}



Tipo efectivo del impuesto sobre sociedades

Ahora que comprendemos cómo afectan los impuestos a las decisiones de inversión de una empresa, podemos resumir matemáticamente el impacto neto del sistema tributario sobre dichas decisiones.

El efecto de los impuestos se resume mediante el tipo efectivo del impuesto sobre sociedades, definido como el porcentaje de incremento en la tasa de rendimiento antes de impuestos del capital que resulta necesario debido a la tributación.

Cuando se introducen impuestos, la tasa de rendimiento obtenida por la empresa sobre sus inversiones debe aumentar para poder financiar el pago de dichos impuestos. El incremento necesario depende del tipo impositivo, del tratamiento fiscal de la depreciación y de la existencia del crédito fiscal. Estos factores determinan conjuntamente el efecto global de la tributación sobre las decisiones de inversión.

Consideremos, en primer lugar, el sistema simplificado del impuesto sobre sociedades representado en la Figura 24-5, con un impuesto del 35 % sobre los beneficios empresariales y sin amortizaciones fiscales ni ITC.

Antes de la tributación, la tasa efectiva de rendimiento de la empresa debe ser, como mínimo, igual a la tasa de rendimiento requerida del 20 %. Después de introducir el impuesto sobre sociedades del 35 %, la tasa efectiva de rendimiento de la empresa debe ser al menos:

0,20 / (1 − 0,35) = 0,307 = 30,7 %.

La empresa debe obtener una tasa de rendimiento un 35 % superior sobre sus inversiones para poder pagar los impuestos y, al mismo tiempo, satisfacer la tasa de rendimiento necesaria para cubrir los costes de depreciación y financiación.

Como consecuencia, el tipo efectivo del impuesto sobre sociedades coincide con el tipo legal del impuesto sobre sociedades, es decir, el 35 %.

De forma más general, el tipo efectivo del impuesto sobre sociedades (ETR) se mide mediante la siguiente expresión:

ETR = [MPK (después de impuestos) − MPK (antes de impuestos)] / MPK (después de impuestos).

Como consecuencia del impuesto sobre sociedades, el producto marginal correspondiente a cada dólar de inversión (MPK) debe ser mayor para poder pagar los impuestos y satisfacer la tasa de rendimiento requerida. Cuando no existen amortizaciones fiscales ni ITC, el MPK debe ser un 35 % mayor como consecuencia de la tributación.

Consideremos ahora el sistema del impuesto sobre sociedades más realista representado en la Figura 24-6, con un tipo del impuesto sobre sociedades del 35 %, así como amortizaciones fiscales y un crédito fiscal por inversión.

En este caso, la empresa debe obtener una tasa de rendimiento más elevada para pagar los impuestos sobre sus beneficios, pero puede conformarse con una tasa de rendimiento menor porque el coste de la máquina está subvencionado mediante las amortizaciones fiscales y el ITC.

Con amortizaciones fiscales y el ITC, el tipo efectivo del impuesto es:

ETR = (t − tz − a) / (1 − tz − a)

Por ejemplo, si z y a son ambos iguales a cero (el caso en el que no existen amortizaciones fiscales ni ITC), el tipo efectivo del impuesto es simplemente t, es decir, el tipo legal.

En nuestro ejemplo, con z = 0,5 y un ITC = 0,1, el tipo efectivo del impuesto es:

ETR = (0,35 − 0,35 × 0,5 − 0,1) / (1 − 0,35 × 0,5 − 0,1)

= 0,075 / 0,725

= 10,3 %.

Debido al valor de las deducciones por depreciación y del ITC, el tipo efectivo del impuesto pasa a ser únicamente del 10,3 %, a pesar de que el tipo legal es del 35 %.

Es decir, la empresa únicamente necesita obtener una tasa de rendimiento antes de impuestos un 10,3 % superior a la que obtenía antes de la tributación para satisfacer la tasa de rendimiento requerida del 20 %.

Este porcentaje es inferior al tipo legal porque el ITC y las amortizaciones fiscales compensan parte del efecto de la tributación.

Tipos efectivos de imposición negativos

Con un valor suficientemente elevado de z y del ITC a, el tipo efectivo del impuesto sobre sociedades puede llegar incluso a ser negativo. Si la empresa de nuestro ejemplo pudiera imputar inmediatamente como gasto la totalidad del coste de la máquina y deducir el coste íntegro de la máquina en el momento de su adquisición, de modo que z = 1, entonces el tipo efectivo del impuesto sería: $ETR = (0,35 − 0,35 × 1 − 0,1) / (1 − 0,35 × 1 − 0,1)= −0,1 / 0,55= −18,2 \%.$

Un tipo efectivo del impuesto negativo significaría que la curva de coste marginal (MC) desciende tanto en la Figura 24-6 que el punto C queda situado a la derecha del punto A, y la empresa invierte más con impuestos (K₃) de lo que invertía sin ellos (K₁).

(Un tipo efectivo negativo del impuesto sobre sociedades es más probable cuando la empresa financia su inversión mediante deuda en lugar de mediante capital propio. Recordemos que los pagos de intereses correspondientes a la deuda de una empresa son deducibles fiscalmente. Esta deducción reduce todavía más el coste de financiar la inversión, porque la financiación mediante deuda cuesta únicamente r × (1 − t) en lugar del coste íntegro r.)



--------------------------------------------------------------------------------------------------------



Implicaciones de política económica del impacto del impuesto sobre sociedades sobre la inversión

El análisis matemático del apartado anterior muestra la diversidad de efectos que las distintas estructuras del impuesto sobre sociedades pueden tener sobre las decisiones de inversión de las empresas.

Para un mismo tipo del impuesto sobre sociedades, el sistema tributario puede diseñarse de forma que ofrezca incentivos muy distintos para la inversión.

Un sistema tributario que simplemente gravara la renta empresarial, sin deducciones por depreciación ni crédito fiscal por inversión, reduciría claramente el nivel de inversión al disminuir el beneficio marginal (el rendimiento después de impuestos) de cada dólar invertido.

Permitir amortizaciones fiscales —especialmente si son aceleradas— y créditos fiscales por inversión puede mitigar, e incluso revertir, el efecto negativo de los impuestos sobre la inversión al reducir el rendimiento exigido a la inversión.

Cuando una empresa compra una máquina, recibe las ventajas fiscales incorporadas a la adquisición de esa máquina. Cuanto más rápida sea la depreciación y mayor sea el crédito fiscal por inversión, mayores serán esas ventajas fiscales.

Como consecuencia, el tipo marginal efectivo del impuesto sobre sociedades ha variado considerablemente en el pasado reciente, desde un máximo del 51 % en 1980 hasta un mínimo del 27 % en 2003.

Evidencia sobre impuestos e inversión

Existe una extensa literatura que analiza el impacto de los impuestos sobre sociedades en las decisiones de inversión de las empresas.

La conclusión de los estudios más recientes es que la inversión responde de forma bastante sensible a los incentivos fiscales, con una elasticidad de la inversión respecto al tipo efectivo del impuesto del orden de −0,5: cuando los impuestos reducen el coste de la inversión en un 10 %, la inversión aumenta aproximadamente un 5 %.

Esta elasticidad, de magnitud considerable, sugiere que la política del impuesto sobre sociedades puede constituir un instrumento muy poderoso para determinar el nivel de inversión y que el impuesto sobre sociedades dista mucho de ser un verdadero impuesto sobre los beneficios puros.


------------------------------------------------------------------------------------------------

\begin{specialblock}{Imposición: Efectos sobre la financiación de la inversión empresarial}
    Las consecuencias del impuesto sobre sociedades para la financiación

    El impuesto sobre sociedades tiene implicaciones en términos de eficiencia que van más allá de las decisiones de las empresas sobre cuánto invertir. Otra decisión potencialmente importante que puede verse influida por la tributación de las sociedades es cómo financiar esa inversión.
    
    Como se señaló anteriormente, las empresas disponen de tres opciones principales para financiar sus inversiones: pueden utilizar beneficios retenidos, pueden aumentar su endeudamiento mediante préstamos o pueden emitir capital propio (participaciones de propiedad). Además, si una empresa opta por financiarse mediante capital propio, dispone de dos formas de retribuir a los inversores: puede distribuir dividendos o puede reinvertir los beneficios e intentar aumentar el valor de las acciones, generando así ganancias de capital.
    
    Por el momento, dejaremos de lado el caso menos interesante de financiar la inversión mediante beneficios retenidos y nos centraremos en el efecto de los impuestos sobre la financiación de nuevas inversiones mediante nuevo capital.
    
    El impacto de los impuestos sobre la financiación
    
    Supongamos que una empresa necesita 10€ para realizar una inversión que producirá 1 dólar de renta empresarial cada año. La empresa desea financiar esta inversión mediante deuda o mediante capital propio y, en cualquiera de los dos casos, transferir ese dólar al inversor.
    
    La Figura 24-7 ilustra el impacto de los impuestos sobre la decisión de la empresa acerca de cómo financiar dicha inversión (ignorando, una vez más y por el momento, la posibilidad de financiarla mediante beneficios retenidos).
    
    En el primer nodo, la empresa decide si endeudarse o emitir acciones para financiar la inversión de 10€.
    
    Si la empresa se financia mediante deuda (la rama superior), pagará el dólar en concepto de intereses a los tenedores de bonos. En este caso, la empresa no pagará impuestos sobre ese dólar de renta empresarial porque los pagos de intereses son deducibles de la renta de la empresa a efectos del impuesto sobre sociedades: el dólar de renta empresarial queda compensado por un dólar de pagos de intereses al calcular los impuestos.
    
    Por consiguiente, el tenedor del bono que ha prestado el dinero a la empresa recibe el dólar íntegro, pero posteriormente debe pagar el impuesto personal correspondiente sobre ese dólar al tipo del impuesto sobre la renta aplicable a los intereses, tᵢₙₜ, de modo que finalmente recibe:
    
    1 dólar × (1 − tᵢₙₜ).
    
    Como alternativa, la empresa puede emitir capital propio para financiar la inversión.
    
    En este caso, cuando la empresa obtiene el dólar procedente de la inversión, ya no dispone del pago de intereses que serviría como deducción compensatoria frente al impuesto sobre sociedades, por lo que deberá pagar el impuesto sobre ese dólar.
    
    En consecuencia, únicamente:
    
    1 dólar × (1 − t꜀),
    
    donde t꜀ representa el tipo del impuesto sobre sociedades, llegará a los accionistas.
    
    Si la empresa utiliza financiación mediante capital propio, debe tomar además otra decisión: ¿cómo debe transferir ese dinero a los accionistas?
    
    Con la deuda no existe elección posible: el dólar simplemente se paga en forma de intereses.
    
    Con el capital propio, en cambio, la empresa puede optar entre pagar un dividendo o retener los beneficios y reinvertirlos.
    
    Si distribuye el dólar como dividendo, el accionista deberá pagar el impuesto correspondiente sobre los dividendos.
    
    En consecuencia, el accionista recibirá finalmente:
    
    1 dólar × (1 − t꜀) × (1 − tdiv),
    
    donde tdiv representa el tipo del impuesto personal sobre la renta por dividendos.
    
    Esta característica suele denominarse doble imposición de los dividendos, ya que estos son gravados tanto por el impuesto sobre sociedades como por el impuesto personal sobre la renta.
    
    Si, por el contrario, la empresa reinvierte ese dólar, el accionista podrá obtener una ganancia de capital cuando venda sus acciones.
    
    Supongamos que cada dólar reinvertido incrementa el precio de la acción en 1 dólar.
    
    En ese caso, el accionista recibe:
    
    1 dólar × (1 − t꜀) × (1 − tcg),
    
    donde tcg representa el tipo efectivo del impuesto sobre las ganancias de capital.
    
    El tipo efectivo del impuesto sobre las ganancias de capital es inferior al tipo legal para la mayoría de los inversores debido al conjunto de ventajas fiscales aplicables a las ganancias de capital analizadas en el Capítulo 23, entre ellas la tributación en el momento de la realización en lugar de en el momento del devengo y el step-up de la base fiscal en el momento del fallecimiento.
    
    Además, tradicionalmente el tipo efectivo sobre las ganancias de capital ha sido considerablemente inferior al tipo aplicable a los dividendos.
    
    Desde 2004, los tipos legales aplicables a las ganancias de capital y a los dividendos son iguales (15 %), pero el tipo efectivo sobre las ganancias de capital continúa siendo inferior al correspondiente a los dividendos.
    
    Como debería quedar claro a partir de la Figura 24-7, la estructura del impuesto sobre sociedades puede desempeñar un papel decisivo en las decisiones de las empresas acerca de cómo financiar sus inversiones.
    
    La financiación mediante deuda evita la tributación societaria sobre los beneficios, pero no permite a los inversores aprovechar los bajos tipos aplicables a las ganancias de capital.
    
    La financiación mediante capital propio y el reparto de dividendos parecen constituir, desde el punto de vista fiscal, la alternativa menos eficiente de todas.
    
    Por ello, esta figura plantea dos cuestiones fundamentales sobre la financiación empresarial.
    
    ¿Por qué no financiarse exclusivamente con deuda?
    
    Una empresa dispone de dos formas de financiar una inversión.
    
    Si aumenta su endeudamiento, los intereses que paga para financiar dicha inversión están exentos del impuesto sobre sociedades.
    
    Si aumenta su capital propio, debe pagar el impuesto sobre sociedades.
    
    Entonces, ¿por qué las empresas no financian todas sus inversiones mediante deuda y evitan así la tributación societaria?
    
    Existen diversas respuestas específicas a esta cuestión, pero todas ellas se reducen a la diferencia fundamental entre deuda y capital propio: la deuda exige pagos fijos, mientras que el capital propio no.
    
    Una empresa puede decidir si paga o no dividendos sobre las acciones, pero no puede decidir si paga o no los intereses de su deuda: está obligada a hacerlo.
    
    Si la empresa deja de efectuar los pagos de intereses, incurre en incumplimiento (default) y puede verse obligada a declararse en quiebra.
    
    En caso de quiebra, todos los inversores de la empresa pueden reclamar derechos sobre los activos que queden disponibles.
    
    Los tenedores de bonos son los primeros acreedores; tienen prioridad para reclamar los activos restantes.
    
    Los accionistas, por el contrario, únicamente tienen derecho a los recursos que queden una vez satisfechos los derechos de los tenedores de bonos.
    
    Imaginemos que usted pudiera elegir el sistema de evaluación de esta asignatura y dispusiera de dos opciones.
    
    En una de ellas, tendría que entregar un problema cada semana y, si suspendiera uno solo, suspendería toda la asignatura.
    
    En la otra, también entregaría un problema cada semana, pero la nota final dependería de la media de todas las calificaciones obtenidas.
    
    La mayoría elegiríamos este segundo sistema porque ofrece una flexibilidad mucho mayor.
    
    Del mismo modo, esa mayor flexibilidad del capital propio puede aumentar suficientemente su valor como para compensar su desventaja fiscal.
    
    La financiación mediante capital propio puede proporcionar una zona de amortiguación (buffer zone) frente al riesgo de quiebra: si la situación económica empeora, la empresa simplemente puede dejar de repartir dividendos en lugar de verse obligada a declararse en quiebra.
    
    Los directivos, en particular, que con frecuencia se preocupan más por reducir el riesgo inmediato de perder su puesto de trabajo que por maximizar el rendimiento para los propietarios de la empresa, pueden preferir la financiación mediante capital propio frente a la financiación mediante deuda.
    
    Esta preferencia por el capital propio, pese a sus desventajas fiscales, se ve reforzada por un tipo particular de problema de agencia dentro de la empresa: el conflicto de intereses entre los acreedores y los accionistas.
    
    Los acreedores reciben un rendimiento fijo (los pagos de intereses), independientemente de lo bien que funcione la empresa, siempre que esta no quiebre.
    
    Los accionistas, por el contrario, reciben un rendimiento vinculado al desempeño de la empresa: la parte de los beneficios reinvertida en la empresa aumenta de valor a medida que aumenta el valor de la propia empresa.
    
    En la Figura 24-7 supusimos que cada dólar reinvertido generaba únicamente un dólar adicional de valor.
    
    En la realidad, cada dólar invertido produce un rendimiento incierto.
    
    Los accionistas pierden cuando una inversión empresarial obtiene malos resultados, pero se benefician cuando dicha inversión tiene éxito.
    
    Sin embargo, existe un límite inferior a las pérdidas que pueden sufrir los accionistas cuando la empresa obtiene malos resultados: como máximo pueden perder la inversión inicial si la empresa quiebra.
    
    En cambio, no existe un límite superior a las ganancias que pueden obtener cuando la empresa obtiene buenos resultados.
    
    Por ello, cuando los accionistas poseen únicamente una pequeña parte de la empresa, tenderán a asumir riesgos excesivos, porque obtienen una mayor rentabilidad potencial soportando relativamente poco riesgo; al poseer solo una pequeña fracción de la empresa, no tienen demasiado capital que perder.
    
    La posibilidad de asumir riesgos excesivos constituye un problema porque son los accionistas quienes toman las decisiones sobre a quién contratar y, potencialmente, sobre qué proyectos emprender.
    
    Considérese una empresa que actualmente tiene 1 millón de€ de capital propio y 5 millones de€ de deuda, como se muestra en el panel superior de la Tabla 24-1, lo que da un valor total de la empresa de 6 millones de€.
    
    El pago de intereses de la empresa sobre esa deuda en cada período es de 500.000€.
    
    La empresa obtiene actualmente unos beneficios de 600.000€ al año, de modo que cada año puede pagar los intereses que debe a sus acreedores y distribuir los 100.000€ restantes de beneficios en forma de dividendos a sus accionistas.
    
    Supongamos que, además de este flujo estable de beneficios, a la empresa se le ofrece un proyecto que tiene un 50 % de probabilidad de generar 3 millones de€ y un 50 % de probabilidad de perder 6 millones de€.
    
    Si la inversión fracasa, la empresa quebrará porque tendría que vender sus 6 millones de€ en activos (5 millones de deuda y 1 millón de capital propio) para cubrir la pérdida de 6 millones de€.
    
    Desde la perspectiva de la empresa en su conjunto, este no es un riesgo que convenga asumir.
    
    El rendimiento esperado de esta inversión —la probabilidad de éxito multiplicada por el beneficio si el proyecto tiene éxito, más la probabilidad de fracaso multiplicada por la pérdida si fracasa— es:
    
    (0,5 × 3 millones de€) + (0,5 × (−6 millones de€)) = −1,5 millones de€.
    
    Por tanto, la empresa estaría mejor si no llevara a cabo la inversión, ya que, en promedio, pierde dinero.
    
    Ahora analicemos el proyecto desde la perspectiva de los accionistas.
    
    Si la inversión tiene éxito, los accionistas reciben la totalidad de los 3 millones de€, porque obtienen todos los beneficios derivados de la mejora en el desempeño de la empresa.
    
    Si la inversión fracasa y la empresa quiebra, pierden el millón de€ que habían invertido, además de los 100.000€ anuales de dividendos que dejarán de percibir en el futuro.
    
    Con un tipo de interés del 10 %, el valor actual descontado de esa corriente futura de dividendos asciende a 1 millón de€, por lo que los accionistas pierden un total de 2 millones de€ si la empresa entra en quiebra.²⁰
    
    Desde la perspectiva de los accionistas, existe un 50 % de probabilidad de perder 2 millones de€ y un 50 % de probabilidad de ganar 3 millones de€, lo que proporciona un valor esperado de:
    
    (0,5 × (−2 millones de€)) + (0,5 × 3 millones de€) = 500.000€.
    
    Por ello, los accionistas votarán a favor de emprender el proyecto.
    
    ¿Y qué ocurre con los acreedores?
    
    Ellos no obtienen ningún beneficio si la apuesta sale bien, porque la empresa ya genera beneficios suficientes para pagar íntegramente los intereses de la deuda.
    
    Pero, si la apuesta sale mal y la empresa quiebra, no solo pierden sus 500.000€ anuales en intereses, sino también los 5 millones de€ que habían prestado inicialmente.
    
    Desde su punto de vista, se trata de una apuesta con un 50 % de probabilidad de no ganar nada (si el proyecto tiene éxito, seguirán recibiendo únicamente sus 500.000€ de intereses) y un 50 % de probabilidad de perder 5 millones de€ en activos, además del valor actual descontado de 5 millones de€ correspondientes a los intereses futuros.
    
    Para los tenedores de bonos, el valor esperado de esta inversión es de −5 millones de€.
    
    En este caso existe un claro conflicto de intereses entre accionistas y acreedores.
    
    Los accionistas, que son los propietarios de la empresa y controlan a quienes toman las decisiones —aunque no necesariamente todas las decisiones concretas—, apoyarían la realización del proyecto, mientras que los acreedores se opondrían.
    
    Este problema surge porque la participación del capital propio es muy reducida.
    
    Si el capital propio ascendiera a 5 millones de€ y la deuda únicamente a 1 millón de€, los accionistas ya no apoyarían este proyecto.
    
    Esto se ilustra en el panel inferior de la Tabla 24-1 (suponiendo ahora que la empresa paga 500.000€ en dividendos a los accionistas y 100.000€ en intereses a los acreedores).
    
    En esta situación, los accionistas no desean asumir la nueva inversión porque tienen demasiado que perder: el rendimiento esperado del proyecto para ellos pasa a ser de −3,5 millones de€.
    
    Como antes, los acreedores tampoco desean realizar esta inversión arriesgada porque su valor esperado continúa siendo negativo.
    
    En este caso, los intereses de accionistas y acreedores quedan alineados (ambos pierden dinero), y el proyecto no se lleva a cabo.
    
    La idea fundamental es que cuanto mayor es la proporción de financiación mediante deuda, mayor es el potencial para este conflicto de intereses.
    
    Ese potencial aumenta porque, a medida que la proporción de deuda crece, los acreedores soportan una parte cada vez mayor de las pérdidas derivadas de los proyectos que terminan en quiebra, mientras que los accionistas asumen un riesgo cada vez menor al aceptar inversiones arriesgadas.
    
    Supongamos ahora que usted trabaja para un banco que está considerando conceder un préstamo a una pequeña empresa.
    
    ¿Preferiría prestar dinero a una empresa con un 50 % de capital propio y un 50 % de deuda, o a otra con un 10 % de capital propio y un 90 % de deuda?
    
    Siguiendo la lógica del ejemplo anterior, preferiría prestar a la empresa con un 50 % de capital propio, porque una mayor parte del capital de los propietarios está en riesgo y, por tanto, estos tomarán decisiones más prudentes.
    
    Le preocuparía que una empresa con solo un 10 % de capital propio asumiera riesgos excesivos con un valor esperado negativo para usted como acreedor.
    
    Como consecuencia de este problema de agencia, los bancos (y otros prestamistas) exigirán tipos de interés más elevados a las empresas a medida que aumente la proporción de financiación mediante deuda.
    
    Estos tipos de interés más elevados compensan la ventaja fiscal de la deuda, de modo que las empresas se enfrentan ahora a una disyuntiva (trade-off): una mayor deuda implica una mayor ventaja fiscal, pero también unos costes de financiación potencialmente más altos porque los bancos desconfían de las empresas excesivamente endeudadas.
    
    La paradoja de los dividendos
    
    El segundo gran misterio de la financiación empresarial es por qué las empresas pagan dividendos.
    
    Una vez que la empresa ha optado por financiarse mediante capital propio, si obtiene 1 dólar de beneficios puede distribuir ese dólar como dividendo, desencadenando el pago del impuesto sobre los dividendos, o puede reinvertirlo, aumentando el valor de las acciones y permitiendo que el accionista se beneficie de los tipos impositivos preferenciales aplicables a las ganancias de capital cuando venda sus acciones.
    
    El hecho de que los tipos efectivos sobre las ganancias de capital hayan sido tradicionalmente mucho más bajos que los tipos aplicables a los dividendos sugiere que las empresas deberían reinvertir sus beneficios en lugar de repartir dividendos, siempre que puedan encontrar oportunidades de reinversión productivas.
    
    Sin embargo, aproximadamente una quinta parte de las empresas que cotizan en bolsa pagan dividendos, aunque este porcentaje ha disminuido durante las dos últimas décadas.
    
    Nos encontramos, por tanto, ante un nuevo enigma: ¿por qué las empresas pagan dividendos cuando las ganancias de capital tributan mucho menos?
    
    La evidencia empírica respalda dos explicaciones distintas acerca de por qué las empresas pagan dividendos, tal como revisan Gordon y Dietz (2006).
    
    La primera es una teoría de agencia.
    
    Los inversores están dispuestos a soportar la ineficiencia fiscal de los dividendos para sacar el dinero del control de unos directivos afectados por problemas de agencia.
    
    Un estudio encontró que si las 25 mayores empresas que repartían dividendos de forma continuada en 2002 nunca hubieran pagado esos dividendos, sus disponibilidades de efectivo, que entonces ascendían a 160.000 millones de€, habrían alcanzado 1,8 billones de€.
    
    Los autores sugieren que unas reservas de efectivo tan enormes podrían constituir una auténtica receta para el desastre debido al comportamiento oportunista de los directivos, dando lugar a inversiones ineficientes o incluso a corrupción manifiesta.
    
    En consecuencia, los accionistas pueden estar dispuestos a soportar el coste derivado de la ineficiencia fiscal de los dividendos con el fin de reducir el control de los directivos sobre los activos de la empresa.²¹
    
    La segunda explicación es una teoría de la señalización (signaling theory).
    
    Los inversores disponen de información imperfecta acerca del verdadero estado de la empresa, por lo que los directivos pagan dividendos para señalar a los inversores que la empresa está obteniendo buenos resultados.
    
    Es decir, el simple hecho de que los directivos estén dispuestos a «quemar dinero» pagando dividendos fiscalmente ineficientes debe demostrar que la empresa tiene dinero de sobra.
    
    Óscar repartir dividendos, los directivos pueden demostrar a unos inversores insuficientemente informados que su inversión está funcionando satisfactoriamente.
    
    Con mucho gusto. Aquí tienes la última parte de este apartado.
    
    ¿Cómo deberían gravarse los dividendos?
    
    Un importante debate permanente en materia de política fiscal se refiere al tratamiento tributario adecuado de la renta procedente de dividendos.
    
    Como hemos visto a lo largo de este capítulo, el tipo impositivo individual aplicable a los dividendos suele ser superior al tipo efectivo aplicable a las ganancias de capital.
    
    De acuerdo con la lógica de la Figura 24-7, un tipo impositivo elevado sobre los dividendos puede tener tres efectos sobre las decisiones de financiación de las empresas.
    
    En primer lugar, puede reducir el uso de los dividendos como forma de retribuir a los accionistas.
    
    Como se revisa en la Aplicación de la página 765, existe evidencia clara de que la reducción de la tributación sobre los dividendos aprobada en 2003 dio lugar a un aumento en el reparto de dividendos.
    
    En segundo lugar, unos impuestos elevados sobre los dividendos pueden inducir a las empresas a elegir financiación mediante deuda en lugar de financiación mediante capital propio.
    
    La evidencia disponible sugiere que unos tipos más altos sobre los dividendos hacen que las empresas dependan en mayor medida de la financiación mediante deuda.²³
    
    Por último, y quizá lo más importante, unos impuestos elevados sobre los dividendos podrían reducir la inversión.
    
    Para aquellos inversores cuya rentabilidad adopta la forma de dividendos, el tipo efectivo de gravamen sobre la inversión aumenta cuando aumenta el tipo impositivo sobre los dividendos.
    
    Este incremento hará que dichos inversores exijan una mayor tasa de rendimiento antes de impuestos sobre las inversiones, desplazando hacia arriba la curva de coste marginal de las Figuras 24-4, 24-5 y 24-6 y reduciendo la inversión (porque, al existir producto marginal decreciente, las empresas deben reducir la inversión para obtener una mayor tasa efectiva de rendimiento sobre el siguiente dólar invertido).
    
    Este resultado ha llevado a muchos autores a sostener que la doble imposición de los dividendos reduce el nivel de inversión empresarial en Estados Unidos.
    
    Por tanto, reducir el impuesto sobre los dividendos podría constituir, en principio, un instrumento muy eficaz para incrementar el nivel de inversión empresarial del país.
    
    Sin embargo, esta predicción se deriva de un modelo que ni siquiera puede explicar por qué las empresas pagan dividendos en primer lugar.
    
    De hecho, tanto el modelo de agencia como el modelo de señalización predicen que unos tipos impositivos más elevados sobre los dividendos conducirán a una mayor inversión, y no a una menor.
    
    Ello se debe a que, en ambos modelos, los dividendos constituyen un mecanismo costoso o ineficiente, y gravarlos con tipos más elevados reduce su utilización; con independencia de que los dividendos sirvan para resolver problemas de agencia o para transmitir señales al mercado, cuando aumenta su coste (como consecuencia de una mayor tributación), las empresas recurrirán menos a ellos para esos fines.
    
    Cuando las empresas distribuyen menos dividendos, disponen de más beneficios retenidos.
    
    La literatura sobre financiación empresarial ha sugerido que la disponibilidad de beneficios retenidos constituye un determinante importante —y quizá el más importante— de la financiación de las inversiones.
    
    (La hipótesis de que los beneficios retenidos por las empresas constituyen una fuente fundamental de financiación de la inversión fue contrastada por primera vez por Fazzari et al. (1988). Lamont (1997) aporta evidencia adicional muy sólida en apoyo de esta hipótesis, mostrando que la inversión de las compañías petroleras disminuyó cuando disponían de menos liquidez debido a la caída de los precios del petróleo. Rauh (2006) mostró asimismo que la inversión de las empresas disminuyó cuando el gobierno les obligó a realizar mayores aportaciones a sus fondos de pensiones.)
    
    Por tanto, gravar con mayor intensidad estos pagos de dividendos ineficientes conduciría a más inversión, y no a menos.
    
    La mejor evidencia disponible hasta la fecha, que se revisará más adelante, sugiere que, en conjunto, los impuestos sobre los dividendos tienen muy poco impacto sobre la inversión.
    
    Integración del impuesto sobre sociedades
    
    Aunque diversas reformas legislativas han reducido el tipo impositivo aplicable a los dividendos, sigue existiendo una cuestión más profunda:
    
    ¿Por qué debería el sistema tributario tratar de forma diferente la renta empresarial dependiendo de cómo se distribuya a los accionistas?
    
    Una alternativa consistiría en implantar la integración del impuesto sobre sociedades (corporate tax integration), eliminando el impuesto sobre sociedades y gravando la renta empresarial directamente en el ámbito de las personas físicas.
    
    Un enfoque habitual de integración consistiría en atribuir cada año los beneficios de la sociedad a sus accionistas, con independencia de que dichos beneficios se distribuyeran o no en forma de dividendos.
    
    Una persona que poseyera el 5 % de una empresa vería imputada, en el momento de presentar su declaración de impuestos, una renta equivalente al 5 % de los beneficios de la empresa, mientras que la propia sociedad no pagaría ningún impuesto.
    
    Un sistema de este tipo eliminaría cualquier trato fiscal preferente asociado a una forma concreta de distribuir los beneficios empresariales a los accionistas, ya que dichos beneficios tributarían directamente en el ámbito del accionista, con independencia de la forma en que esos fondos le fueran entregados.
    
    Además, este enfoque gravaría a las sociedades exactamente igual que se grava a las entidades no societarias, como las sociedades personalistas (partnerships).
    
    Asimismo, eliminaría cualquier sesgo que el impuesto sobre sociedades pudiera introducir en la decisión de adoptar la forma jurídica de sociedad o la de entidad no societaria.
    
    Óscargunas investigaciones sugieren que las empresas deciden abandonar el sector societario cuando aumenta la carga tributaria de las sociedades en relación con la soportada por las entidades no societarias.²⁷
    
    Esta elevada elasticidad sugiere que puede existir una pérdida irrecuperable de eficiencia (deadweight loss) considerable cuando las empresas adoptan una forma jurídica inadecuada únicamente por motivos fiscales.
    
    La integración del impuesto sobre sociedades reduciría, por tanto, aún más esa pérdida irrecuperable de eficiencia derivada de la tributación societaria.
    
    No obstante, una reforma de este tipo reduciría los ingresos del gobierno federal, ya que toda la renta empresarial tributaría una sola vez, en lugar de hacerlo dos veces, como sucede actualmente con la inversión financiada mediante capital propio.
    
    El ahorro fiscal derivado de la integración beneficiaría principalmente a los inversores con mayores niveles de renta, que son quienes poseen la mayor parte de las acciones de las sociedades.
    
    Por ello, como ocurre con muchas otras reformas del impuesto sobre la renta de las sociedades, las ganancias de eficiencia derivadas de la integración deben ponderarse frente a la reducción de la equidad vertical del sistema tributario.
\end{specialblock}

\subsection{Aproximación macroeconómica: Incidencia sobre la función tributaria}
El interés por el lado de la oferta en macroeconomía es una de las principales características del cambio de paradigma económico experimentado en los años 70's y 80's. Dentro de ese amplio movimiento científico, se enmarcan los Economistas de la Oferta, que centraron su atención en la política fiscal y determinaron el diseño de las políticas públicas en la década de 1980 (especialmente, a través de las conocidas \textit{reaganomics}).

La mayoría de los economistas consideraban hasta entonces que las reducciones de impuestos afectaban a la producción a través únicamente de su impacto sobre la demanda agregada, en línea con los postulados keynesianos que defendían las políticas de demanda; ignorando, conscientemente, los posibles efectos de los impuestos en el lado de la oferta. 

Sin embargo, en la década de 1 970, a medida que la inflación supuso un aumento nunca antes experimentado de la progresividad en frío, algunos economistas estadounidenses (ROBERTS, TURE, LAFFER y MUNDELL) argumentaron que los impuestos altos afectaban negativamente a los ingresos tributarios, la producción y la eficiencia en el uso de recursos. Como concpto central para apuntalar sus desarrollos utilizaron la tasa impositiva marginal, que tiene efectos (des)incentivadores sobre el deseo de obtener mayor renta. 

Desde su punto de vista, un aumento en las tasas impositivas marginales afectaría negativamente la producción de una economía de dos maneras:
\begin{la}
    \item \textbf{Incentivo: Agentes}. En primer lugar, tasas marginales altas y crecientes reducen la recompensa que las personas obtienen por el trabajo y por otras actividades productivas gravables. Por tanto, cuando aumentan se reduciría tanto la oferta de trabajo como las inversiones individuales que incorporen riesgo. Por ejemplo, personas con cónyuges que trabajan optarán por salir de la fuerza laboral, disponer de más vacaciones, jubilarse anticipadamente o renunciar a oportunidades de horas extras. Además, en un marco abierto se generarían movimientos transfronterizos de factores de produccción, en especial de aquellos altamente productivos (como el trabajo cualificado). Todos estos ajustes en la oferta de trabajo darían lugar a una caída de la producción;
    \item \textbf{Recaudación fiscal}. En segundo lugar, tasas impositivas marginales altas afectarían negativamente el volumen de recaudación fiscal. De manera directa, favorecerían la aparición de estrategias para eludir impuestos y otras formas de evasión fiscal, lo que resulta ineficiente. De manera indirecta, un aumento en las tasas impositivas marginales, reduciría la base impositiva por sus efectos sobre las rentas laborales, la inversión, la inmigració y el crecimiento económico, conduciendo a su vez a un aumento en los ingresos fiscales menor que proporcional. 
\end{la}

Adicionalmente, suponiendo que cualquier política resulta distorsionante, los Economistas de la Oferta defendieron al menos tres importantes proposiciones con importantes implicaciones de política económica: (i) la protección de consumidores y trabajadores es costosa y no se justifica desde el punto de vista de su relación coste-beneficio; (ii) los programas de protección social desincentivan el esfuerzo laboral; y, (iii) el impuesto sobre la renta está intrínsecamente sesgado contra el esfuerzo, el ahorro y la inversión, de forma que reducir los tipos marginales del impuesto sobre la renta provocará un aumento de la oferta de trabajo, el ahorro y la inversión. 

Veamos algunos de los desarrollos analíticos más importantes.

\subsubsection{Recaudación fiscal: Presión fiscal e incentivos}
En tanto la imposición condiciona la conducta de los agentes, resulta esencial estudiar la relación entre los impuestos y la recaudación tributaria. 


LAFFER (1979), no fue el primero, pero sí el que introdujo de manera más sonada (y un tanto absurda), esta idea en el debate político:\footnote{%
    Si bien es cierto que se tradicionalmente no se cuestionaba la existencia de una relación positiva entre tipo impositivo y recaudación. Esta idea se basaba en el supuesto de que los impuestos distorsionaban moderadamente los incentivos de los agentes, y, por tanto, al nivel de renta y a la base imponible. Existen matices a esta alirnrnciún. Por ejemplo, SMITH defiende que impuestos altos pueden conducir a una menor recaudación fiscal que impuestos más moderados, bien sea por la reducción en el consumo de bienes y servicios, bien sea por incentivar el contrabando y la creación de mercados informales.
} reducciones del tipo impositivo no conducen necesariamente a pérdidas de recaudación. 

La idea es simple. Con tipos suficientemente elevados, se introducen distorsiones de tal magnitud sobre las ofertas factoriales (que a su vez determinan la renta y la base imponible) que la recaudación disminuye. Gráficamente la función de recaudación en función del tipo impositivo tendrá la siguiente forma:



Cuando el tipo impositivo es cero o $100\%$, el impuesto no genera recaudación. A medida que el tipo impositivo aumenta desde cero, la recaudación aumenta, pero una vez alcanza el máximo con el tipo impositivo $t^*$, los ingresos disminuyen y vuelven a caer a cero con un tipo impositivo del $100\%$. Por tanto, el objetivo del diseño del sistema fiscal es identificar un esquema de tipos impositivos que maximice el bienestar social, y que al mismo tiempo considere que el aumento de los tipos impositivos tiene efectos contradictorios sobre los ingresos. Que la economía se encuentre o no en el tipo impositivo óptimo, no obstante, es difícil de sabe. 

Desde un enfoque analítico, el objetivo a la hora de fijar el impuesto óptimo sobre la renta es elegir el patrón de tipos impositivos entre los individuos que maximice el bienestar social, sujeto a un objetivo de recaudación. ¿Hasta qué punto debería ser progresivo el impuesto sobre la renta? Como demostró el debate en torno a la propuesta del presidente Obama, difícilmente existe una cuestión más controvertida en las finanzas públicas. El economista del siglo XIX John McCulloch, que se oponía a la tributación progresiva, argumentó que una vez que se abandona la tributación proporcional, «se está en el mar sin timón ni brújula, y no hay cantidad de injusticia y necedad que no se pueda cometer». El objetivo de la teoría de la tributación óptima de la renta es proporcionar un timón, es decir, proporcionar una forma sistemática de pensar sobre la compensación «correcta» entre equidad y eficiencia.

\paragraph{Modelo EDGEWORTH: Función de bienestar social}
La primera aproximación analítica relativa a la construcción de un sistema de recaudación óptima basada en la imposición sobre la renta fue propuesto a finales del siglo XIX, por EDGEWORTH (1959/1897). Éste examinó la cuestión utilizando un modelo sencillo basado en los siguientes supuestos:
\begin{lnum}
    \item \textbf{Optimización: FBS y euilibrio presupuestario}. Sujeto a los ingresos requeridos, el objetivo es hacer que la suma de las utilidades de los individuos sea lo más alta posible. Óscargebraicamente, si $U_i$ es la utilidad del individuo $i$-ésimo y $W$ es el bienestar social, el sistema tributario debería maximizar: $W = U_1 + U_2 + \dots+ U_n$, donde $n$ es el número de personas de la sociedad;
    \item \textbf{Funciones de Utilidad: Utilidad marginal decreciente en renta}. Los individuos tienen funciones de utilidad idénticas que dependen únicamente de sus rentas. Estas funciones de utilidad presentan una utilidad marginal decreciente de la renta; a medida que aumenta la renta, un individuo mejora su situación, pero a una tasa decreciente;
    \item La cantidad total de renta disponible es fija.
\end{lnum}

Dando forma a las funciones y resolviendo el problema llegamos a un resultado bastante obvio: implica una estructura tributaria radicalmente progresiva. En el modelo de EDGEWORTH, las rentas se nivelan desde arriba hasta que se alcanza la igualdad completa. En efecto, los tipos impositivos marginales sobre los individuos de renta alta son del $100\%$. Esto es así porque, con estos supuestos, la maximización del bienestar social requiere que la utilidad marginal de la renta de cada persona sea la misma. Cuando las funciones de utilidad son idénticas, las utilidades marginales son iguales únicamente si las rentas son iguales. Las implicaciones para la política tributaria son claras: los impuestos deberían fijarse de modo que la distribución de la renta después de impuestos sea tan igual como sea posible. En particular, la renta debería tomarse primero de los ricos porque la utilidad marginal perdida es menor que la de los pobres. Si el gobierno requiere más ingresos, incluso después de alcanzar la igualdad completa, la carga tributaria adicional debería distribuirse uniformemente.

\textcolor{red}{Este resultado supone un equilibrio entre: (i) equidad vertical, el bienestar social se maximiza cuando aquellos que tienen un elevado nivel de consumo, y por tanto una baja utilidad marginal, suportan una mayor carga tributaria mientras que aquellos que tienen un nivel inferior de consumo, y por tanto una utilidad marginal mayor, tienen una menor carga impositiva; y, (ii) el comportamiento de los agentes en respuesta a los impuestos, ya que el aumento de la carga impositiva de un grupo de individuos puede dar lugar a una reducción de su oferta laboral, lo que reduce la base imponible del impuesto y la recaudación tributaria. }

Sin embargo, cada uno de los supuestos que subyacen a este análisis es cuestionable. En las últimas décadas, los economistas han investigado cómo cambian los resultados de Edgeworth cuando se relajan determinados supuestos.

\textcolor{red}{Los supuestos de Edgeworth son prácticamente idénticos a los supuestos que subyacen al modelo de distribución óptima de la renta presentado en el Capítulo 12 bajo «Justificaciones para la redistribución de la renta». Óscarlí mostramos que, con estos supuestos, la maximización del bienestar social requiere que la utilidad marginal de la renta de cada persona sea la misma. Cuando las funciones de utilidad son idénticas, las utilidades marginales son iguales únicamente si las rentas son iguales. Las implicaciones para la política tributaria son claras: los impuestos deberían fijarse de modo que la distribución de la renta después de impuestos sea tan igual como sea posible. En particular, la renta debería tomarse primero de los ricos porque la utilidad marginal perdida es menor que la de los pobres. Si el gobierno requiere más ingresos incluso después de alcanzar la igualdad completa, la carga tributaria adicional debería distribuirse uniformemente.}
    
\paragraph{Efecto incentivo al trabajo: Endogeneidad de la renta (STERN, 1987)}
Uno de los problemas más desconcertantes del análisis de Edgeworth es el supuesto de que la cantidad total de renta disponible para la sociedad es fija: los tipos impositivos confiscatorios no tendrían ningún efecto sobre la cantidad de producción generada. 

De manera más realista, supóngase que las utilidades de los individuos dependen no solo de la renta sino también del ocio. Entonces los impuestos sobre la renta distorsionan las decisiones de trabajo y crean excesos de gravamen [\authorfont{Ver Tema 4.B.9}]. Una sociedad con una función de bienestar social aditiva se enfrenta así a un dilema ineludible. Por una parte, desea asignar la carga tributaria para igualar la distribución de la renta después de impuestos. Sin embargo, en el proceso de hacerlo, reduce la cantidad total de renta real disponible. Un sistema óptimo de imposición sobre la renta, que maximiza el bienestar social, debe tener en cuenta los costes (en exceso de gravamen) de lograr una mayor igualdad. En el modelo de EDGEWORTH, el coste de obtener una mayor igualdad es cero, lo que explica la prescripción de un resultado perfectamente igualitario. ¿Cómo cambia el resultado de Edgeworth cuando se tienen en cuenta los incentivos al trabajo? 

STERN (1987) estudió un modelo similar al de EDGEWORTH, excepto en que los individuos eligen entre renta y ocio. Para simplificar el análisis, STERN supuso que los ingresos tributarios recaudados de una persona vienen dados por:

\eqblock{
\begin{aligned}
    \text{Ingresos} = −\alpha + t × \text{Renta}
\end{aligned}
}{Efecto incentivo al trabajo: Endogeneidad de la renta (STERN, 1987). Dimensión conductual: Aproximación macroeconómica}
    
Donde $\alpha$ y $t$ son números positivos, supóngase que $\alpha = 3.000€$ y $t = 0,25$. Entonces una persona con una renta de $20.000€$ tendría una obligación tributaria de $2.000€$ ($=−3.000+ 0,25 \cdot 20.000 $). Una persona con una renta de $6.000€$ tendría una obligación tributaria de menos $1.500€$ ($= −3.000+ 0,25 \cdot 6.000 $), y recibiría una transferencia de $1.500€$ del gobierno. Representando gráficamente la ecuación de STERN en un diagrama con la renta medida en el eje horizontal y los ingresos tributarios en el vertical obtenemos:
    
\imagenfit{Fig4.png}{Esquema de tributación de la renta lineal (STERN, 1987)}{ROSEN y GRABER ()}
    
Cuando la renta es cero, la carga tributaria es negativa: el individuo recibe una transferencia del gobierno de $\alpha€$. Después, por cada euro de renta, el individuo debe pagar $t€$ al gobierno. Así, $t$ es el tipo impositivo marginal, la proporción de un euro adicional que debe pagarse en impuestos. 

En las discusiones populares, un esquema lineal del impuesto sobre la renta se denomina a menudo impuesto sobre la renta de tipo único. Obsérvese que, aunque el tipo impositivo marginal para un esquema tributario lineal es constante, el esquema es progresivo en el sentido de que cuanto mayor es la renta de un individuo, mayor es la proporción de la renta pagada en impuestos. Exactamente cuán progresivo sea depende de los valores precisos de $\alpha$ y $t$. Valores mayores de $t$ están asociados con sistemas tributarios más progresivos. Sin embargo, al mismo tiempo que valores altos de $t$ conducen a una mayor progresividad, crean mayores excesos de gravamen. El problema de la tributación óptima de la renta consiste en encontrar la mejor combinación de $\alpha$ y $t$: los valores que maximizan el bienestar social sujetos a la restricción de que se recaude una cantidad determinada de ingresos (por encima de las transferencias requeridas).
    
STERN (1987) encuentra que, permitiendo una cantidad modesta de sustitución entre ocio y renta, y con unos ingresos públicos requeridos iguales a alrededor del $20\%$ de la renta, un valor de $t$ de alrededor del $19\%$ maximiza el bienestar social. Esto es considerablemente menor que el valor del $100\%$ implícito en el análisis de EDGEWORTH. Incluso unos efectos sobre los incentivos bastante modestos parecen tener implicaciones importantes para los tipos impositivos marginales óptimos. Incidentalmente, el tipo calculado por STERN también es mucho menor que los tipos impositivos marginales reales encontrados en muchos países occidentales. Por ejemplo, bajo el IRPF actual, el tipo impositivo marginal legal más alto sobre la renta en 2025 era del $51\%$.
    
De manera más general, STERN mostró que cuanto más elástica sea la oferta de trabajo, menor será el valor óptimo de $t$, ceteris paribus. Intuitivamente, el coste de la redistribución es el exceso de gravamen que crea. Cuanto más elástica sea la oferta de trabajo, mayor será el exceso de gravamen derivado de gravarla. Por tanto, una oferta de trabajo más elástica significa un coste mayor de la redistribución, de modo que debería llevarse a cabo una menor cantidad de ella.\footnote{%
    STERN también investigó cómo afectan a los resultados funciones alternativas de bienestar social, centrándose en el impacto de otorgar diferentes ponderaciones sociales a las utilidades de los ricos y los pobres. Las Funciones de Bienestar Social con preferencias más igualitarias están representadas asignando a las utilidades de las personas pobres ponderaciones mayores que a las utilidades de los ricos. Un caso extremo interesante es el criterio maximin, según el cual el único individuo que recibe alguna ponderación en la función de bienestar social es la persona con la utilidad mínima. STERN encontró que el criterio maximin exige un tipo impositivo marginal de alrededor del $80\%$. No sorprende que, si la sociedad tiene objetivos extremadamente igualitarios, se requieran tipos impositivos altos. Incluso aquí, sin embargo, los tipos quedan por debajo del $100\%$.
}\textsuperscript{,}\footnote{%
    Una limitación del análisis de STERN es que restringe el sistema del impuesto sobre la renta a tener únicamente un tipo impositivo marginal. GRUBER y SAEZ (2002) investigaron un modelo más general que permitía cuatro tipos impositivos marginales. Su hallazgo más interesante es que las personas situadas en tramos de renta más altos deberían enfrentarse a un tipo impositivo marginal menor que las personas situadas en los tramos inferiores. La intuición detrás del resultado es que, al reducir el tipo impositivo marginal sobre las personas de renta alta, se las induce a ofrecer más trabajo, y el aumento de los ingresos tributarios puede utilizarse para reducir las cargas tributarias sobre los individuos de renta baja. Es importante señalar que, aunque los tipos impositivos marginales disminuyen con la renta, los tipos impositivos medios aumentan con la renta, por lo que el sistema tributario óptimo sigue siendo progresivo. Recientemente, un cantón suizo implementó de hecho un sistema tributario que impone tipos impositivos marginales más bajos a quienes obtienen mayores ingresos (RABUSHKA, 2003).
}

Esta catalogación de resultados puede transmitir una sensación algo falsa de precisión sobre lo que los economistas saben realmente acerca del sistema tributario óptimo. Después de todo, hay muchos juicios de valor controvertidos detrás del bienestar social aditivo que el sistema tributario óptimo busca maximizar. Además, existe una incertidumbre considerable acerca de las elasticidades de comportamiento que son cruciales para analizar la disyuntiva entre eficiencia y equidad. No obstante, calcular los tipos impositivos óptimos bajo conjuntos alternativos de supuestos es extremadamente informativo. 


\subsubsection{Recaudación fiscal: Incertidumbre y evasión fiscal}

\href{https://www.revistasice.com/index.php/ICE/article/view/7956/8213}{\textit{LA UTILIZACIÓN DEL ENFOQUE CONDUCTUAL EN LAS ADMINISTRACIONES TRIBUTARIAS}. Revistas ICE (2025)}

\href{https://www.funcas.es/wp-content/uploads/Migracion/Articulos/FUNCAS_PEE/154art07.pdf}{\textit{La economía del comportamiento en el análisis del cumplimiento fiscal}. PRIEGO, 2017. Revista FUNCAS}


En la práctica, es díficil probar que exista una relación directa entre el establecimiemo de impuestos y la recaudación tributaria. El sistema impositivo se ve afectado por problemas de información asimétrica de forma que los agentes pueden tener incentivos a llevar a cabo comportamientos estratégicos para alcanzar sus propios intereses, lo cual puede tener un impacto negativo en la recaudación fiscal.

El que más atención recibe, de entre los comportamientos estratégicos posibles, es la evasión fiscal. La evasión fiscal consiste en no pagar los impuestos que legalmente se adeudan.\footnote{%
    No es un problema nuevo. Hace siglos, Platón observó: \textit{Cuando existe un impuesto sobre la renta, el hombre justo pagará más y el injusto menos por la misma cantidad de ingresos.}
}

El fraude fiscal es extremadamente difícil de medir. Por ejemplo, para Estados Unidos, el Internal Revenue Service (IRS) estima que los contribuyentes pagan voluntariamente solo alrededor del $80\%$ de la obligación tributaria real correspondiente al impuesto sobre la renta. Si esta estimación es siquiera aproximadamente correcta, sugiere que la evasión constituye un problema de enorme importancia.

Las personas cometen fraude fiscal de diversas maneras:
\begin{la}
    \item Llevar dos contabilidades para registrar las operaciones del negocio. Una refleja la actividad real y la otra es la que se presenta a las autoridades tributarias. Óscargunos evasores incluso utilizan dos cajas registradoras;
    \item Trabajar en empleos secundarios cobrando en efectivo. Evidentemente, tener un segundo empleo es completamente legal. Sin embargo, los ingresos obtenidos en estos trabajos suelen pagarse en efectivo, en lugar de mediante cheque. Como consecuencia, no existe un registro oficial de dichos ingresos y estos no se declaran a las autoridades fiscales;
    \item Declarar ingresos inferiores a los realmente obtenidos. La omisión de ingresos constituye una forma muy común de evasión fiscal. Según la Government Accountability Office, la infradeclaración de ingresos es especialmente frecuente entre las personas que trabajan por cuenta propia (Herman, 2007). En 2009, el secretario del Tesoro de Estados Unidos, Timothy Geithner —quien supervisa el IRS—, tuvo problemas durante su audiencia de confirmación por no haber pagado impuestos sobre parte de sus ingresos como trabajador autónomo; y/o
    \item Realizar operaciones en efectivo. Pagar bienes y servicios en efectivo y emitir cheques al portador («cash») dificulta enormemente que el IRS pueda rastrear las transacciones.
\end{la}

En otro tiempo, la evasión fiscal se asociaba con millonarios que ocultaban su patrimonio en cuentas bancarias suizas. No obstante, la imagen actual del evasor fiscal probablemente sea la de un reparador cuyos ingresos proceden de trabajos informales, no declarados a efectos fiscales, o la de unos padres que evaden impuestos sobre los salarios pagados a una niñera.\footnote{%
    De hecho, en Estados Unidos, quienes pagan a empleadas domésticas, niñeras u otros trabajadores del hogar más de aproximadamente $1.500€$ al año están obligados a ingresar las cotizaciones a la Seguridad Social correspondientes. Sin embargo, menos del $0,25\%$ de los hogares pagan este denominado \textit{nanny tax} (HERMAN, 2004b).
} Evidentemente, la sensación de que «todo el mundo lo hace» está muy extendida.

Para abordar esta importante cuestión, en primer lugar analizaremos la teoría positiva de la evasión fiscal, para después abordar la cuestión normativa de cómo debería enfrentarse la política pública a este fenómeno.

\paragraph{Análisis positivo: }

Supongamos que Óscar solo se preocupa por maximizar su renta esperada. Dispone de una determinada cantidad de ingresos y trata de elegir $R$, la cantidad que ocultará a las autoridades tributarias. Supongamos también que su tipo marginal del impuesto sobre la renta es $0,3$ 

Por cada euro que consigue excluir de la renta imponible, su factura fiscal disminuye en treinta céntimos. Este es el beneficio marginal que obtiene al ocultar un dólar de ingresos a las autoridades tributarias. Más generalmente, cuando Óscar se enfrenta a un tipo marginal del impuesto sobre la renta igual a $t$, el beneficio marginal de ocultar cada dólar es precisamente $t$. La administración tributaria desconoce cuál es la renta real de Óscar, pero realiza auditorías aleatorias sobre las declaraciones de los contribuyentes. Como consecuencia, existe una cierta probabilidad, $P$, de que Óscar sea inspeccionado. Si es descubierto haciendo trampas, Óscar deberá pagar una sanción que aumenta con $R$, y además lo hace a una tasa creciente.\footnote{%
    Obsérvese que, si fuera posible controlar a Óscar gratuitamente cada segundo de cada día, las oportunidades para evadir impuestos simplemente no existirían. El hecho de que ese control permanente sea inviable constituye la fuente fundamental del problema. De allí que una cuestión esencial sea la introducción de información incompleta, o incertidumbre, en el contexto analítico.
}

Suponiendo que Óscar conoce el valor de $P$ y el régimen de sanciones, tomará su decisión comparando los costes marginales y los beneficios marginales de hacer trampas. Gráficamente, podemos representar la cantidad de renta no declarada en el eje horizontal y los euros en el eje vertical:

\imagenfit{Fig13.png}{Evasión fiscal: Óptimo positivo (incentivo)}{ROSEN y GARBER}

El beneficio marginal ($MB$) de cada euro no declarado es $t$, es decir, el importe del impuesto ahorrado. Por el contrario, el coste marginal esperado ($MC$) es igual al incremento de la sanción por cada dólar adicional ocultado (la sanción marginal) multiplicado por la probabilidad de ser descubierto. Por ejemplo, si la sanción adicional por ocultar el dólar número mil es de 1,50€ y la probabilidad de ser descubierto es de 1 entre 3, entonces la sanción marginal esperada asciende a 50 centavos. Con ello, encontramos que la cantidad óptima de evasión se encuentra en el punto donde ambas curvas se cruzan, es decir, en $R^*$. Es óptima en el sentido de que, en promedio, constituye la estrategia que maximiza la renta de Óscar. 

En un mundo caracterizado por la incertidumbre, encontrar la mejor estrategia en términos de valor esperado constituye una forma razonable de proceder. Es posible, naturalmente, que la estrategia óptima consista en no hacer trampas en absoluto. Por ejeplo, puede darse el caso  en que el coste marginal de hacer trampas supera al beneficio marginal para cualquier valor positivo de $R$, de modo que el nivel óptimo de evasión es igual a cero. Gráficamente:

\imagenfit{Fig14.png}{Evasión fiscal: Óptimo nulo (desincentivo)}{ROSEN y GARBER}

En esencia lo que nos dice el modelo es que la evasión disminuye cuando bajan los tipos marginales del impuesto. Ello se debe a que una reducción del valor de $t$ disminuye el beneficio marginal de evadir impuestos, desplazando hacia abajo la curva de beneficio marginal, de modo que su intersección con la curva de coste marginal se produce para un valor menor de $R$.\footnote{%
    Esta predicción es coherente con la evidencia anecdótica. Considérese, por ejemplo, el caso de Estonia, una de las repúblicas bálticas, que sustituyó recientemente su sistema de elevados tipos marginales progresivos por un impuesto proporcional sobre la renta del $26\%$. El antiguo primer ministro Mart Laar afirmó que esta reforma redujo drásticamente la evasión porque: \textit{En el mundo real, las personas ricas encuentran la manera de evitar los impuestos elevados. Con un impuesto proporcional dejan de preocuparse por ocultar sus ingresos o por trabajar en la economía sumergida.} (TIERNEY, 2006). Además, la predicción del modelo también se ve confirmada por estudios econométricos. Por ejemplo, FISMAN y WEI (2004) encuentran que unos aranceles elevados en China provocan una evasión fiscal considerable. Según sus estimaciones, un incremento del $1\%$ en el tipo impositivo sobre las importaciones induce a los importadores a aumentar en un $3\%$ la cantidad de impuestos que evaden.
}

Aunque este modelo proporciona ideas útiles, ignora algunas consideraciones potencialmente importantes;\footnote{%
    Por ejemplo, los costes psicológicos de hacer trampas. Además, en cualquiera de las dos representaciones suponemos que las personas solo se preocupan por la renta esperada y que el riesgo, en sí mismo, no les importa. En la medida en que los individuos sean aversos al riesgo, sus decisiones respecto a participar en lo que esencialmente constituye una apuesta pueden verse modificadas. También suponemos que la única decisión consiste en cuánto ingreso declarar; es decir, que el tipo de empleo y el nivel de ingresos antes de impuestos se consideran dados. Sin embargo, en realidad, el sistema tributario puede afectar tanto al número de horas trabajadas como a la elección de la ocupación. Por último, en nuestro análisis simplificado la probabilidad de ser inspeccionado es independiente tanto de la cantidad evadida como del volumen de renta declarada. Sin embargo, en muchos países la probabilidad de auditoría depende tanto de la ocupación del contribuyente como del nivel de renta declarado, lo que complica el modelo, aunque no modifica sus aspectos esenciales.
} es evidente que el fraude fiscal constituye un fenómeno mucho más complejo de lo que sugieren las representaciones. No obstante, el modelo proporciona un marco útil para analizar los factores que influyen en las decisiones de evasión.

\paragraph{Análisis normativo: ¿Por qué perseguir la evasión fiscal?}
La mayoría de los debates públicos sobre la economía sumergida parten del supuesto de que se trata de un fenómeno negativo y de que la política pública debe orientarse a reducir su tamaño. Aunque esta afirmación puede ser correcta, merece un examen detenido.

Una cuestión importante en este contexto es si debemos preocuparnos o no por el bienestar de quienes evaden impuestos. En la terminología de la economía del bienestar, la pregunta consiste en saber si las utilidades de quienes participan en la economía sumergida deben formar parte de la función de bienestar social. 

Supongamos, por el momento, que sí forman parte de ella. En ese caso, bajo determinadas condiciones, la existencia de una economía sumergida puede incrementar el bienestar social. Por ejemplo, si la oferta de trabajo es más elástica en la economía sumergida que en la economía formal, la teoría de la imposición óptima sugiere que la primera debería gravarse con un tipo relativamente bajo. No se trata más que de una aplicación de la regla de la elasticidad inversa (inverse elasticity rule), expresada en la Ecuación (16.9).

Alternativamente, supongamos que quienes participan en la economía sumergida tienden a ser más pobres que quienes trabajan en la economía formal. De hecho, muchos observadores consideran que la economía sumergida constituye un componente esencial de la vida en los barrios marginales de las ciudades estadounidenses. En la medida en que la sociedad persiga objetivos igualitarios de redistribución de la renta, podría ser deseable dejar intacta la economía sumergida.

Consideremos ahora las implicaciones para la política pública cuando los evasores no reciben ningún peso en la función de bienestar social y el objetivo consiste simplemente en eliminar el fraude con el menor coste administrativo posible. La misma representación anterior sugiere una forma sencilla de alcanzar ese objetivo.

\imagenfit{Fig13.png}{Evasión fiscal: Óptimo positivo (incentivo)}{ROSEN y GARBER}

El coste marginal esperado de hacer trampas es el producto entre la sanción y la probabilidad de ser descubierto. La probabilidad de detección depende de los recursos destinados a la administración tributaria. Si el auditor dispone de un presupuesto elevado, podrá detectar a un gran número de evasores. \footnote{%
    Sin embargo, incluso si las autoridades tributarias cuentan con un presupuesto reducido y, por tanto, la probabilidad de detección es baja, el coste marginal de la evasión puede hacerse arbitrariamente elevado si las sanciones son suficientemente severas. Por ejemplo, si solo se capturara a un evasor fiscal al año, pero este fuera ahorcado públicamente por su delito, el coste esperado de la evasión fiscal bastaría para disuadir a muchas personas. El hecho de que una política tan draconiana nunca haya sido seriamente propuesta en Estados Unidos indica que los sistemas sancionadores existentes tratan de incorporar el principio de la retribución justa.
}

En definitiva, contrariamente a los supuestos del marco utilitarista, la sociedad no se preocupa únicamente por el resultado final (eliminar a los evasores), sino también por los procedimientos mediante los cuales se alcanza dicho resultado.

\textcolor{red}{\textbf{OJO ->} RECAP ICEX-CECO}
Las autoridades fi scales disponen de dos métodos eficaces para disuadir la evasión fiscal. pero ambos métodos no son equivalentes. Es probable que el aumento de los esfuerzos de control y aplicación de la ley para aumentar p conlleven un mayor coste respecto al aumento de s''. No obstante, la sociedad podría no estar dispuesta a aumentar s lo suficientemente como para reducir la evasión fiscal.\footnote{%
    Asimismo. el principio jurídirn de que "la pena debe aj ustarse al delito'' (principio de proporcionalidad de la pena) podría explicar por qué las autoridades tiscale�� realizan .:ostosos esfuerzos de control e inspección en lugar de impon��r sanciones extremadamente duras.
}

En definitiva, ante el establecimiento de impuestos los agentes tienen incentivos para evadirlos, aunque éstos pueden reducirse con mecanismos de control y sanción. No obstante, parte de la literatura económica sefiala que estos modelos obtienen una tasa de cumplimiento de la norma t1ibutaria inferior a la observada en tanto ignora cuestiones como el sentimiento de deber social o la percepción en torno a la justicia del sistema tributario.

Tornando en consideración este análisis cabe seilalar la importancia del estudio contra la evasión fi scal. En primer lugar, la evasión fiscal provoca efectos disuibutivos y atenta contra la equidad en tanto da lugar a transferencias desde quienes contribuyen hacia quienes def raudan. Adicionalmente, conlleva pérdidas de eficiencia en tanto el sector público puede verse obligado a aumentar los tipos impositivos para mantener los mismos niveles de recaudación que se obtendrían si todos los individuos fueran honestos en su contribución. Y. por último. como ya se ha apuntado. la evasión fiscal conlleva importantes costes administrativos de inspección y sanción.




\subsubsection{Recaudación fiscal: Economía abierta y empresas multinacionales}
Existe, no obstante, otro comportamiento estratégico que merma (muy notablemente) la recaudación fiscal del Estado: la elusión fiscal. La elusión fiscal, a la que KEYNES se refirió en una ocasión como \textit{el único pasatiempo intelectual que proporciona alguna satisfacción} consiste en modificar el comportamiento para reducir la deuda tributaria.\footnote{%
    Los contribuyentes se aprovechan de lagunas legales para reducir su obligación tributaria. Por ejemplo modificando la obtención de ingresos en forma de ganancias del capital, como ya vimos, o de beneficios complementarios no gravados, en lugar de hacerlo en forma de sueldos y salarios totalmente gravados
} 

La elusión fiscal, que John Maynard Keynes calificó en una ocasión como «la única actividad intelectual que siempre proporciona una recompensa», consiste en modificar el propio comportamiento con el fin de reducir la carga tributaria. No hay nada ilegal en la elusión fiscal: \textit{Una y otra vez los tribunales han afirmado que no hay nada siniestro en organizar los propios asuntos de manera que los impuestos sean lo más bajos posible. Todo el mundo lo hace, ricos y pobres; y todos hacen bien, porque nadie tiene el deber público de pagar más de lo que exige la ley. }[...]\textit{ Exigir más en nombre de la moral no es más que pura hipocresía.} Juez Learned Hand, Commissioner v. Newman (1947).

En todo caso, de fonna simplificada, se puede decir que la elusión fiscal es lícita mientras que la evasión fiscal es ilícita y sus consecuencias son inciertas ya que dependen de la probabilidad de que el contribuyente sea descubierto. Por tanto. a través de la regulación el sector público puede establecer distintos mecanismos que interactúen con los agentes para así incentivar la declaración de su renta verdadera.

24.6 Tratamiento de la renta internacional de las sociedades

Las sociedades de todo el mundo operan en un mercado mundial de productos cada vez más integrado.

Como resultado, puede haber diversas razones por las que las empresas deseen producir en otros países, como unos costes laborales más bajos, menores costes de transporte al vender sus productos en esos países o una mayor capacidad para adaptar los productos a las preferencias locales.

Los costes de producción o las ventajas comerciales no son las únicas razones por las que las empresas producen bienes en el extranjero. La estructura del impuesto sobre sociedades en muchos países desarrollados puede desempeñar un papel importante al fomentar la producción internacional como medio para evitar la tributación. De hecho, algunas de las mayores empresas ven reducida de forma muy significativa su carga tributaria gracias a su diversificación internacional.

General Electric, por ejemplo, obtuvo unos beneficios mundiales de 14.200 millones de€ en 2010, de los cuales 5.100 millones procedían de Estados Unidos. Sin embargo, cuando llegó el momento de pagar impuestos, en lugar de pagar, GE solicitó una devolución fiscal de 3.200 millones de€.

En esta sección analizamos el tratamiento de la renta internacional dentro de la estructura del impuesto sobre sociedades de Estados Unidos.

(Otro tema interesante tratado en, aunque no se revisa aquí, es la sensibilidad de la inversión extranjera frente a los tipos impositivos de cada país, que se puede consultar en Gordon y Hines (2002). La evidencia disponible sugiere que las empresas son muy sensibles a los tipos impositivos al decidir dónde localizar sus inversiones entre distintos países).

Cómo gravar la renta internacional

Existen dos enfoques básicos para gravar a las sociedades que obtienen rentas en el extranjero.

(Aunque esta exposición se refiere al impuesto sobre sociedades, muchos de los conceptos introducidos aquí también son aplicables al tratamiento fiscal de las rentas obtenidas por personas físicas en el extranjero).

El primero es un sistema territorial, mediante el cual las sociedades que obtienen rentas en el extranjero pagan impuestos únicamente al gobierno del país donde se genera esa renta.

El segundo es un sistema global, mediante el cual las sociedades son gravadas por su país de residencia sobre toda su renta, con independencia del lugar donde haya sido obtenida.

Aproximadamente la mitad de los países de la OCDE, entre ellos Estados Unidos, utilizan el enfoque global.

Por tanto, las empresas estadounidenses pagan el impuesto sobre sociedades de Estados Unidos con independencia del lugar donde estén situadas sus plantas o fábricas.

No obstante, las empresas pueden computar los impuestos pagados a gobiernos extranjeros como un crédito fiscal frente a los impuestos que deben pagar en Estados Unidos.

Gracias a este crédito fiscal por impuestos extranjeros (foreign tax credit), las empresas deberían, en principio, soportar el mismo tipo impositivo (el tipo del impuesto sobre sociedades estadounidense) con independencia del país en el que localicen sus filiales.

En la práctica, sin embargo, esto no sucede, lo que genera una ventaja fiscal asociada a la obtención de rentas en otros países.

Esta ventaja fiscal tiene dos causas:

* la ganancia en valor actual derivada de retrasar el pago de impuestos sobre las rentas obtenidas en el extranjero;
* la posibilidad de trasladar beneficios desde países con impuestos elevados hacia países con impuestos reducidos.

Repatriación de dividendos procedentes del extranjero

La primera ventaja fiscal de la renta obtenida en el extranjero consiste en que las empresas solo tributan por esa renta internacional cuando esta es devuelta, o repatriada, desde la filial extranjera a la empresa matriz estadounidense.

Como consecuencia, las empresas estadounidenses nunca estarán sujetas al impuesto estadounidense mientras mantengan sus beneficios en sus filiales extranjeras.

En el caso de una filial situada en un país con tipos del impuesto sobre sociedades inferiores a los de Estados Unidos, esta estructura fiscal proporciona un poderoso incentivo para retrasar la repatriación de beneficios.

Por ejemplo, si el país extranjero aplica un tipo del 10 % sobre los beneficios empresariales, mientras que el tipo estadounidense es del 35 %, la empresa puede aplazar la tributación correspondiente al 25 % de los beneficios simplemente evitando repatriarlos desde la filial extranjera hacia Estados Unidos.

Este aplazamiento constituye una ventaja fiscal porque esos impuestos estadounidenses pendientes (el 25 %) se pagarán en el futuro, cuando los beneficios sean finalmente repatriados.

En consecuencia, el valor actual de los impuestos pagados sobre la renta extranjera es inferior al que existiría si dichos impuestos se pagaran en el momento en que la renta fue obtenida.

Si el dinero nunca se repatría, la empresa evita por completo el pago del impuesto estadounidense sobre esos beneficios obtenidos en el extranjero.

(Desai et al. (2001) comparan el comportamiento de las filiales extranjeras de empresas estadounidenses, que no pagan impuestos estadounidenses sobre sus rentas extranjeras hasta que estas se repatrian, con el de las sucursales extranjeras de empresas estadounidenses, que sí tributan cuando obtienen la renta, con independencia de si esta se repatria o no. Encuentran que las filiales situadas en países con baja tributación tienen una probabilidad significativamente menor de repatriar dividendos que las sucursales localizadas en esos mismos países, lo cual es coherente con una repatriación motivada por razones fiscales).

Esta situación es muy similar a la ventaja fiscal derivada de gravar las ganancias de capital en el momento de su realización, en lugar de hacerlo conforme se devengan: el tipo efectivo disminuye porque los impuestos se pagan en el futuro.

Precios de transferencia

La segunda razón por la que las empresas multinacionales pueden pagar menos impuestos es que, cuando un bien se produce utilizando factores procedentes de varios países, resulta difícil atribuir correctamente los beneficios generados por ese bien a un país concreto.

En particular, las empresas tendrán incentivos para declarar que los beneficios se obtienen en países con tipos bajos del impuesto sobre sociedades, mientras que declararán que los gastos se producen en países con tipos elevados, donde dichos gastos reducen la base imponible.

Una empresa puede conseguir este objetivo manipulando sus precios de transferencia (transfer prices), es decir, el importe que cualquiera de las filiales de una empresa paga a otra filial de la misma empresa por los bienes transferidos entre ellas.

Imaginemos que Francia aplica un impuesto del 50 % sobre los beneficios empresariales, frente al 35 % vigente en Estados Unidos.

Una empresa estadounidense de ordenadores posee una filial francesa que fabrica microchips con un coste de 100€ por unidad.

La empresa estadounidense transfiere esos microchips a Estados Unidos, donde gasta otros 500€ en completar el ordenador y vende el producto terminado por 1.000€.

Por tanto, el beneficio total de la operación asciende a 400€.

¿Cómo decide la empresa repartir esos 400€ de beneficio entre la matriz estadounidense y la filial francesa?

La empresa estadounidense podría afirmar que la totalidad del beneficio de 400€ procede del microchip y transferir 500€ a su filial francesa por cada chip (los 100€ de coste más 400€ de beneficio).

Sobre el papel, la filial francesa obtendría un beneficio de 400€ (los 500€ recibidos desde Estados Unidos menos los 100€ que cuesta fabricar el chip) y pagaría 200€ en impuestos (0,50 × 400€) al gobierno francés.

La empresa estadounidense no registraría ningún beneficio, ya que vende el ordenador por 1.000€ y registra como costes 500€ pagados a la filial francesa por el chip y 500€ correspondientes a la fabricación del ordenador en Estados Unidos.

En el extremo opuesto, la empresa estadounidense podría afirmar que ninguna parte del beneficio de 400€ procede del microchip y transferir únicamente 100€ a su filial francesa por cada chip.

En ese caso, la filial francesa no obtendría ningún beneficio, mientras que la empresa estadounidense obtendría 400€ de beneficio por ordenador (1.000 − 500 − 100 = 400€).

En consecuencia, la empresa pagaría únicamente 140€ en impuestos (0,35 × 400), esta vez al gobierno de Estados Unidos.

Así, la empresa estadounidense puede reducir su carga tributaria en 60€ por ordenador transfiriendo únicamente 100€ a la filial francesa.

El problema de los precios de transferencia derivado de esta práctica consiste en que las empresas tienen incentivos para hacer que los beneficios parezcan haberse obtenido en el país con menor tributación.

Estados Unidos y los demás países de la OCDE son conscientes de este problema.

Por ello, las autoridades tributarias de esos países exigen que, en las operaciones entre una empresa y sus filiales extranjeras, los precios de transferencia se registren como si la operación hubiera tenido lugar en condiciones de plena competencia (arm’s length), es decir, como si dos empresas independientes y sin relación entre sí hubieran negociado el precio.

En la práctica, esta regla resulta extraordinariamente difícil tanto de aplicar como, en muchos casos, de interpretar.

Existe abundante evidencia de que las empresas manipulan los precios y los beneficios entre sus distintas filiales con el fin de reducir su carga fiscal.

Por ejemplo, Collins et al. (1998) encontraron que las multinacionales estadounidenses declaran una rentabilidad extranjera mayor cuando los tipos impositivos en el extranjero son inferiores a los de Estados Unidos.

Quizá el ejemplo más importante de este tipo de práctica fue el caso de GlaxoSmithKline, una empresa farmacéutica con sede en el Reino Unido que comercializaba el medicamento contra la úlcera Zantac a través de filiales situadas en numerosos países.³⁸

Glaxo tenía una filial en Estados Unidos, que en este caso era el país con mayor tributación.

Por ello, podía reducir sus impuestos trasladando gastos hacia Estados Unidos y beneficios hacia el Reino Unido.

Según el Internal Revenue Service (IRS) estadounidense, Glaxo hizo precisamente eso al pagar unos royalties excesivamente elevados por el medicamento Zantac a su empresa matriz británica, que era quien había realizado la investigación y el desarrollo del medicamento.

Esos elevados pagos por derechos de propiedad intelectual dieron lugar a un mayor pago de impuestos en el Reino Unido y a un menor pago de impuestos en Estados Unidos, ya que la filial estadounidense podía deducir esos royalties como un gasto de explotación.

Como consecuencia, en enero de 2004, el IRS reclamó a la empresa el pago de 5.200 millones de€: 2.700 millones en impuestos atrasados correspondientes al período 1989-1996 y 2.500 millones en intereses.

La disputa sobre los impuestos de Glaxo se prolongó durante años ante el tribunal tributario estadounidense y estaba previsto que llegara a juicio en 2007.

Sin embargo, en septiembre de 2006, Glaxo anunció que resolvería el conflicto con el gobierno estadounidense mediante un acuerdo por importe de 3.100 millones de€, que cubría los ejercicios comprendidos entre 1989 y 2005.

(Hilzenrath (2006). Aunque la mayor parte del debate sobre tributación internacional se centra en la comparación entre los sistemas territorial y global, Auerbach (2010) propuso una tercera alternativa: la tributación basada en el destino (destination-based taxation), mediante la cual las empresas estadounidenses tributarían únicamente por sus ventas realizadas en Estados Unidos y solo podrían deducir los gastos efectuados en ese país. Este enfoque tendría el inconveniente de no gravar los beneficios obtenidos por las multinacionales en otros países, pero tendría la ventaja de eliminar la posibilidad de trasladar a paraísos fiscales extranjeros los beneficios obtenidos en el mercado nacional).




\subsubsection{Nuevos desarrollos}

\paragraph{Imposición negativa sobre el trabajo}

Típicamente. en una economía perfectamente compet1t1va. sin ninguna intervención gubernamental o cualquier otra distorsión. el libre juego del mercado da lugar a agentes con bajos ingresos. Hasta ahora, hemos analizado cómo los impuestos incidían en las decisiones de oferta de los individuos. En este apai1ado, nos centraremos en una serie de enfoques alternativos que buscan estudiar cómo los subsidios o transferencias monetarias afectan a las decisiones de oferta en el mercado de trabajo y hasta qué punto los programas de bienestar pueden reducir la participación en el mercado laboral al fomentar la dependencia del Estado. Para ello nos centraremos en un primer momento en aquellos agentes de bajos ingresos que trabajan mediante el estudio de los subsidios salariales y, posteriormente, analizaremos en un recuadro el caso de los programas de asistencia social para aquellos individuos sin empleo.

El impuesto negativo sobre la renta puede definirse como un beneficio pagado por la administración fiscal al contribuyente. calculado sobre la base de los ingresos de dicho contribuyente. Por tanto. es una transferencia que adopta la forma de un crédito fiscal reembolsable. a partir del cual cualquier otra renta obtenida es imponible. El carácter reembolsable del crédito fiscal es lo que pen nite el pago de un impuesto negativo sobre la renta a las personas sin ingresos o con ingresos bajos.

En las últimas décadas, los subsidios salariales han ganado popularidad en las economías avanzadas. Un primer análisis realizado por lmmervoll y Pearson (2009) reveló que. en 2009. 16 de los 30 países de la OCDE habían puesto en marcha programas de subsidios salariales, y en algunos países estos rrogramas se han ampliado hasta convertirse en un componente básico de los sistemas de protección social. En particular. los subsidios salariales para los trabajadores de bajos ingresos suelen considerarse un enfoque atractivo para abordar una doble preocupación: reforzar los incentivos de participación en el mercado laboral y reducir la pobreza. 

En adelante. consideramos la posibilidad de un subsidio a los ingresos de las familias de renta baja. En particular, estudiaremos el caso del '·Earned lncome Ta>. Credit'º (EITC) de Estados Unidos. Se trata de una transferencia monetaria para los trabajadores de bajos ingresos que adopta la forma de un crédito fiscal. reduciendo la deuda tributaria del individuo. Además. la cuantía del subsidio depende del número de hijos que integren la familia y se suprime gradualmente cuando las rentas superan un umbral de ingresos.

Así. el EITC presenta una serie de diferencias con las prestaciones sociales tradicionales: i ) el EITC se proporciona a través del sistema fiscal y no del sistema de protección social; ii) el EITC está disponible para todas las familias de bajos ingresos con hijos. independientemente de su estado civil; iii) la recepción del crédito fiscal requiere que los ingresos de la familia sean positivos y, por tanto, el EITC fomenta la participación en el mercado laboral.

Cabe preguntarse si estos subsidios mejoran los incentivos a u·abajar de las personas con rentas bajas o si únicamente cumplen con un objeto de redistribución de recursos ( incremento de equidad vertical). Tomando como punto de partida la estructura del EITC en el ai'io 2008 para una familia con dos hijos. el crédito fiscal asciende al 40\% de los ingresos salariales hasta un máximo de 1 2.570€. de fomia que el crédito máximo asciende a 5.0:28€. Hasta los 16.400€ este crédito se mantiene constante. pero se reduce gradualmente cuando los ingresos del individuo superan esta cuantía hasta los 40.295€, cuando el crédito se agota. En este tramo decreciente, por cada dólar de ingresos que se obtenga, el crédito se reduce en 2 1 centavos. En consecuencia, en el intervalo creciente, el Estado añade 40 centavos a cada dólar de ingresos. lo cual equivale a un tipo impositivo marginal negativo del 40\%.\footnote{%
    El tipo impositivo es marginal porque es el que se aplica a un dúlar adicional de ingresos.
} Sin embargo, el hecho de que se rcdu7ca el crédito supone implícitamente la aplicación de un tipo marginal positivo del 21 % en el intervalo decreciente (gráfico I O).

Como hemos analizado anteriormente en relación al impuesto sobre la renta, el EITC produce tanto un efecto renta como un efecto sustitución en las decisiones de oferta de trabajo. Para analizar esta cuestión, incluimos el programa de subsidios salariales Erre al diagrama previamente analizado de consumo y ocio ( gráfico 1 1 ). La línea continua representa la restricción presupuestaria original antes de introducir el EITC, suponiendo una tasa salarial de 20€/hora; mientras que la línea discontinua representa la nueva restricción presupuestaria después de ai'iadir el Erre ( hasta 40.295€, momento en el que se vuelve a la restricción presupuestaria original).

Este gráfico nos permite ilustrar el distinto impacto del EITC según los colectivos de renta:

1. Para las personas que no forman parte de la población activa, como la persona A. es probable que el EITC aumente su oferta de trabajo. ya que la única forma de obtener el crédito es percibir rentas de trabajo. Por tanto, el EITC 110 tiene ningún efecto sobre los ingresos de este grupo, ya que comienzan sin ingresos laborales. Sin embargo, el EITC tiene un efecto de sustitución al fomentar la participación laboral, pasando al punto A1 .

2. Las personas que ya forman parte de la población activa y que ganan menos de 12.570€. como la persona B (con ingresos de 5.000€), se encuentran en la parte de creciente de la estructura del EITC, lo cual equivale a un tipo marginal negativo. Los efectos sobre la oferta de trabajo de este grupo son ambiguos: se produce un efecto de sustitución que conducen a un aumento de la oferta laboral, ya que cada hora de tralx��io aporta un salario más alto, pero al mismo tiempo se genera un efecto de renta que conducen a una disminución de la oferta laboral, ya que los tralx��jadores son más ricos como resultado de este subsidio. Por tanto, si domina el efecto de sustitución, los trabajadores se desplazarán a un punto como el B , , con una disminución del ocio y un aumento de la oferta de trabajo; por el contrario, si domina el efecto renta, los trabajadores se moverán al punto B1, con un aumento del ocio y una disminución de la oferta de trabajo.

3. Aquellos trabajadores con ingresos entre 1 2.570 y 1 6.400€, como la persona C (que gana 1 4.000€). se encuentran en la parte plana de la estructura del ElTC, recibiendo la misma cantidad de EITC independientemente de lo poco o mucho que trabajen dentro de este rango de rentas. Esta configuración de las prestaciones equivale a un desplazamiento paralelo hacia el exterior de la restricción presupuestaria, que reducirá la oferta de trabajo. Dado que el EITC no aumenta su sala1io por horas adicionales de trabajo. no hay efecto de sustitución. Sin embargo. al ser más ricos por el subsidio obtenido, el efecto renta puede reducir el número de horas trabaja das, lo que desplazaría al individuo al punto C1.
4. Aquellos trnbaja dores con ingresos entre 1 6.400 y 40.29.5 dóla res. como la persona D (que gana 20.000 dóla res). se encuentran en la parte decreciente del E rre, donde el impo11c del subsidio disminuye a medida que trabajan más. Esto hace que la pendiente de la restricción presupuestaria se reduzca en relación con los segmentos a nteriores. Así, el efecto sustitución y el efecto renta reducen la oferta de trabajo. moviéndose el individuo a un punto D1. 

En resumen. el efecto de la introducción de subsidios sa lariales para colectivos de renta s bajas sobre la oferta de trabajo resulta ambiguo. Esta ambigüedad ha dado lugar a numerosos trabajos de evidencia empírica que pasamos a analizar.

La investigación reciente es prácticamente unánime a l señalar que el EITC ha ii1crementado las tasas de pa rticipación de la fuerza de trabajo. En particular. los estudios en este ámbito se centran en el efecto sobre las madres solteras. Por ejemplo, Meyer y Rosenbaum (200 1 ) a na lizaron los datos del comporta miento laboral de las madres solteras entre 1 984 y 1 996, un periodo en el que se produjeron importantes expansiones del EJTC. Durante dicho periodo, la tasa de empleo de las madres solteras se incrementó desde el 58,5 al 64,5%,. En concreto, las estima ciones de Meyer y Rosenbaum indican que el 60% de este aumento se debió a los subsidios salariales. Por otra p arte, Eissa y Leibman ( 1 996) estimaron que las ma dres solteras tenían entre 1.4 a 3,7 puntos porcentuales más de probabilidades de trabajar corno resultado de la a mpliación del EITC en 1 986. Aunque el EITC parece tener efectos positivos sobre la oferta de trabajo de la s madres solteras, podría tener efectos nega tivos en las decisiones de oferta de trabajo de una pareja c asada , en tanto el EITC se calcularía sobre los ingresos de la unidad familiar. Eissa y Hoynes ( 1 998) estudiaron el impacto del Erre en las parejas casadas y descubrieron que no había ningún efecto sobre el trabajo de los hombres. Este resultado es coherente con los resultados de las ma dres solteras (el EITC aumentó la pai1icipación en la fuerza de tra ba,jo de las madres solteras, pero no a fectó a sus horas una vez que esta ban trabajando) porque casi todos los hombres ya participaban en el mercado laboral. Óscar mismo tiempo, los investigadores descubrieron una modesta reducción de la oferta de trabajo de las mujeres casadas, un efecto que compensa en cierta medida el a umento de la oferta de trabajo de las mujeres solteras.

Otra línea de investigacíón se centra en el tamaño de los efectos renta y sustitución en función de las preferencias de los individuos en relación a la renta y el ocio. Brewer et a l. (2003) muestran que, en el Reino Unido, cuanto mayor es el número de hijos, mayor es la preferencia por los ingresos en relación con las horas de trabajo. Además, los niveles más altos de educación se asocian a una menor va loración de los ingresos con respecto a l_a s horas de trabajo, es decir, una hora de trabajo conlleva una menor desutilidad para los padres con a ltos niveles de educación en comparación con los que tienen bajos niveles de educación. En cuanto a la edad, los resultados difieren entre las parejas y las famili a s monoparentales. En las parejas, las preferencias por los ingresos a costa de mayores hora s trabajadas disminuyen con la edad de los progenitores, mientras que, en el caso de las ma dres solteras. el efecto de la edad en la preferencia por los ingresos no pa rece esta r determina do. No obstante. los estudios empíricos constatan que los individuos con una edad superior a la media tienen una ma yor preferencia por las horas de trabajo (MaCurdy, 1992). Adicionalmente, Laroque y Salanié (2000) estiman el salario a l que las mujeres casadas y desemplea das en Francia estarían dispuestas a aceptar un trabajo (salario de reserva ). Estos autores encuentran que el salario de reserva tiende a aumentar con el número de hijos y con la edad, pero no parece verse afectado por la educación. En general, estos conceptos representan una referencia útil a la hora de estimar el impacto de los tipos impositivos marginales sobre comportamiento en el mercado laboral de los diferentes grupos sociodemográficos, así como de los distintos tipos de familia (part:ias o familias monoparentales).

\textbf{lmplicacio11es }

Tomando en consideración el análisis teórico y empírico. podemos destacar una serie de puntos clave a la hora de disefiar estos subsidios salariales.

En primer lugar, la cuestión relativa a la cuantía y el ritmo de retirada de estos impuestos negativos es crucial y, al mismo tiempo, muy dificil de abordar. Como ya se ha mencionado, para los individuos situados en la parte decreciente del EHC, el importe del subsidio disminuye a medida que trabajan más horas o reciben un mayor salario. En general, se ha demostrado que ambas cuestiones dependen de la oferta de trabajo: un aumento de la participación en el mercado laboral o un aumento de las horas trabajadas. A este respecto, Saez (2002) muestra que cuando la mayor parte del efecto se observa en términos de participación en el mercado laboral. el programa óptimo de transferencias debería parecerse al Erre, caracterizado por una pequeña renta garantizada y bajas tasas de retirada del crédito fiscal. Por otro lado, cuando se espera que la mayor parte del efecto se produzca en términos de cambios en las horas trabajadas, la mejor opción sería un programa con una renta garantizada alta y una tasa de retirada elevada.\footnote{%
    Sin embargo. el modelo de Saa no tiene en cuenta el coste adicional Jel desempleo de larga duración en 1érn1inos de pérdida de capital humano. Óscar introducir esta cuestión en la función de bienestar social. probablemente se rcfor1aria la superioridad de los programas con ingresos garantizados relativamente bajos y tasas de retirada bajas.
} 

En segundo lugar, la limitación temporal de estas medidas no solo reduce su impacto en las arcas públicas. sino que permite evitar el fuerte desincentivo sobre la progresión salarial o la inversión en capital humano.

En tercer 1 ugar, autores como Akerlof ( 1978), Parsons ( 1996), Bes ley and Coa tes ( 1992) destacan la necesidad de focalizar estas medidas en colectivos vulnerables que cumplan una serie de características como bajos niveles educativos, familias monoparentales o familias con dependientes a su cargo. En particular, la focalización puede ayudar a reducir los costes públicos de dos maneras: i) de forma más directa. reduce el número de beneficiarios potenciales del programa al centrar la ayuda en los grupos más necesitados; ii) reduce el número de beneficia1ios imprevistos, por ejemplo, al centrarse en individuos que ya han recibido asistencia social durante un determinado periodo de tiempo se reduce la probabilidad de que personas que no están necesitadas intenten beneficiarse de los subsidios.

En cuaito lugar, en caso de coexistir diferentes programas dirigidos a unos mismos potenciales beneficiarios podrían generarse duplicidades que incrementasen los costes públicos, por lo que resulta esencial reforzar la coordinación entre administraciones públicas.

Por último, para evitar í:'l riesgo de que las empresas realicen devaluaciones salariales aprovechando la existencia de subsidios para los trabajadores, resulta necesario establecer una política de salarios mínimos apropiada.

\begin{specialblock}{Evidecia empírica: Programas de asistencia social (ingreso mínimo vital y RBU)}
    Por definición. los subsidios salariales buscan fomentar el empleo de los individuos ocupados de rentas bajas y, por tanto. aquellos desem pleados de bajos ingresos no serían un colectivo beneficiario de esta medida. Por ello. surge la necesidad de completar el sistema con una red de segurídad dirigida a los colectivos vulnerables sin empleo. 
    
    A contínuación, nos centramos en el impacto de la introducción de una renta minuna garantizada y una renta básica universal sobre los incentivos laborales de los individuos. en particular de los trabajadores poco cualificados.\footnote{%
        Renta mínima garantizada: Definida como un programa de asistencia social establecido por el gobierno para proporcionar a cada persona la garantía de unos determinados ingresos mínimos.

    Renta básica Universal: Entendida como una asignación monetaria pública incondicional a toda la población, es decir. independienle de la edad, el nivt>l de r<:nta del individuo. que éste haya cotiü1do pre\'iamenlc. que tenga o no un empleo, o cualquier otra consideración.
    } La crítica principal a la introducción de una renta básica universal se centra en el posible abandono masivo del trabajo remunerado. en esrecial de aquellos trabajadores que ocupan puestos más precarios. Tanto la generosidad como el diseño de las prestaciones sociales pueden desincentivar el trabajo tanto por un efecto renta (mayores ingresos fuera del trabajo) como por un efecto de sustitución (menor rendimiento de incrementos de ingresos debido a la retirada de las prestaciones). Estos desincentivos al trabajo pueden manifestarse tanto en términos de un menor incentivo para aceptar un empleo, es decir, en términos de participación en el mercado laboral (margen "extensivo") como en términos de aumentar los ingresos una vez empleado, principalmente a través del aumento de las horas trabajadas (margen "intensivo"). Los resultados empíricos seíialan que estos desincentivos laborales son especialmente importantes para los trabajadores de bajos ingresos y secunda1ios (Meghir y Phillips, 2010). Por tanto, un reto clave para el diseí'ío de los programas de asistencia social consiste en hacer frente al trade-off entre los objetivos de equidad (lucha contra la pobreza o exclusión social) y eficiencia (incentivos laborales).

    Numerosos autores se centran en el estudio de las denominadas trampas de pobreza. La trampa de la pobreza consiste en el hecho de que una persona que está cobrando una renta mínima garantizada tiene un fuerte desincentivo para aceptar un empleo dado que podría estar sometida a una tasa impositiva marginal del 100\%, es decir, por cada euro de ingreso que reciba como resultado de su nuevo empleo se le descontará otro euro de su prestación.\footnote{%
        En algunos casos de renta mínima. la tasa imposíti\ a marginal es menor ) a que s.: establecen períodos transitorios de eliminación de la prestación
    } Para eliminar este efecto, el salario del nuevo empleo tiene que superar ampliamente la cuantía de la prestación. cosa que no suele suceder dado que los empleos a los que optan los beneficiarios de las rentas mínimas suelen ser de carácter precario. Adicionalmente, tanto los beneficiarios de una renta mínima garantizada como de una renta básica universal se ven afectados por la erosión de su capital humano lo que impacta negativamente en su búsqueda de empleo.\footnote{%
        Esta pérdida de capital humano puede darse como con,ccucncia de la pérdida de conocimi<:ntos técnicos o experiencia adquirida asociada a largos periodos de desempleo o puede ser fruto de los avances tecnológicos de la economía.
    }
    
    Implicaciones
    Para activar la participación en el mercado laboral de los beneficiaiios de estos programas. los países occidentales han adoptado una serie de medidas:
    l . Reducción de la generosidad del programa: los bajos niveles de prestaciones pueden incentivar fuertemente el empleo, aunque esto puede socavar el objetivo principal de aliviar la pobreza.
    2. Reducción de las prestaciones: muchos países introducen un tramo decreciente en las prestaciones. reduciendo gradualmente los derechos a medida que aumentan los ingresos. En algunos casos, la reducción de las prestaciones está vinculada al tiempo y no a los ingresos, con independencia de los ingresos obtenidos. La reducción ayuda a evitar los efectos denominados "cliff-edge·• en los que la obtención de una cantidad marginal de ingresos puede llevar a la pérdida total de una prestación y. por tanto, a un efecto desincentivo desmesurado.
    
    Vinculación al empleo: mientras que la reducción gradual de la prestación suaviza los desincentivos laborales al pem1i tir que los beneficiarios conserven una parte de sus prestaciones, la vinculación al empleo tiende a recompensar a los beneficiarios con una ayuda adicional ("subsidios salariales") por participar en el mercado laboral o aumentar su oferta de trabajo.
\end{specialblock}

\paragraph{Imposición sobre el capital}
Tras las últimas crisis económicas se ha acentuado el debate acerca del papel de la imposición sobre la riqueza o capital. Las necesidades de financiación del sector público o la creciente desigualdad en la distribución de la riqueza son factores que se suelen esgiimir a la hora de justificar un papel creciente de la imposición sobre la riqueza.

Desde un punto de vista teórico, la teoría de la imposición óptima se centra en el trade-off entre eficiencia-equidad. En concreto, este trade-off se acentúa en relación con la imposición sobre el capital en tanto la concentración de la propiedad del capital es mayor que aquella de la renta salarial. Existen d iversos impuestos sobre el capital: impuestos sobre la propiedad y el patrimonio que gravan las existencias de activos, el impuesto sobre la renta de las personas fisicas que grava las rentas del capital percibidas por los pa11iculares sobre los beneficios de las empresas y los impuestos de sucesiones que gravan las transmisiones por causa de muerte.

Siguiendo el trabajo de Piketty, Sacz y Zucman (2013), podemos distinguir cuatro justificaciones para gravar el capital:

1. Argumento de la frontera difusa: en la práctica, la frontera entn.: los impuestos sobre las rentas del capital y los impuestos sobre la propiedad puede ser a veces difusa. Por ejemplo, los autónomos y los propietarios de empresas pueden decidir en gran medida cuánto cobran en salarios y cuánto reciben en dividendos. Esto también se aplica a un gran número de ejecutivos de empresas, cuyos paquetes de remuneración a menudo implican diversos flujos de ingresos. En ocasiones. no resulta sencillo descomponer estos flujos en un componente puramente laboral (pago por servicios laborales) y un componente puramente de capital (compensación por propiedad de capital). De acuerdo con Piketty, Saez y Zucman (2013), esta imprecisión permite justificar que el tratamiento fiscal de los fü1jos de rentas del trabajo y el capital no sean demasiado diferente.
2. Argumento de la capacidad fiscal: estos autores sefialan que los propios flujos de renta (y no sólo su descomposición en componentes de renta del capital y del trabajo) no suelen ser observables para los grandes patrimonios. En consecuencia, la función principal de un impuesto a la riqueza no debe ser hacer progresiva la imposición sobre el conjunto de la renta de las personas físicas. sino que debe gravar de forma indirecta las rentas del patrimonio que son difíciles de observar y que por equidad horizontal deben ser gravadas.

3. Argumento de los incentivos: al gravar el stock de capital en lugar del flujo de renta se incentiva a los agentes a obtener mayores rendimientos.

4. Argumento meritocrático: los individuos no son responsables de su riqueza heredada. por lo que ésta debería gravarse más que las rentas del trabajo. Los mercados imperfectos de capital implican que pa1te del impuesto sobre la herencia debería trasladarse al impuesto sobre el capital a lo largo de la vida.

En concreto, Piketty, Saez y Zucman observan una creciente desigualdad en la distribución de la riqueza, especialmente a través de las herencias. y recomiendan aumentar la importancia de la imposición sobre sucesiones y donaciones como una manera de garantizar la estabilidad social. Su estudio se centra en la tributación óptima de la herencia en un modelo en el que las personas difieren en su capacidad de obtener ingresos y en sus preferencias por las herencias, lo que se traduce en una desigualdad de dotaciones iniciales para futuras generaciones.\footnote{%
    Por ejemplo. incluso los individuos que no reciben ninguna hcr.;ncia podrían prelcrir no gravar dema��iado la herencia. porque ellos mismos valoran mucho la posibilidad de dejar un legado a sus propios hijos.
} Desde el punto de vista de la igualdad de oportunidades. defienden una menor tributación para las rentas del tTabajo o la riqueza generada por uno mismo que para la riqueza heredada. Cabe mencionar que estos resultados sobre la tributación óptima de las herencias también tienen implicaciones sobre los impuestos sobre el capital en vida. Es decir. si se introducen imperfecciones en el mercado de capitales, entonces podría ser óptimo dividir la carga fiscal de la herencia entre un impuesto pagado en el momento de la herencia y un impuesto pagado durante la vida del heredero (ya sea en fom,a de un impuesto sobre el flujo de ingresos del capital o un impuesto sobre la propiedad o el patrimonio aplicado a las acciones).

Por último. en relación con el problema ele la evasión fiscal ya comentado, dada la alta movilidad del capital, el FMI alerta sobre las dificultades de brravar la riqueza en un contexto globalizado. Por ello, recomienda el establecimiento de acuerdos multilaterales de intercambi·o de información entre países y la eliminación de los paraísos fiscales.


\subsubsection{Ahorro nacional}
Ahorro privado frente a ahorro nacional

La discusión de los efectos de los incentivos al ahorro hasta ahora se ha centrado en los impactos sobre el ahorro privado. Lo que importa en última instancia para la determinación de la inversión (y potencialmente del crecimiento) es el ahorro nacional, la suma del ahorro privado y del ahorro público. Recuérdese del Capítulo 4 que aumentar el déficit público puede reducir la inversión al reducir el fondo de capital disponible. Aunque los incentivos fiscales a la jubilación pueden aumentar el ahorro privado, tienen un efecto negativo compensatorio sobre el ahorro nacional porque se financian mediante una desgravación fiscal. La reducción de los ingresos fiscales disminuye el ahorro nacional y compensa cualquier aumento del ahorro privado.

Supongamos que encontráramos que los 401(k) aumentaran el ahorro privado en 30 centavos por cada dólar de aportación: de cada dólar ahorrado en los 401(k), 70 centavos son ahorro para la jubilación que se habría ahorrado incluso en ausencia del 401(k) (pero que ahora se ahorra a través del 401(k) por razones fiscales), y 30 centavos son ahorro nuevo inducido por el 401(k). Supongamos también que el tipo impositivo típico de los contribuyentes fuera del 43 %. Bajo estos supuestos, el coste de cada dólar de aportación a una IRA en términos de ingresos fiscales dejados de recaudar sobre los ahorros existentes reasignados a los 401(k) sería de 30 centavos (0,43 × 70 centavos de ahorro existente). Así, no habría ningún impacto de los 401(k) sobre el ahorro nacional: el aumento del ahorro privado debido al 401(k) sería completamente compensado por la reducción del ahorro nacional debida a la desgravación fiscal sobre ahorros que se habrían producido incluso sin el 401(k) (y que habrían sido gravados en ese caso).

La comparación entre ahorro privado y ahorro nacional vuelve a la noción de impactos marginales de los incentivos fiscales (nuevo comportamiento fomentado) frente a impactos inframarginales de los incentivos fiscales (comportamiento antiguo recompensado) que introdujimos en el Capítulo 18. En el ejemplo aquí, existe cierta respuesta marginal: cada dólar de aportación aumenta el ahorro privado en 30 centavos. Pero también existe una gran respuesta inframarginal: 70 centavos de cada dólar de ahorro en 401(k) iban a producirse incluso si los 401(k) nunca hubieran existido. Esta gran respuesta inframarginal reduce significativamente los ingresos fiscales, compensando cualquier ganancia de ahorro nacional derivada de los incentivos de los 401(k). Así, a menos que exista un gran aumento marginal del ahorro derivado de la disponibilidad de planes de ahorro para la jubilación, estos planes reducirán el ahorro nacional total.

El tamaño de la respuesta marginal e inframarginal a los incentivos fiscales al ahorro dependerá de dos factores. El primero es el tamaño de los efectos renta y sustitución para los ahorradores para la jubilación por debajo del límite de ahorro (por ejemplo, para aquellos con menos de 5.000€ de ahorro para la jubilación). El segundo es la proporción de ahorradores para la jubilación que están por encima del límite de ahorro, para quienes solo existe una respuesta inframarginal: no se fomenta ningún nuevo ahorro para esta población, puesto que no existe efecto sustitución.












\clearpage
\section{Dimensión conductual: Marco de análisis empírico}







\subsection{Aproximación microeconómica}



\subsubsection{Oferta de trabajo}

A continuación, nos centraremos en una de las políticas tributarias más importantes del gobierno federal de Estados Unidos para fomentar la oferta de trabajo: el \textit{Earned Income Tax Credit} (EITC), un programa de subsidio salarial destinado a familias con bajos ingresos.

El estudio del EITC ofrece la oportunidad de aplicar el análisis teórico de la tributación y la oferta de trabajo y de comprender, a partir de una sólida evidencia empírica, cómo la política tributaria del gobierno puede afectar en la práctica a la oferta de trabajo. Finalmente, examinaremos cuál debería ser el tratamiento tributario adecuado de los gastos de cuidado infantil, que pueden constituir un determinante fundamental del comportamiento laboral de los padres.

Tributación y oferta de trabajo: evidencia

Existe una amplia literatura econométrica que estima los efectos de los impuestos sobre la oferta de trabajo. Esta literatura suele distinguir entre perceptores principales de ingresos (\textit{primary earners}) y perceptores secundarios de ingresos (\textit{secondary earners}).

Los perceptores principales son los miembros de la familia que constituyen la principal fuente de ingresos laborales del hogar, mientras que los perceptores secundarios son los demás trabajadores de la familia. Dado que los perceptores principales son los miembros más vinculados al mercado de trabajo, es probable que los perceptores secundarios sean quienes se ocupen de otras actividades del hogar, como el cuidado de los hijos. Tradicionalmente, los perceptores principales eran los maridos y las perceptoras secundarias eran las esposas, que asumían el cuidado de los hijos.

(Nota al pie 15: Existen exenciones a la normativa sobre horas extraordinarias para determinadas categorías de trabajadores, como los empleados ejecutivos, administrativos y profesionales. Los profesores universitarios, por ejemplo, se consideran empleados profesionales porque sus actividades requieren «conocimientos avanzados». Puede encontrarse más información sobre esta normativa en el sitio web del Departamento de Trabajo de los Estados Unidos.)

La conclusión general de esta literatura es doble.

En primer lugar, las decisiones laborales de los perceptores principales responden muy poco a los cambios en sus salarios (como los provocados por los impuestos). Por cada reducción del 10 % en los salarios después de impuestos, los perceptores principales trabajan aproximadamente un 1 % menos de horas, lo que implica una elasticidad de la oferta de trabajo con respecto al salario después de impuestos de 0,1.

Estos estudios también muestran que los perceptores secundarios son mucho más sensibles a las variaciones salariales (y, por tanto, a los impuestos), estimándose normalmente elasticidades de la oferta de trabajo con respecto al salario después de impuestos comprendidas entre 0,5 y 1; para los perceptores secundarios, cada aumento del 1 % en el salario después de impuestos incrementa la oferta de trabajo entre un 0,5 % y un 1 %.

La mayor parte de la respuesta de los perceptores secundarios proviene de la decisión de trabajar o no (participación en el mercado laboral), mientras que una parte menor procede de la decisión sobre cuántas horas trabajar.

Estos resultados son coherentes con el contexto histórico. La elasticidad de la oferta de trabajo respecto al salario después de impuestos viene determinada por la disponibilidad de alternativas al trabajo: cuanto mejores sean las alternativas, más elástica será la oferta de trabajo.

Tradicionalmente, los hombres que eran perceptores principales tenían pocas alternativas al trabajo remunerado, dada la expectativa de que el principal sostén económico del hogar en Estados Unidos trabajara a tiempo completo.

Las mujeres que eran perceptoras secundarias, por el contrario, disponían de una alternativa natural: dedicarse al cuidado de los hijos. Por ello, sus decisiones de ofrecer trabajo eran más elásticas.

Estos resultados también son coherentes con las restricciones sobre las horas de trabajo analizadas en la sección anterior: tradicionalmente, los hombres ocupaban empleos del sector productivo en los que las restricciones horarias eran importantes, de modo que tenían poca capacidad para ajustar sus horas trabajadas al nivel deseado. Las mujeres, por el contrario, trabajaban tradicionalmente en el sector servicios, donde los horarios laborales eran más flexibles.

⸻

Limitaciones de los estudios existentes

Una cuestión importante planteada por esta literatura es que la distinción entre perceptores principales y secundarios se ha ido difuminando.

Mientras que en 1970 solo el 31,9 % de las mujeres casadas trabajaba, en 2013 lo hacía el 59 %.

En la actualidad, algo más de la mitad de los matrimonios estadounidenses tienen a ambos cónyuges trabajando. Entre el 27 % de las familias con un único perceptor de ingresos, en el 19 % el perceptor era el marido y en solo el 7 % era la esposa. Únicamente en el 29 % de los matrimonios en los que ambos cónyuges obtienen ingresos la mujer gana más al año que el hombre.

Todo ello dificulta determinar quién debe considerarse el perceptor principal y quién el secundario al analizar el mercado de trabajo.

Blau y Kahn (2007) encontraron que la elasticidad de la oferta de trabajo de las mujeres casadas se redujo a la mitad entre 1980 y 2000, pasando de una elasticidad de 0,8 en 1980 a 0,4 en 2000.

Este resultado indica que, a medida que los perceptores secundarios se consolidan en el mercado laboral, sus elasticidades salariales se aproximan a las de los perceptores principales y sus decisiones laborales se vuelven menos sensibles a la tributación.

Investigaciones más recientes respaldan esta conclusión y muestran que las respuestas de los perceptores secundarios a los impuestos se están haciendo cada vez más similares a las de los perceptores principales (Chetty et al., 2011).

Otra limitación importante de esta literatura ha sido su concentración en solo un subconjunto de las posibles medidas de la oferta de trabajo: la participación en el mercado laboral y las horas trabajadas.

En realidad, otras dimensiones de la oferta de trabajo también pueden responder a la tributación.

Por ejemplo, supóngase que usted puede obtener un salario más elevado realizando un mayor esfuerzo por cada hora trabajada, pero que ese esfuerzo tiene un coste. En ese caso, solo proporcionará ese esfuerzo adicional si la recompensa en términos de un mayor salario es suficientemente elevada.

Cuando los tipos impositivos son bajos, puede merecer la pena realizar ese esfuerzo porque el salario después de impuestos es suficientemente alto para compensarlo. Sin embargo, cuando aumentan los impuestos, el salario después de impuestos puede dejar de ser suficiente para compensar el coste del esfuerzo adicional, y usted dejará de realizarlo.

De forma similar, supóngase que está considerando dos empleos: en uno es más productivo, pero también está sometido a un mayor nivel de estrés; en el otro es menos productivo, pero trabaja en un entorno más relajado (que, suponemos, usted prefiere).

Solo aceptará el empleo más estresante si el salario después de impuestos es suficientemente elevado para compensar el mayor nivel de estrés.

En este ejemplo, a medida que aumentan los tipos impositivos, será menos probable que acepte el empleo más estresante y mejor remunerado.

Otra dimensión de la oferta de trabajo que puede responder a los cambios impositivos es la decisión sobre cuánto capital humano adquirir. Los efectos de los impuestos en este ámbito, que son más complejos, se analizan en el Capítulo 23.

Lo que realmente importa para las consecuencias de eficiencia social de la tributación es cómo afectan los impuestos al producto total de la sociedad.

Si los individuos no modifican sus horas de trabajo como consecuencia de los impuestos, pero aceptan empleos menos productivos o trabajan con menor intensidad durante cada hora trabajada, la tributación sigue distorsionando sus decisiones laborales y reduce la eficiencia social total.

En el Capítulo 25 volveremos a analizar cómo afectan los impuestos a la renta total obtenida por los individuos (y no solo a las horas trabajadas).

Conviene señalar, no obstante, que la investigación de Chetty et al. (2011), comentada anteriormente, se centró en los ingresos laborales totales y no únicamente en las horas trabajadas, y aun utilizando esta medida más amplia encontró elasticidades reducidas.

Política tributaria para fomentar la oferta de trabajo: el Earned Income Tax Credit (EITC)

En las secciones anteriores analizamos, de forma general, cómo los tipos del impuesto sobre la renta pueden afectar a las decisiones de oferta de trabajo de los perceptores principales y secundarios de ingresos. En esta sección aplicamos esas enseñanzas para estudiar los efectos del Earned Income Tax Credit (EITC), una política del impuesto sobre la renta dirigida específicamente a los trabajadores con bajos ingresos.

El EITC subvenciona los salarios de los trabajadores de bajos ingresos con el fin de alcanzar dos objetivos: redistribuir recursos hacia los grupos de renta más baja y aumentar la cantidad de trabajo ofrecida por dichos grupos.

En el Capítulo 17 analizamos la disyuntiva habitual de los programas redistributivos: incrementan la equidad vertical al redistribuir la renta, pero pueden reducir la eficiencia social al disminuir los incentivos de los grupos de bajos ingresos para obtener rentas.

El EITC ofrece la posibilidad de romper este «triángulo de hierro». Óscar redistribuir la renta mediante subvenciones salariales, el programa pretende aumentar la equidad vertical y fomentar el trabajo entre la población de bajos ingresos.

Antecedentes del EITC

El EITC fue introducido en 1976 y ha experimentado un enorme crecimiento con el paso del tiempo, como muestra la Figura 21-5.

En la actualidad, el gobierno federal destina aproximadamente 58.000 millones de€ anuales al EITC, lo que lo convierte en el mayor programa nacional de lucha contra la pobreza basado en transferencias monetarias.

El EITC ha tenido un claro éxito desde el punto de vista de la equidad vertical: más del 90 % de sus beneficios se destinan a personas con ingresos inferiores a 30.000€.

En esta sección revisaremos la evidencia relativa al otro objetivo del programa: promover la oferta de trabajo de los grupos con bajos ingresos.

Para poder acogerse al EITC, una familia debe tener ingresos laborales anuales superiores a cero e inferiores aproximadamente a 43.941€ si mantiene un hijo, 49.186€ si mantiene dos hijos o 52.427€ si mantiene más de dos hijos.

Una familia sin hijos debe tener ingresos laborales superiores a cero e inferiores aproximadamente a 20.020€.

En el caso de las familias sin hijos, el EITC es considerablemente menor (aproximadamente el 15 % del importe máximo que puede recibir una familia con un hijo).

Es importante señalar que, para todos los tipos de familias, el EITC es reembolsable (refundable): incluso aunque la familia no deba pagar ningún otro impuesto, puede seguir teniendo derecho al crédito, y el gobierno le enviará un cheque por dicho importe en forma de devolución fiscal.

⸻

Estructura actual del EITC

La estructura vigente del EITC se ilustra en la Figura 21-6 para Stacey, una trabajadora soltera con dos hijos.

Durante los primeros 13.870€ de ingresos de Stacey, el gobierno le paga 40 centavos por cada dólar ganado (40 %), hasta alcanzar un crédito máximo de 5.548€ (el 40 % de 13.870€). Por ello, la pendiente del primer tramo de la Figura 21-6 es igual a 0,4.

Durante los 4.240€ siguientes de ingresos (hasta un total de 18.110€), el crédito permanece constante en 5.548€, por lo que el gráfico es horizontal entre 13.870 y 18.110€.

Sin embargo, una vez que los ingresos de Stacey superan los 18.110€, el gobierno comienza a reducir el crédito a razón de aproximadamente 21 centavos por cada dólar adicional ganado (21 %), de modo que quienes obtienen 44.454€ o más dejan de percibir el EITC.

⸻

Efecto del EITC sobre la oferta de trabajo: teoría

Como ocurre con la mayoría de las políticas del impuesto sobre la renta, el EITC produce tanto efectos renta como efectos sustitución sobre las decisiones de oferta de trabajo.

Esto puede observarse incorporando el EITC al diagrama habitual de la elección entre consumo y ocio, como en la Figura 21-7 (la figura no está dibujada a escala para facilitar la ilustración de los efectos sobre la oferta de trabajo).

La línea azul representa la restricción presupuestaria original antes de introducir el EITC, suponiendo un salario de 20€ por hora; la línea verde representa la nueva restricción presupuestaria tras la incorporación del EITC (hasta unos ingresos de 44.454€, punto a partir del cual vuelve a coincidir con la restricción presupuestaria azul).

La figura ilustra el impacto del EITC sobre cuatro grupos diferentes.

1. Personas que no participan en el mercado de trabajo

Para quienes no forman parte del mercado laboral, como la persona A, es probable que el EITC aumente su oferta de trabajo, ya que la única forma de recibir el crédito es incorporándose al mercado laboral.

El EITC no produce ningún efecto renta sobre este grupo porque inicialmente no obtiene ingresos laborales.

Sí produce, en cambio, un efecto sustitución: al aumentar la remuneración derivada de participar en el mercado laboral para cualquier nivel de ingresos inferior a 44.454€, incentiva a los no trabajadores a incorporarse al mercado laboral, desplazándose hasta un punto como A₁.

2. Personas que ya trabajan y ganan menos de 13.870€

Las personas que ya participan en el mercado laboral y obtienen menos de 13.870€, como la persona B, que gana 5.000€, se encuentran en el tramo ascendente (phase-in) del EITC, donde reciben un crédito mayor por cada hora adicional trabajada.

Los efectos sobre la oferta de trabajo de este grupo son ambiguos.

La subvención produce un efecto sustitución que incentiva a trabajar más, porque cada hora de trabajo proporciona una remuneración mayor.

Sin embargo, también produce un efecto renta que induce a trabajar menos, ya que los trabajadores son más ricos gracias a la subvención.

Si predomina el efecto sustitución, los trabajadores se desplazarán hacia un punto como B₁, donde disminuye el ocio y aumenta la oferta de trabajo.

Si predomina el efecto renta, se desplazarán hacia un punto como B₂, donde aumenta el ocio y disminuye la oferta de trabajo.

3. Personas que trabajan y ganan entre 13.870 y 18.110€

Las personas que ya trabajan y obtienen ingresos comprendidos entre 13.870 y 18.110€, como la persona C, que gana 14.000€, se encuentran en el tramo horizontal del EITC, donde reciben exactamente el mismo crédito con independencia de que trabajen más o menos dentro de ese intervalo.

Esta configuración equivale a un desplazamiento paralelo hacia fuera de la restricción presupuestaria, lo que reduce la oferta de trabajo.

Como el EITC no incrementa el salario por hora para las horas adicionales trabajadas, no existe efecto sustitución.

Sin embargo, dado que los trabajadores son más ricos gracias al subsidio obtenido por las horas ya trabajadas, el efecto renta puede reducir el número de horas trabajadas, desplazándolos hacia un punto como C₁.

4. Personas que trabajan y ganan entre 18.110 y 44.454€

Las personas que ya participan en el mercado laboral y obtienen ingresos entre 18.110 y 44.454€, como la persona D, que gana 20.000€, se encuentran en el tramo descendente (phase-out) del EITC, donde el importe del crédito disminuye conforme aumentan los ingresos laborales.

Esta situación provoca que la pendiente de la restricción presupuestaria disminuya con respecto a los tramos anteriores.

El efecto sustitución pasa ahora a reducir la oferta de trabajo, ya que la transferencia pública disminuye con cada hora adicional trabajada.

Óscar mismo tiempo, el efecto renta continúa reduciendo la oferta de trabajo porque el trabajador sigue recibiendo una parte de la transferencia.

En consecuencia, la oferta de trabajo disminuye de forma inequívoca para este grupo, desplazándose hacia un punto como D₁.

Efecto agregado

Si se consideran conjuntamente estos cuatro grupos, la teoría predice que el efecto neto del EITC sobre la oferta total de trabajo de la población con bajos ingresos es ambiguo.

Esta ambigüedad ha motivado la realización de numerosos estudios destinados a analizar los efectos del EITC sobre la oferta de trabajo.


Impacto del EITC sobre la oferta de trabajo: evidencia

La literatura empírica sobre el impacto del EITC se ha centrado en estudiar cómo los cambios en la estructura del EITC a lo largo del tiempo han afectado a la oferta de trabajo.

Un ejemplo de ello se muestra en la Figura 21-8, que ilustra el cambio en el EITC debido a la TRA de 1986. Antes de este cambio legislativo, el EITC era mucho más modesto de lo que es hoy, con una tasa de subsidio de solo el 11 %, un crédito máximo de solo 550€ y una tasa de eliminación gradual (phase-out) del 12,2 %.

La ley de 1986 elevó la tasa de subsidio al 14 % y el crédito máximo a 851€, y redujo la tasa de eliminación gradual al 10 %. La combinación del mayor crédito máximo y de una eliminación gradual más lenta dio lugar a una gran ampliación de la elegibilidad para el crédito: los trabajadores con ingresos comprendidos entre 11.000 y 15.432€ pasaron a ser elegibles por primera vez.

Obsérvese, sin embargo, que ambos programas eran mucho menos generosos que el EITC actual (Figura 21-6), porque desde 1986 se han producido ampliaciones todavía mayores, particularmente en 1993. Además, a diferencia de la situación actual, tanto antes como después de la TRA de 1986, el EITC solo estaba disponible para familias con hijos.

La ampliación de la generosidad del EITC después de 1986 tuvo efectos ambiguos sobre la oferta de trabajo.

Para quienes no trabajaban, hubo un aumento claro del rendimiento del trabajo, por lo que la oferta de trabajo debería haber aumentado. Para quienes ya trabajaban, el efecto sustitución hacia un mayor trabajo se reforzó, tanto por el aumento de la tasa de subsidio en el tramo de entrada (phase-in) como por la reducción de la tasa a la que el subsidio se eliminaba gradualmente (phase-out).

Óscar mismo tiempo, el efecto renta hacia un menor trabajo se reforzó al aumentar el nivel máximo del EITC de 550 a 851€.

Además, recordando la discusión del «triángulo de hierro» del Capítulo 17, un nuevo conjunto de individuos para quienes el EITC había sido anteriormente irrelevante (aquellos con ingresos entre 11.000 y 15.432€) pasó a ser elegible y, para esos individuos, existían efectos renta y efectos sustitución, ambos conduciendo a una menor oferta de trabajo.

La literatura que evalúa el efecto del EITC ha llegado a varias conclusiones claras, que se analizan en las siguientes secciones.

Efectos sobre la participación en la población activa

La mayoría de los estudios en este ámbito han examinado el impacto del EITC sobre la participación en la población activa de las madres solteras (es decir, sobre su decisión de trabajar).

Existe un fuerte consenso entre estos estudios en que el EITC ha desempeñado un papel importante en el aumento de la proporción de madres solteras que trabajan.

Por ejemplo, como se describe en el recuadro Empirical Evidence, Eissa y Leibman (1996) estimaron que las madres solteras tenían entre 1,4 y 3,7 puntos porcentuales más de probabilidad de trabajar como resultado de la ampliación del EITC en 1986.

Existe una creciente evidencia internacional sobre el impacto de subsidios al mercado de trabajo de este tipo.

Blundell y Hoynes (2004) resumieron la evidencia procedente de un programa similar al EITC en el Reino Unido y encontraron que los efectos iban en la misma dirección que los del EITC, pero eran menores. Argumentan que ello se debe a que las ampliaciones del programa británico fueron acompañadas por ampliaciones simultáneas de programas que beneficiaban a quienes no trabajaban, compensando las ganancias netas derivadas del aumento de los incentivos al trabajo.

Bravo y Rau (2013) estudiaron un programa de subsidios salariales para jóvenes en Chile y encontraron aumentos muy grandes en la participación en la población activa durante el primer año, que se atenúan en cierta medida con el paso del tiempo.

Por tanto, el resultado general parece ser que los subsidios salariales del tipo del EITC aumentan la participación en la población activa, aunque la magnitud del efecto puede variar.

Efectos sobre las horas de trabajo

Siguiendo la teoría que acabamos de analizar, cabría esperar que la ampliación del EITC en 1986 aumentara la participación en la población activa, pero que también pudiera reducir las horas trabajadas por quienes ya estaban en la población activa, especialmente las personas cuyos ingresos se encontraban en el tramo de eliminación gradual (phase-out).

En la mayor parte de ese intervalo, el tipo marginal impositivo sobre los ingresos adicionales era del 42,5 % (un tipo marginal federal del 15 %, un impuesto sobre las nóminas del 15,3 % y una tasa de eliminación gradual del EITC del 12,2 %).

Este elevado tipo impositivo podría haber desincentivado la oferta de trabajo.

Eissa y Leibman también examinaron el impacto de la ampliación del EITC sobre las horas trabajadas entre quienes ya participaban en la población activa y, sorprendentemente, no encontraron ningún efecto.

Es decir, el subsidio salarial del EITC ha tenido un gran efecto para incorporar personas a la población activa; sin embargo, una vez que las personas ya participan en ella y trabajan, los impuestos sobre los ingresos marginales a través del EITC no tienen un efecto importante sobre el número de horas que trabajan.

Este resultado puede producirse porque la decisión de participar en la población activa con respecto a los subsidios salariales es elástica, mientras que la decisión sobre cuántas horas ofrecer con respecto a la tributación de los salarios no es elástica.

Óscarternativamente, este resultado puede producirse porque los individuos cuyos ingresos se encuentran en el tramo de eliminación gradual (phase-out) no comprenden realmente los complicados desincentivos de dicho tramo y, por tanto, no responden reduciendo sus horas de trabajo, a pesar de los incentivos financieros para hacerlo.

Impacto sobre los matrimonios

Aunque el EITC parece tener efectos positivos sobre la oferta de trabajo de las madres solteras, puede tener efectos negativos sobre las decisiones de oferta de trabajo de los matrimonios.

Considérese un matrimonio que toma una decisión secuencial sobre la oferta de trabajo: primero el marido decide cuántas horas trabajará y después la esposa decide cuántas horas trabajará.

Para el marido, los efectos teóricos del EITC son los mismos que para una madre soltera (quienes no participan en la población activa comenzarán a trabajar, pero los efectos sobre la oferta de trabajo de quienes ya trabajan son ambiguos).

Para la esposa, sin embargo, los efectos del EITC serán más negativos porque el EITC se calcula sobre la base de los ingresos totales de la familia.

Para la mayoría de las familias de bajos ingresos con hijos, esto significa que la familia ya habrá alcanzado el crédito máximo con los ingresos del marido, de modo que la esposa solo afrontará el tramo descendente de eliminación gradual (phase-out) al decidir cuánto trabajar.

En efecto, entre los matrimonios con bajo nivel educativo y con hijos, el 9 % se encuentra en el tramo de entrada (phase-in), el 6 % en la meseta (beneficio máximo) y el 43 % en el tramo de eliminación gradual (phase-out).

Como se ha señalado, en el tramo de eliminación gradual existen fuertes desincentivos para trabajar, porque cada dólar adicional de ingresos reduce los pagos del EITC en 21 centavos.

Por tanto, los perceptores secundarios de ingresos tienen un fuerte desincentivo para trabajar a través del EITC.

Eissa y Hoynes (1998) estudiaron el impacto del EITC sobre los matrimonios.

Encontraron que no existía ningún efecto sobre el trabajo de los hombres.

Este resultado es coherente con los resultados obtenidos para las madres solteras (según los cuales el EITC aumentó la participación en la población activa de las madres solteras, pero no afectó a sus horas de trabajo una vez que ya estaban trabajando), porque casi todos los hombres participan en la población activa desde el principio.

Óscar mismo tiempo, los investigadores encontraron una reducción moderada de la oferta de trabajo de las mujeres casadas, un efecto que compensa en cierta medida los aumentos de la oferta de trabajo de las mujeres solteras.

Resumen de la evidencia

En conjunto, la experiencia del EITC en Estados Unidos parece haber sido bastante satisfactoria.

Es un poderoso instrumento redistributivo que actualmente transfiere más dinero en efectivo a las familias de bajos ingresos que cualquier otro programa asistencial de Estados Unidos.

Y lo ha hecho sin reducir la oferta total de trabajo, que era el problema de los programas tradicionales de asistencia monetaria.

En su lugar, esta redistribución se ha asociado con un aumento de la oferta de trabajo entre las madres solteras (un aumento de la participación en la población activa sin una reducción compensatoria de las horas trabajadas), ningún efecto sobre los padres y una reducción moderada de la oferta de trabajo entre las madres casadas.

(Estos resultados no implican que la asistencia monetaria deba abolirse y sustituirse por un EITC de mayor tamaño. Óscargunos cabezas de familia simplemente no pueden obtener ingresos suficientes para vivir —quizá porque poseen muy poca cualificación—, por lo que sigue existiendo la necesidad de una transferencia basada exclusivamente en la renta que no requiera trabajar. Los ejercicios de simulación de Saez (2000) sugieren que, si existe un EITC, el esquema redistributivo óptimo debería seguir incluyendo una renta garantizada, como la analizada en la sección sobre el impuesto óptimo sobre la renta del Capítulo 20, junto con incentivos al trabajo similares a los del EITC para los perceptores de rentas más bajas.)

Además, una serie de estudios recientes ha demostrado beneficios de largo plazo para los hijos derivados de unos pagos más elevados del EITC: una menor incidencia de nacimientos con bajo peso al nacer, una mejora del rendimiento escolar y mayores tasas de matriculación en la universidad.

El tratamiento fiscal del cuidado infantil y su impacto sobre la oferta de trabajo

Una enseñanza principal de la literatura empírica sobre la oferta de trabajo es que las decisiones de oferta de trabajo de los perceptores secundarios de ingresos son muy sensibles a los salarios después de impuestos. Sin embargo, los salarios claramente no son el único factor que determina tales decisiones de trabajo. Otro factor importante probablemente sea el coste del cuidado infantil para los hijos de padres que trabajan. Actualmente, solo el 21,3 % de los niños en edad preescolar son cuidados principalmente por uno de sus padres; el 33,7 % son cuidados por otros familiares, y el 32,9 % son puestos al cuidado de personas no familiares, como guarderías, centros de cuidado diurno y hogares de cuidadores infantiles. En hogares con ingresos superiores a 54.000€, solo el 24 % de los niños en edad preescolar son cuidados principalmente por sus padres, mientras que el 57 % son colocados en centros de cuidado infantil. Entre los gastos de los padres y los subsidios gubernamentales, los gastos totales en cuidado infantil en Estados Unidos superan ampliamente los 80.000 millones de€ anuales.

Teóricamente, los gastos en cuidado infantil para los padres que trabajan operan exactamente de la misma manera que la tributación salarial. Cuanto mayores sean los gastos en cuidado infantil, menores serán los rendimientos netos del trabajo. Así, los gastos en cuidado infantil tienen efectos sustitución que conducen a menos trabajo y efectos renta que conducen a más trabajo; para los perceptores secundarios de ingresos, la literatura empírica sugiere que dominan los efectos sustitución, haciendo que los perceptores secundarios trabajen menos (Anderson y Levine, 2000).

El tratamiento fiscal del cuidado infantil

El tratamiento fiscal del cuidado infantil también plantea interesantes cuestiones adicionales relacionadas con la discusión del estándar de renta integral de Haig-Simons introducido en el Capítulo 18. Bajo el sistema del impuesto sobre la renta de Estados Unidos, el trabajo prestado a través del mercado está gravado, mientras que el trabajo prestado a través de actividades no de mercado, como el cuidado infantil en el hogar, no está gravado. Este enfoque es inequitativo porque las familias que eligen proporcionar ellas mismas el cuidado infantil, en lugar de obtener ingresos y después comprar servicios de cuidado infantil, pagan menos impuestos. También es ineficiente porque subsidia el cuidado infantil en el hogar frente al cuidado infantil de mercado.

Considérese mi familia, donde yo soy el perceptor principal de ingresos. Tenemos la opción de que mi esposa cuide de nuestros tres hijos o de enviarlos a cuidado infantil y que ella vaya al mercado a trabajar. Supongamos que mi esposa puede ganar 25€ por hora y que el cuidado infantil de cada uno de nuestros hijos cuesta 5€ por hora, para un total de 15€ por hora. Supongamos además que el tipo impositivo sobre los ingresos de mercado de mi esposa es del 50 %, incluidos los impuestos federales sobre la renta, los impuestos sobre las nóminas y los impuestos estatales sobre la renta. Finalmente, supongamos que mi esposa es igualmente feliz trabajando o cuidando de nuestros hijos, de modo que decidimos tomar esta decisión puramente por razones financieras. Este ejemplo se ilustra en la Tabla 21-1.

Si mi esposa trabaja 40 horas por semana, ganaría 1.000€ cada semana, pero tendríamos que poner a nuestros hijos en cuidado infantil durante esas 40 horas, lo que nos costaría 600€. Deberíamos 500€ en impuestos sobre sus ingresos de mercado, de modo que, después de impuestos, estaría ganando 500€ durante la semana, lo cual es menos que los 600€ que pagaríamos en cuidado infantil. Así, aunque el salario antes de impuestos de mi esposa de 1.000€ es superior al coste del cuidado infantil, su salario después de impuestos es inferior. En otras palabras, en términos después de impuestos, el cuidado infantil que mi esposa proporciona en casa vale 600€ por semana (la última columna), mientras que el salario después de impuestos que lleva a casa procedente del trabajo de mercado es solo de 500€ (la penúltima columna). La tributación del trabajo de mercado pero no del trabajo doméstico (su valor al cuidar de nuestros hijos) ha creado una cuña fiscal que pone al trabajo de mercado en desventaja. Las cuñas fiscales se crean siempre que dos actividades comparables son gravadas a tipos diferentes.

Pista rápida

Hemos utilizado el término cuña fiscal de dos maneras diferentes en el Capítulo 19 y en esta discusión. En el Capítulo 19, analizamos las cuñas fiscales dentro de un único mercado; en ese contexto, las cuñas fiscales son la diferencia que los impuestos causan entre los precios del productor y del consumidor en ese mercado. Aquí, nos referimos a cuñas fiscales entre mercados de factores, donde las cuñas fiscales son la diferencia entre los rendimientos del factor (oferta de trabajo) en los distintos mercados. Ambas son definiciones válidas y constituyen un subconjunto de la definición más amplia de cuñas fiscales: cualquier diferencia entre los rendimientos antes y después de impuestos de una actividad causada por los impuestos.

Opciones para resolver las cuñas fiscales

Este ejemplo plantea un punto general sobre las cuñas fiscales entre actividades gravadas y no gravadas: tales cuñas fiscales distorsionan el comportamiento al animar a las personas a realizar las actividades no gravadas y causan pérdida irrecuperable de eficiencia. Óscar gravar el trabajo de mercado pero no el trabajo doméstico, redujimos la productividad total en 400€ porque mi esposa renunció a la oportunidad de ganar 1.000€ para quedarse en casa y proporcionar 600€ de cuidado infantil. Esta reducción de la productividad aumenta la pérdida irrecuperable de eficiencia: el pastel de eficiencia social es ahora 400€ menor porque mi esposa está realizando la actividad para la que tiene un producto marginal menor. Los economistas de hacienda pública dicen a menudo que tales cuñas fiscales crean un terreno de juego desigual entre actividades económicas, donde los individuos son tratados de manera diferente debido a las elecciones que realizan.

Imputación de ingresos domésticos

En este caso, hay dos maneras mediante las cuales los responsables de las políticas podrían nivelar el terreno de juego a favor del trabajo de mercado, como se ilustra en las dos filas siguientes de la Tabla 21-1. Una sería gravar el trabajo en el hogar del mismo modo que gravamos el trabajo en el mercado. El gobierno lo haría imputando ingresos domésticos, o asignando un valor de ingresos domésticos a mi esposa basado en la cantidad de cuidado infantil que proporciona. Así, el gobierno podría imputar que ella proporciona 600€ de cuidado infantil doméstico porque eso es lo que nos costaría si tuviéramos a nuestros hijos en cuidado de mercado. Mi familia sería entonces gravada por ese cuidado infantil doméstico, de modo que si ella no trabajara en el mercado, pagaríamos 300€ en impuestos (el 50 % de los 600€ de ingresos domésticos imputados) y tendríamos un valor después de impuestos del trabajo doméstico de 300€. Con este sistema fiscal en vigor, aumentamos nuestros ingresos en 200€ si ella trabaja (500€ en lugar de 300€), por lo que irá a trabajar.

Como se puede sospechar, esta es una solución impracticable. Hay una serie de cuestiones abrumadoras a las que se enfrentan los responsables de las políticas cuando intentan gravar un bien sin un valor de mercado claramente asignado, como el cuidado infantil proporcionado en el hogar.

Costes de cuidado infantil deducibles

La otra alternativa es hacer deducibles los costes del cuidado infantil de mercado. Supongamos ahora que el trabajo en el hogar no está gravado, pero que el gobierno permite a cada familia deducir el coste del cuidado infantil de su renta imponible. Bajo esta regla, la renta imponible de mi esposa caería de 1.000€ a 400€ (porque deduciríamos los 600€ pagados por cuidado infantil), y nuestra factura fiscal sobre sus ingresos caería a 200€. Sus ingresos después de impuestos procedentes del trabajo serían de 800€, superando los 600€ de costes de cuidado infantil; en términos netos, ganaríamos 200€ por semana por su trabajo de mercado, la misma cantidad que si el trabajo doméstico estuviera gravado. Así, hay dos maneras de igualar el terreno de juego: gravando todas las actividades por igual o subsidiando todas las actividades por igual.

Comparación de las opciones

Aunque ambas soluciones nivelan el terreno de juego, sus efectos sobre la base imponible no son idénticos: permitir una deducción por costes de cuidado infantil ha reducido la base imponible. Cuando el gobierno grava el trabajo doméstico y mi esposa decide ir a trabajar, el gobierno recauda 500€ en impuestos. Pero cuando el gobierno subsidia el cuidado infantil y mi esposa decide ir a trabajar, el gobierno recauda solo 200€ en impuestos. Así, el coste de nivelar el terreno de juego para esta decisión de cuidado infantil es la reducción de 300€ del tamaño total de la base imponible. A medida que disminuye el tamaño de la base imponible, los tipos deben aumentar para recaudar un nivel dado de ingresos, y la pérdida irrecuperable de eficiencia aumentará a medida que aumenten los tipos impositivos.

Así, nos enfrentamos a tres opciones, todas las cuales tienen inconvenientes. Podemos continuar teniendo un terreno de juego desigual, lo que reduce la eficiencia social al disuadir a las madres del trabajo de mercado. Podemos igualar el terreno de juego gravando el trabajo doméstico, lo que tiene el mayor sentido económico pero es una pesadilla administrativa. Podemos nivelar el terreno de juego ofreciendo subsidios al trabajo de mercado, lo que reduce la eficiencia global del sistema fiscal.

Descartando la segunda opción (gravar el trabajo doméstico) por razones administrativas, ¿cómo decidimos cuál de las dos opciones restantes causa menos distorsión global a la economía y es, por tanto, más eficiente? Un sistema que no tiene deducción por costes de cuidado infantil tiene el coste de distorsionar a las madres alejándolas del trabajo de mercado y orientándolas hacia el trabajo doméstico, pero un sistema que tiene una deducción por costes de cuidado infantil reduce la base imponible y, por tanto, la eficiencia global del sistema fiscal. ¿Cuál es peor?

Como se destacó en el capítulo anterior, estos costes de eficiencia dependerán de las elasticidades de las actividades gravadas. El coste de eficiencia de reducir la base imponible global estará determinado por la elasticidad global de la actividad económica respecto a la tributación: como analizamos en el Capítulo 25, la mayoría de las estimaciones sugieren que esta elasticidad es modesta. El coste de eficiencia de distorsionar a las madres alejándolas del trabajo de mercado estará determinado por la elasticidad de la oferta de trabajo de mercado de las madres respecto al salario: como se analizó anteriormente en este capítulo, la mayoría de las estimaciones sugieren que esta es bastante grande (aunque está disminuyendo con el tiempo). Así, parece que subsidiar el cuidado infantil para reducir la cuña fiscal es la mejor de las dos opciones.

En Estados Unidos, hemos elegido una posición entre estas dos alternativas subsidiando parcialmente el cuidado infantil de mercado. En particular, las familias reciben un crédito fiscal por costes de cuidado infantil de hasta 3.000€ por un hijo o 6.000€ por dos o más. El crédito se calcula como un porcentaje de esos costes dependiente de la renta, de modo que los hogares con ingresos inferiores a 15.000€ pueden reclamar el 35 % de los costes de cuidado infantil (hasta 1.050 o 2.100€), mientras que los hogares con ingresos superiores a 43.000€ pueden reclamar el 20 % de los costes de cuidado infantil (hasta 600 o 1.200€).

Conclusión

La discusión sobre la tributación óptima de la renta en el Capítulo 20 ilustró el papel clave desempeñado por las respuestas conductuales a los tipos impositivos. Si impuestos más altos llevan a las personas a cambiar su comportamiento para ofrecer menos trabajo, estos cambios pueden compensar las ganancias derivadas de los aumentos de impuestos y podría existir un límite natural a los ingresos que pueden recaudarse mediante la tributación sobre la renta.

En este capítulo, revisamos la evidencia sobre cómo responde la oferta de trabajo a la tributación. La mayoría de los estudios muestran que los tipos impositivos tienen poco impacto sobre la oferta de trabajo de los perceptores principales de ingresos, pero un impacto más sustancial sobre los perceptores secundarios. También hemos revisado una de las principales políticas fiscales para promover la oferta de trabajo, el EITC, y analizado evidencia que muestra que el EITC ha aumentado la oferta de trabajo de los perceptores de bajos ingresos. Finalmente, analizamos el tratamiento fiscal apropiado del cuidado infantil, uno de los principales impedimentos para la oferta de trabajo de los perceptores secundarios de ingresos.


\subsection{Ahorro}
Evidencia: ¿Cómo afecta el tipo de interés después de impuestos al ahorro?

En contraste con el caso de la oferta de trabajo, existe poco consenso sobre el impacto de los impuestos o del tipo de interés en las decisiones de ahorro. De hecho, los economistas ni siquiera están de acuerdo sobre si los impuestos sobre la renta tienen un impacto negativo, nulo o incluso positivo sobre el ahorro. Hall (1988) utilizó evidencia de series temporales para sostener que el tipo de interés después de impuestos no tiene efecto sobre el consumo ni sobre el ahorro. Otros estudios sugieren que las decisiones de consumo sí responden de forma importante al tipo de interés después de impuestos (Attanasio y Weber, 1995; Gruber, 2006).

Para ser justos, estudiar las relaciones entre los tipos de interés después de impuestos y el ahorro es un problema difícil. Determinar el tipo de interés apropiado que debe utilizarse es uno de los aspectos más difíciles de este problema: aunque podemos medir el salario de un trabajador determinado, es difícil medir el tipo de interés apropiado para cualquier ahorrador determinado. Qué tipo de interés es «apropiado» depende del conjunto de oportunidades de ahorro disponibles para un individuo, que van desde cuentas bancarias hasta planes de pensiones subsidiados fiscalmente. Además, el interés que puede obtenerse sobre cualquier tipo de ahorro cambia normalmente a lo largo del tiempo de la misma manera para todos los individuos, lo que hace difícil encontrar grupos de tratamiento y de control apropiados para estudiar cómo responde el ahorro a los cambios en el tipo de interés. No obstante, la elasticidad del ahorro con respecto al tipo de interés después de impuestos sigue siendo un parámetro crucial para el análisis de políticas, y claramente se necesita más investigación para evaluar su magnitud.


Evidencia a favor del modelo precautorio

En apoyo del modelo de ahorro precautorio, varios estudios han mostrado que una mayor incertidumbre conduce a un mayor ahorro y que reducir la incertidumbre reduce el ahorro. Otros estudios han mostrado que las ampliaciones de programas de seguro social que reducen la incertidumbre sobre la renta también reducen el ahorro, un resultado que es coherente con la idea de que el ahorro está motivado por la precaución frente al riesgo.

Evidencia a favor del modelo de autocontrol

Existe una evidencia creciente que es coherente con los modelos de autocontrol del ahorro. Recuérdese de nuestra discusión de los modelos de autocontrol en el Capítulo 6 que los individuos con problemas de autocontrol demandarán mecanismos de compromiso para ayudar a limitar sus problemas de autocontrol. Tales mecanismos de compromiso están muy extendidos en el contexto del ahorro. Un ejemplo clásico es el «Christmas Club», una cuenta bancaria en la que los individuos depositan dinero a lo largo del año, con un interés bajo o nulo, para asegurarse de que tienen dinero disponible en Navidad para comprar regalos. Otro ejemplo son las cuentas de ahorro para la jubilación del tipo descrito en la siguiente sección.

Una segunda evidencia que es coherente con las explicaciones de autocontrol es el patrón de acumulación de activos en Estados Unidos, que tiene una característica muy extraña: los individuos tienen ahorros sustanciales en formas a las que es difícil acceder (vivienda o cuentas de jubilación), pero muy poco en formas más fácilmente accesibles, como las cuentas corrientes. Aún más desconcertante, muchos individuos ahorran dinero en formas a las que es difícil acceder, como las cuentas de jubilación, a tipos de interés bastante modestos (5 % o menos), mientras que simultáneamente piden prestado dinero con sus tarjetas de crédito a tipos de interés elevados (10 % o más). En el modelo estándar, estos individuos podrían mejorar su situación ahorrando algo menos al 5 % y no pidiendo prestado al 10 %. El hecho de que no lo hagan es muy coherente con el modelo de autocontrol: los individuos están preocupados de que, si tienen el dinero en sus manos, lo gastarán imprudentemente, pero si pueden comprometerse a mantenerlo lejos de ellos mismos, será ahorrado. Así, los individuos pueden ahorrar con éxito en formas como la vivienda o las cuentas de jubilación, ¡pero no pueden ahorrar con éxito una vez que tienen el dinero en sus manos!

Finalmente, una evidencia muy convincente a favor de los modelos de autocontrol se encuentra en un experimento llevado a cabo por los economistas Richard Thaler y Shlomo Benartzi (2004), quienes ofrecieron un plan único de ahorro para la jubilación a empleados de una empresa manufacturera de tamaño mediano. Con su plan «ahorra más mañana» (save more tomorrow), los empleados comprometían una parte de sus futuros aumentos salariales a sus ahorros para la jubilación. A los empleados se les daba la opción de asumir este compromiso de ahorro mucho antes de que apareciera el aumento salarial real, de modo que la decisión parecía menos difícil. Una vez que los empleados se incorporaban, su contribución al plan de ahorro aumentaba con cada subida salarial que recibían.

Un plan de este tipo no debería haber tenido ningún atractivo para el tipo de ahorradores racionales y orientados al futuro del modelo estándar de consumo intertemporal. Podían ahorrar tanto como quisieran de su renta, de modo que no había necesidad de comprometerse previamente con un plan de ahorro fijo. En otras palabras, en el modelo estándar, rara vez existe una razón para restringir el propio comportamiento de esta manera; si una trabajadora quiere ahorrar una parte de los aumentos salariales futuros, simplemente puede hacerlo cuando se produzca el aumento. Sin embargo, este plan tendría un gran atractivo para quienes tienen problemas de autocontrol, que querrían mantener parte de sus futuros aumentos salariales fuera de las manos de sus yoes impacientes de corto plazo. Los individuos con problemas de autocontrol podrían temer que, si no se comprometen ahora a ahorrar sus aumentos, entonces, cuando esos aumentos realmente se produzcan, los gastarán en un arrebato de impaciencia de corto plazo.

De hecho, cuando se introdujo este plan, el 78 % de los empleados a quienes se ofreció decidió incorporarse, y el 80 % de esos empleados permaneció en el plan durante cuatro aumentos salariales. Lo más llamativo es que las tasas de ahorro de los participantes en el plan aumentaron del 3,5 % al 13,6 % a lo largo de 40 meses. El hecho de que una forma de compromiso que es inútil en el modelo estándar tenga un efecto tan grande sobre el comportamiento de ahorro es, una vez más, coherente con la idea de que los empleados tienen problemas de autocontrol respecto a su ahorro. Ashraf et al. (2006) documentaron un resultado similar para un mecanismo de compromiso para el ahorro en Filipinas, demostrando que los problemas de autocontrol son un fenómeno mundial.

Evidencia sobre los incentivos fiscales y el ahorro

Existen varios estudios sobre los impactos de las pensiones proporcionadas por el empleador, las IRA y los 401(k) sobre el comportamiento de ahorro.

Como se analiza en el recuadro de Evidencia Empírica, identificar una estimación causal del impacto de estos incentivos sobre el ahorro ha resultado ser sorprendentemente difícil. Pero la evidencia de un importante estudio reciente indicó que los incentivos de precio fiscal para aumentar el ahorro no tienen mucho efecto sobre el ahorro privado total (al tiempo que reducen el ahorro nacional). Óscar mismo tiempo, este estudio y otros han sugerido que no son los ahorros fiscales, sino otros factores del diseño del programa, los que tienen el mayor impacto sobre el efecto de los incentivos a la jubilación.

Por ejemplo, Madrian y Shea (2001) siguieron la participación en 401(k) en una empresa que pasó de un sistema en el que los trabajadores tenían que inscribirse activamente en un plan 401(k) si querían participar, a un sistema en el que los trabajadores eran incluidos por defecto en el 401(k) y tenían que salirse si no querían participar. Según la teoría económica estándar, estos dos sistemas deberían tener efectos idénticos porque los individuos tienen la opción de participar o no en cada caso. En la práctica, sin embargo, cambiar la opción por defecto de una en la que los individuos no están inscritos automáticamente (y tienen que inscribirse activamente) a una en la que los individuos están inscritos (y tienen que darse de baja activamente) tuvo efectos enormes, elevando las tasas de participación en los planes 401(k) de la empresa entre los nuevos empleados de alrededor del 50 % a casi el 90 %. Estos efectos fueron particularmente grandes para los grupos desfavorecidos: para aquellos con ingresos inferiores a 20.000€ al año, la participación en los planes 401(k) aumentó del 13 % al 80 % a partir de este simple cambio.

Este tipo de resultados motivó al presidente Obama a ampliar aún más el enfoque que se analizó en la introducción. Una propuesta posterior de la administración Obama planteó exigir que las empresas que no tuvieran planes de jubilación establecieran e inscribieran automáticamente a los empleados en cuentas de tipo IRA. Los empleados tendrían el 3 % de su salario depositado directamente en una cuenta de jubilación a menos que optaran por no participar. El presidente también propuso ampliar el crédito del ahorrador para ayudar a las familias de menores ingresos a contribuir a estas nuevas cuentas de jubilación. Sin embargo, la propuesta del presidente no fue aprobada por el Congreso y, en consecuencia, los estados están elaborando sus propias leyes relativas a cuentas de jubilación obligatorias.

Conclusión

Una de las decisiones más importantes tomadas por los contribuyentes en Estados Unidos es cuánto ahorrar, y parece probable que los impuestos influyan en esa decisión. Desafortunadamente, ni la teoría ni la evidencia empírica existente ofrecen una enseñanza clara sobre la magnitud (o incluso la dirección) del efecto de los impuestos sobre el ahorro.

A pesar de esta falta de evidencia, los incentivos fiscales al ahorro continúan creciendo en importancia. En 1975, el gasto fiscal en incentivos al ahorro era inferior a 20.000 millones de€; en 2014, había crecido hasta 141.000 millones de€.¹⁸ Claramente, los responsables de las políticas creen que los incentivos fiscales pueden marcar una diferencia en las decisiones de ahorro de los individuos. Se necesita investigación futura para evaluar la validez de esa creencia.


\subsection{Curva LAFFER}
Las primeras investigaciones se centraron en modelos estáticos de oferta de trabajo. STUART (1981) construye un modelo simple con un sector gravado y un sector no gravado en Suecia. El sector no gravado incluye la evasión fiscal ilícita, así como el ocio y la producción doméstica (como el cultivo de un huerto doméstico), de forma que el valor de $t^*$ se encuentra alrededor del $70\%$. Posteriomente, FULLERTON (1982) construye un modelo multisectorial con datos estadounidenses y estima un tipo impositivo óptimo del $79\%$.

Estas investigaciones están sujetas a una serie de problemas. En primer lugar, no conocemos realmente la elasticidad de la oferta de trabajo y la heterogeneidad de trabajadores dificulta su estimación. En segundo lugar, tampoco conocemos el tipo impositivo actual ya que los sistemas fiscales actuales son combinaciones complejas de impuestos sobre la renta, el consumo y las ventas. Además, el impuesto sobre la renta es progresivo, lo que da lugar a diferentes tipos según la renta de los individuos. Toda esta heterogeneidad significa que no existe un único tipo impositivo óptimo que maximice la recaudación. Incluso si ignoramos la heterogeneidad, un sistema progresivo implica que el tipo impositivo marginal (que afecta a los incentivos) supera el  tipo impositivo medio (que afecta a los ingresos) de forma que la cuestión de cómo afecta a los ingresos un cambio en el tipo impositivo marginal no está bien definida, porque también hay que especificar cómo afecta la reforma a los tipos medios. Por último, de acuerdo con autores como MALCOMSON la propia curva de LAFFER no está bien definida con los ingresos fiscales en el eje vertical ya que, para él, lo relevante es cómo se gastan esos ingresos. Para ilustrar esta concepción, MALCOMSON (1986) muestra que la curva de LAFFER puede seguir inclinándose hacia arriba hasta llegar a un tipo impositivo del $100\%$, lo que en esencia explica que no hay ningún rango prohibitivo. Sin embargo, GAHVARI (1989) muestra cómo este resultado depende del supuesto de que los ingresos fiscales se utilicen para proporcionar un bien público que sea separable en utilidad. De este modo, la subida de impuestos tiene un efecto sobre la renta que aumenta el esfuerzo laboral, y la recaudación puede seguir aumentando. Sin embargo, si el aumento de los ingresos se destina a transferencias a tanto alzado, esta transferencia tiende a compensar el efecto renta, dejando sólo el efecto sustitución. 

Desde un punto de vista dinámico, AGELL y PERSSON (2001) construyen un modelo de crecimiento endógeno unisectorial con el capital como único insumo y sin mano de obra, obteniendo un \textit{efecto LAFFER dinámico} según el cual el gobierno puede reducir el tipo impositivo y aun así aumentar al menos un año los gastos o ingresos públicos.

Los cambios de la oferta de trabajo no son la única vía a través de la que el aumento de los tipos impositivos puede afectar a la recaudación. Cuando los tipos aumentan, las personas pueden sustituir salarios por formas de renta no gravadas, de modo que, aunque la oferta de trabajo no se altere es posible que los ingresos impositivos se reduzcan. Del mismo modo, las personas (especialmente las que perciben rentas elevadas) pueden sustituir rentas del capital sujetas a gravamen por ciertos tipos de rentas de capital no gravadas (como los intereses de la deuda) o llevar a cabo comportamientos estratégicos que reduzcan su deuda tributaria. 

A partir de un análisis de datos procedentes de declaraciones de renta, GRUBER y SÁEZ (2002) concluyen que, especialmente para las personas de rentas altas, los tipos impositivos generan un impacto sustancial sobre la renta gravable. Sus estimaciones implican, por ejemplo, que una reducción del tipo impositivo marginal del $40\%$ al $30\%$ a una persona de renta alta, aumentaría su base imponible en más de un $9\%$. De esta forma, la disminución de los ingresos sería menor que si no hubiera cambios en el comportamiento. Si bien investigaciones posteriores han encontrado una respuesta tanto superior como inferior, la literatura sugiere que el tipo impositivo que hace máximos los ingresos tributarios no es el mismo ni para todos los grupos de renta ni para todas las formas de renta.




\subsubsection{Evasión fiscal}
Por ejemplo, unos tipos marginales elevados pueden inducir a las personas a escoger ocupaciones que ofrecen abundantes oportunidades para evadir impuestos, lo que se conoce como la economía sumergida. Esta comprende tanto actividades económicas legales pero fáciles de ocultar a las autoridades tributarias (como las reparaciones domésticas) como actividades que son ilegales en sí mismas (como la prostitución o la venta de drogas). El tamaño de la economía sumergida es, por definición, muy difícil de medir. Por ejemplo, algunas estimaciones la sitúan en torno al $14\%$ del PIB para Estados Unidos (Friedman et al., 2000), o del $7\%$ en Reino Unido, mientras que en Rusia alcanzaría el $42\%$. 

DAVIS y HENREKSON (2004) analizaron datos de varios países desarrollados durante la década de 1990 y comprobaron que, cuando aumentan los tipos marginales, también aumenta la probabilidad de participar en el sector de la economía sumergida. Resultado que coincide con las informaciones periodísticas relativas a lo ocurrido en la ciudad de Nueva York cuando el incremento del impuesto sobre los cigarrillos elevó el precio por paquete hasta aproximadamente $7,50€$. La subida del impuesto alimentó un próspero mercado negro de cigarrillos procedentes de otros estados con impuestos más bajos, cuyos vendedores incluían no solo a contrabandistas experimentados, sino también a «aficionados en busca de ingresos adicionales» (Fairclough, 2002, p. B1).

En consecuencia, no se sabe con certeza si las multas elevadas o las auditorías frecuentes constituyen el medio más eficaz para disuadir el fraude. Un resultado provisional que emerge de diversos estudios econométricos es que, para la mayoría de los grupos, un aumento en la probabilidad de auditoría incrementa la renta declarada, aunque la magnitud del efecto es reducida (BLUMENTHAL et al., 2001).

\begin{specialblock}{Evidencia empírica: Evasión fiscal}
    La medida del nivel de evasión fiscal es fundamental para un país. y existen diferentes formas de hacerlo: i ) mediante encuestas, aunque existen incentivos a no revelar la realidad ; i i ) mediante la inferencia de la evasión fiscal (se mide la actividad total y posterionnente se elimina la actividad efectivamente medida, para obtener la economía sumergida).

    El trabajo econométrico de mayor relieve en este terreno corresponde a Clotfelter (1983). que utilizando datos fiscales de los EE.UU. estimó que la evasión aumenta con la renta disponible después de impuestos y con el tipo impositivo marginal. Sin embargo, Feldstein ( 1991 ). utilizando también datos fi scales. estimó que la evasión se reduce cuando aumenta el tipo marginal. La razón es que con un tipo marginal más elevado, si bien la ganancia de defraudar (por el ahorro fiscal) es mayor. la sanción en caso de ser descubierto es más elevada. Cuando el segundo efecto domina al primero, se producen resultados como el de Feldstein. Eso pone de mani fiesto una vez más que en la práctica las respuestas de los individuos a las medidas fiscales pueden ser menos obvias de lo que podría parecer a primera vista.

    La relación entre fraude y probabilidad de detección ha sido analizada por. entre otros. Blumenthal, Christian y Slemrod (2001 ). Utilizando los datos de un expe1imento fiscal realizado por Hacienda del Estado de Minnesota, llegaron a la conclusión de que al aumento en la probabilidad de inspección aumenta de forma moderada la renta declarada para las rentas medias y bajas. pero no para las rentas elevada<;. Según su conjetura. la razón es que las rentas altas consideran que una inspección es una ocasión para el regateo con la Administración y optan por iniciar la negociación con una renta inicial declarada baja.
\end{specialblock}



\clearpage
\section{Conclusión} 


\subsection{Recapitulación}


\subsection{Extensiones y relación con otras partes del Temario}



\subsection{Opinión}

\subsubsection{Idea final (Salida o cierra)}

\clearpage
\pagenumbering{gobble}
\MyBibliography



\MyBibItem{DAWID y GATTI}{Chapter 2 - Agent-Based Macroeconomics}{2018}{https://doi.org/10.1016/bs.hescom.2018.02.006}


% ANEXOS
\clearpage










\end{document}